\documentclass[12pt,a4paper]{report}

\usepackage[utf8]{inputenc}
\usepackage[T1]{fontenc}
\usepackage{mathptmx} % Times New Roman
\usepackage[left=1.5in, right=1in, top=1in, bottom=1in]{geometry}
\usepackage{setspace}
\usepackage{graphicx}
\usepackage{hyperref}
\usepackage{booktabs}
\usepackage{caption}
\usepackage{float}
\usepackage{listings}
\usepackage{xcolor}
\usepackage{titlesec}
\usepackage{tabularx}
\usepackage{lipsum} % For some extra text if needed
\usepackage{textcomp}

\onehalfspacing

% Formatting for code
\definecolor{codegreen}{rgb}{0,0.6,0}
\definecolor{codegray}{rgb}{0.5,0.5,0.5}
\definecolor{codepurple}{rgb}{0.58,0,0.82}
\definecolor{backcolour}{rgb}{0.95,0.95,0.92}

\lstdefinelanguage{Rust}{
  keywords={true,false,catch,return,null,catch,switch,if,in,while,do,else,case,break,match,enum,struct,impl,trait,pub,fn,let,mut,for,loop,match,ref,return,unsafe,use,super,self,crate,where,impl,type,const,static,dyn,as,break,continue,extern,macro_rules,box,yield,await,async,move},
  keywordstyle=\color{magenta}\bfseries,
  ndkeywords={bool,u8,u16,u32,u64,u128,i8,i16,i32,i64,i128,f32,f64,char,str,String,Vec,Option,Result,Box,Rc,Arc,Cell,RefCell,HashMap,BTreeMap,SmallVec},
  ndkeywordstyle=\color{codepurple}\bfseries,
  identifierstyle=\color{black},
  sensitive=false,
  comment=[l]{//},
  morecomment=[s]{/*}{*/},
  commentstyle=\color{codegreen}\ttfamily,
  stringstyle=\color{codepurple}\ttfamily,
  morestring=[b]',
  morestring=[b]"
}

\lstdefinestyle{mystyle}{
    backgroundcolor=\color{backcolour},   
    commentstyle=\color{codegreen},
    keywordstyle=\color{magenta},
    numberstyle=\tiny\color{codegray},
    stringstyle=\color{codepurple},
    basicstyle=\ttfamily\footnotesize,
    breakatwhitespace=false,         
    breaklines=true,                 
    captionpos=b,                    
    keepspaces=true,                 
    numbers=left,                    
    numbersep=5pt,                  
    showspaces=false,                
    showstringspaces=false,
    showtabs=false,                  
    tabsize=2,
    extendedchars=true,
    literate={‘}{{'}}1 {’}{{'}}1 {“}{{""}}1 {”}{{""}}1 {§}{{\S}}1
}
\lstset{style=mystyle}

% Heading formatting
\titleformat{\chapter}[display]
  {\normalfont\bfseries\Large\centering}
  {\chaptertitlename\ \thechapter}
  {10pt}{\Large}
\titlespacing*{\chapter}{0pt}{-20pt}{20pt}

\begin{document}

% -------------------------------------------------------------------
% 1. TITLE PAGE
% -------------------------------------------------------------------
\begin{titlepage}
    \begin{center}
        \vspace*{1cm}
        
        \Large
        \textbf{RaftKV: A Deterministic Raft-Replicated Key-Value Store in Rust}
        
        \vspace{1.5cm}
        
        \normalsize
        \textit{Submitted in partial fulfillment of the requirements for the degree of}\\
        \textbf{B.Tech Information Science \& Engineering}
        
        \vspace{1.5cm}
        
        \textbf{Submitted by}\\
        Jeevan U Gowda\\
        
        \vspace{1.5cm}
        
        \textbf{Under the Guidance of}\\
        Guide/Supervisor Name\\
        Designation\\
        
        \vfill
        
        \vspace{1cm}
        
        \textbf{Department of Information Science \& Engineering}\\
        \textbf{BMS College of Engineering}\\
        \textbf{Autonomous VTU}\\
        \textbf{Academic Year 2025--26}
        
    \end{center}
\end{titlepage}

\pagenumbering{roman}

% -------------------------------------------------------------------
% 2. CERTIFICATE / DECLARATION
% -------------------------------------------------------------------
\chapter*{Certificate}
\addcontentsline{toc}{chapter}{Certificate}
This is to certify that the project report entitled \textbf{"RaftKV: A Deterministic Raft-Replicated Key-Value Store in Rust"} submitted by \textbf{Jeevan U Gowda} in partial fulfillment of the requirements for the award of the degree of B.Tech in Information Science \& Engineering from \textbf{BMS College of Engineering}, is a bonafide record of the work carried out under my supervision and guidance.

\vspace{2cm}
\noindent
\begin{tabular}{@{}p{0.5\textwidth}p{0.5\textwidth}@{}}
\textbf{Signature of Guide} & \textbf{Signature of HOD} \\
Name of Guide & Name of HOD \\
Designation & Designation \\
\end{tabular}

\newpage
\chapter*{Student Declaration}
\addcontentsline{toc}{chapter}{Student Declaration}
We hereby declare that the project work entitled \textbf{"RaftKV: A Deterministic Raft-Replicated Key-Value Store in Rust"} is an authentic record of our own work carried out as requirements of the B.Tech degree in Information Science \& Engineering. The matter embodied in this report has not been submitted in part or full to any other university or institute for the award of any degree or diploma.

\vspace{2cm}
\noindent
\textbf{Signatures:}\\
1. \rule{4cm}{0.4pt} \quad (Jeevan U Gowda)

\newpage

% -------------------------------------------------------------------
% 3. ACKNOWLEDGEMENT
% -------------------------------------------------------------------
\chapter*{Acknowledgement}
\addcontentsline{toc}{chapter}{Acknowledgement}

We would like to express our profound gratitude to our project guide, \textbf{Guide Name}, for their invaluable support, continuous guidance, and encouragement throughout the course of this project. Their deep insights into distributed systems and Rust programming were instrumental in shaping this work.

We extend our sincere thanks to the Head of the Department and all faculty members of the Department of Information Science \& Engineering for providing us with the necessary resources and an environment conducive to learning.

Finally, we would like to thank our families and peers for their constant support and motivation during the development of this project.

\newpage

% -------------------------------------------------------------------
% 4. ABSTRACT
% -------------------------------------------------------------------
\chapter*{Abstract}
\addcontentsline{toc}{chapter}{Abstract}

Building robust distributed systems that maintain high availability and strong consistency is notoriously difficult. Network partitions, arbitrary machine crashes, and non-determinism often lead to split-brain scenarios and data loss. This project addresses the challenge by developing \textbf{RaftKV}, a distributed, fault-tolerant key-value store built entirely in Rust. 

RaftKV implements the Raft consensus algorithm with a unique architectural approach: the "pure core" pattern. By rigorously decoupling the consensus state machine from all side-effects (such as disk I/O, network communication, timers, and operating-system randomness), the core logic becomes fully deterministic and highly testable. 

The system leverages gRPC for inter-node communication and offers a Redis-compatible (RESP) protocol frontend, allowing seamless integration with existing Redis clients. Key features include leader election with pre-vote optimization, log replication, a persistent segmented log, and linearizable reads. The separation of the pure consensus brain from the runtime environment allows for exhaustive deterministic simulation, proving the system's resilience under severe simulated network partitions and crashes.

\newpage

% -------------------------------------------------------------------
% 5. TABLE OF CONTENTS
% -------------------------------------------------------------------
\tableofcontents
\newpage

% -------------------------------------------------------------------
% 6. LIST OF FIGURES & TABLES
% -------------------------------------------------------------------
\listoffigures
\addcontentsline{toc}{chapter}{List of Figures}
\newpage

\listoftables
\addcontentsline{toc}{chapter}{List of Tables}
\newpage

% -------------------------------------------------------------------
% 7. LIST OF ABBREVIATIONS
% -------------------------------------------------------------------
\chapter*{List of Abbreviations}
\addcontentsline{toc}{chapter}{List of Abbreviations}

\begin{description}
    \item[API] Application Programming Interface
    \item[CAP] Consistency, Availability, Partition Tolerance
    \item[CLI] Command Line Interface
    \item[CRC] Cyclic Redundancy Check
    \item[DFD] Data Flow Diagram
    \item[gRPC] gRPC Remote Procedure Calls
    \item[KV] Key-Value
    \item[MVC] Model View Controller
    \item[PRNG] Pseudo-Random Number Generator
    \item[RESP] REdis Serialization Protocol
    \item[RPC] Remote Procedure Call
    \item[TTL] Time To Live
    \item[UAT] User Acceptance Testing
    \item[UML] Unified Modeling Language
\end{description}

\newpage
\pagenumbering{arabic}

% -------------------------------------------------------------------
% CHAPTER 1 — INTRODUCTION
% -------------------------------------------------------------------
\chapter{Introduction}

\section{Background / Overview}
Distributed systems are the backbone of modern computing infrastructure. As applications scale beyond the capacity of a single machine, they must distribute their state across multiple servers. However, ensuring that all servers agree on the system's state in the presence of failures—such as network partitions, machine crashes, and message delays—is a fundamental challenge. The Raft consensus algorithm was introduced to solve this problem by managing a replicated log and ensuring that a cluster of machines functions as a single, coherent state machine.

\section{Problem Statement}
Implementing a consensus algorithm like Raft is historically prone to subtle concurrency bugs, primarily because consensus logic is often entangled with asynchronous I/O operations, networking, and system clocks. The lack of testability in such entangled designs leads to vulnerabilities in distributed databases. Thus, the problem is to design and build a fault-tolerant, replicated key-value store where the core consensus logic is completely decoupled from side-effects, making it fully testable and deterministically simulable.

\section{Objectives}
The primary objectives of this project are:
\begin{itemize}
    \item To implement the Raft consensus algorithm from first principles in Rust.
    \item To achieve a "pure/impure" architectural split, completely isolating the consensus logic from network and disk I/O.
    \item To build a persistent, segmented log and a deterministic key-value state machine.
    \item To expose a Redis-compatible (RESP) server interface so clients can interact with the cluster using standard tools.
    \item To build a deterministic simulator that can verify the system's safety invariants without asynchronous non-determinism.
\end{itemize}

\section{Scope of the Project}
The project covers the core Raft components: leader election, log replication, safety invariants, and the pre-vote extension. It includes a functional gRPC networking layer, persistent storage via segmented logs and sled, and a RESP frontend supporting basic KV operations (GET, SET, DEL, INCR, EXPIRE). It also supports joint consensus for membership changes. Out of scope for this phase is dynamic load balancing across multiple Raft groups (sharding).

\section{Motivation}
The motivation behind this project stems from the desire to deeply understand distributed consensus and to leverage Rust's strict type system and memory safety guarantees. Real-world systems like etcd and TiKV are complex; building a clean, pure-core Raft implementation provides profound educational value and serves as a blueprint for highly reliable backend infrastructure.

\section{Organization of the Report}
Chapter 2 reviews related work and existing systems. Chapter 3 covers system analysis and requirements. Chapter 4 details the architectural design and pure/impure split. Chapter 5 discusses the Rust implementation of the modules. Chapter 6 outlines the deterministic testing strategy. Chapter 7 presents results and performance discussions. Finally, Chapter 8 concludes the report and discusses future scope.

% -------------------------------------------------------------------
% CHAPTER 2 — LITERATURE REVIEW
% -------------------------------------------------------------------
\chapter{Literature Review}

\section{Overview of Related Work}
Consensus algorithms have been studied extensively. The Paxos algorithm, introduced by Leslie Lamport, was long the industry standard. However, Paxos is famously difficult to understand and implement correctly in a replicated log setting. In response, Diego Ongaro and John Ousterhout introduced Raft in 2014, designing it specifically for understandability without compromising safety or performance. Raft breaks consensus into independent subproblems: leader election, log replication, and safety.

\section{Review of Existing Systems}
Several prominent distributed systems utilize consensus algorithms:
\begin{itemize}
    \item \textbf{etcd:} Written in Go, etcd is the primary configuration and state store for Kubernetes. It uses Raft and has heavily influenced the community. However, its Raft core, while largely isolated, still interacts with Go's asynchronous runtime in ways that complicate purely deterministic simulation.
    \item \textbf{TiKV:} A distributed transactional key-value database originally built in Rust. It also uses a variant of the etcd Raft implementation ported to Rust. TiKV is highly performant but exceptionally complex.
    \item \textbf{Redis:} Redis itself provides replication and clustering but does not provide strong consistency guarantees by default (using asynchronous replication), though Redis Raft modules exist.
\end{itemize}

\section{Research Gap / Justification}
While there are production-grade Raft implementations (like etcd and TiKV), their codebases are vast and often blend business logic with side-effects, making educational study and exhaustive simulation difficult. \textbf{RaftKV} addresses this gap by rigorously enforcing a pure functional boundary for the Raft core. It accepts inputs (Messages, Ticks) and produces outputs (Actions) without executing any I/O itself. This allows for an elegant deterministic simulator to verify safety invariants without Tokio or wall-clock dependencies.

\section{Summary Table}
\begin{table}[H]
\centering
\begin{tabularx}{\textwidth}{|l|X|X|l|}
\hline
\textbf{System} & \textbf{Consensus} & \textbf{Key Feature} & \textbf{Language} \\ \hline
ZooKeeper & ZAB & High read throughput & Java \\ \hline
etcd & Raft & Core Kubernetes datastore & Go \\ \hline
TiKV & Raft (Multi-Raft) & HTAP, highly scalable & Rust \\ \hline
\textbf{RaftKV (Proposed)} & Raft (Pure Core) & Deterministic testing, RESP API & Rust \\ \hline
\end{tabularx}
\caption{Comparison of Distributed Key-Value Stores}
\end{table}

% -------------------------------------------------------------------
% CHAPTER 3 — SYSTEM ANALYSIS
% -------------------------------------------------------------------
\chapter{System Analysis}

\section{Feasibility Study}
\begin{itemize}
    \item \textbf{Technical Feasibility:} The project relies on mature Rust libraries (Tokio, Tonic, Sled, Bincode). Building a pure state machine in Rust is highly feasible due to its algebraic data types and borrow checker.
    \item \textbf{Operational Feasibility:} The system exposes a Redis-compatible API, meaning any standard Redis client (like \texttt{redis-cli}) can interact with the cluster without modification.
    \item \textbf{Economic Feasibility:} The project utilizes open-source technologies exclusively, incurring zero licensing costs.
\end{itemize}

\section{Requirements Analysis}
\subsection{Functional Requirements}
\begin{itemize}
    \item The system must elect a single leader per term.
    \item Clients must be able to connect via TCP and issue RESP commands.
    \item Write operations must be routed to the leader and appended to the replicated log.
    \item Commands must only be applied to the state machine once a quorum has persisted them.
    \item Read operations must be linearizable, ensuring clients see the latest committed state.
    \item The system must support dynamic membership changes via joint consensus.
\end{itemize}

\subsection{Non-Functional Requirements}
\begin{itemize}
    \item \textbf{Determinism:} The core consensus logic must be free of asynchronous timing dependencies.
    \item \textbf{Fault Tolerance:} A cluster of $2f+1$ nodes must continue to operate if up to $f$ nodes fail.
    \item \textbf{Data Durability:} Committed entries must survive node restarts.
\end{itemize}

\section{Software \& Hardware Requirements}
\begin{itemize}
    \item \textbf{Software:} Rust toolchain (cargo, rustc), Protocol Buffers (\texttt{protoc}), Docker (optional for deployment), Redis-CLI for testing.
    \item \textbf{Hardware:} Minimal hardware required. A multi-core processor and 4GB RAM are sufficient for running a local 3-node cluster. SSDs are recommended to minimize \texttt{fsync} latency for the durable log.
\end{itemize}

\section{Use Case Diagram}
\textit{(The Use Case Diagram visualizes the interaction between the Client, the Leader Node, and Follower Nodes. The Client submits commands. The Leader handles client requests, appends entries, and broadcasts heartbeats. Followers receive replication RPCs and vote in elections.)}

\section{Data Flow Diagram (DFD)}
\begin{itemize}
    \item \textbf{Level 0 DFD (Context Diagram):} Client sends RESP Frame $\rightarrow$ RaftKV System $\rightarrow$ Client receives RESP Response.
    \item \textbf{Level 1 DFD:} 
    \begin{enumerate}
        \item Client Request parsed into \texttt{ClientCommand}.
        \item Command proposed to \texttt{RaftNode} (Core).
        \item Core emits \texttt{Action::AppendEntries} and \texttt{Action::SendMessage}.
        \item Runtime persists log to \texttt{SegmentedLog}.
        \item Runtime sends RPC via \texttt{RaftNet} to peer nodes.
        \item Quorum achieved; Core emits \texttt{Action::ApplyCommitted}.
        \item \texttt{KvStateMachine} updates state.
    \end{enumerate}
\end{itemize}


% -------------------------------------------------------------------
% CHAPTER 4 — SYSTEM DESIGN
% -------------------------------------------------------------------
\chapter{System Design}

\section{System Architecture}
The defining characteristic of RaftKV is the strict separation between pure logic and side effects.
\begin{itemize}
    \item \textbf{\texttt{raft-core}:} The pure Raft brain. No I/O, no async, no clock. It receives \texttt{Message} and \texttt{Tick} inputs, mutates in-memory state, and emits \texttt{Action} directives.
    \item \textbf{\texttt{raftkv-server} (Runtime):} The event loop that drives the core. It executes \texttt{Action}s emitted by the core: persisting to disk, making network calls, and manipulating timers.
    \item \textbf{\texttt{raft-storage}:} Manages the segmented log on disk, sled-backed metadata (\texttt{HardState}), and snapshots.
    \item \textbf{\texttt{raft-net}:} Handles tonic/gRPC transport and protobuf serialization.
    \item \textbf{\texttt{resp-server}:} Translates raw TCP streams of Redis protocol into semantic client commands.
\end{itemize}

\section{Database Design}
The system does not use a traditional SQL database. It relies on embedded storage mechanisms.
\subsection{Segmented Log Structure}
Log entries are stored in append-only files (segments) named by their starting index (e.g., \texttt{00000000000000000001.log}).
Each record contains:
\begin{itemize}
    \item \texttt{u32}: Little-endian payload length.
    \item \texttt{u32}: Little-endian CRC32C of the payload.
    \item \texttt{payload}: Bincode-encoded \texttt{raft\_core::Entry}.
\end{itemize}
\subsection{Key-Value State Machine}
The \texttt{KvStateMachine} leverages the \texttt{sled} embedded database to maintain three trees:
\begin{enumerate}
    \item \texttt{kv}: Maps keys to values.
    \item \texttt{ttl}: Maps \texttt{expire\_at\_ms\_be || key} to empty values, functioning as a sorted index for expirations.
    \item \texttt{meta}: Stores the \texttt{applied\_index} and the rolling state hash.
\end{enumerate}

\section{Module Design}
\subsection{Raft Node (\texttt{raft-core})}
Contains essential state properties: \texttt{current\_term}, \texttt{voted\_for}, \texttt{commit\_index}, \texttt{role} (Follower, PreCandidate, Candidate, Leader), and an in-memory \texttt{RaftLog}.
\subsection{Raft RPCs}
Implemented via Protocol Buffers. Primary services: \texttt{AppendEntries} (replication/heartbeat), \texttt{RequestVote} (elections), \texttt{PreVote}, \texttt{InstallSnapshot}.

\section{UML Diagrams}
\textit{Sequence Diagram for Client Write:}
1. Client $\rightarrow$ Runtime: \texttt{SET K V}
2. Runtime $\rightarrow$ Core: \texttt{propose(SET K V)}
3. Core $\rightarrow$ Runtime: \texttt{Action::AppendEntries (local)}
4. Core $\rightarrow$ Runtime: \texttt{Action::SendMessage (to peers)}
5. Runtime $\rightarrow$ Disk: Write to segment.
6. Runtime $\rightarrow$ Peer: gRPC \texttt{AppendEntries}
7. Peer $\rightarrow$ Runtime: gRPC \texttt{AppendEntriesResponse}
8. Runtime $\rightarrow$ Core: \texttt{step(Response)}
9. Core $\rightarrow$ Runtime: \texttt{Action::ApplyCommitted}
10. Runtime $\rightarrow$ KV Store: Apply command.
11. Runtime $\rightarrow$ Client: RESP \texttt{+OK}


% -------------------------------------------------------------------
% CHAPTER 5 — IMPLEMENTATION
% -------------------------------------------------------------------
\chapter{Implementation}

\section{Technology Stack}
\begin{itemize}
    \item \textbf{Programming Language:} Rust. Chosen for memory safety without garbage collection, and fearless concurrency.
    \item \textbf{Asynchronous Runtime:} Tokio.
    \item \textbf{RPC Framework:} Tonic (gRPC) and Prost (Protocol Buffers).
    \item \textbf{Storage Engine:} Sled (embedded Key-Value store).
    \item \textbf{Serialization:} Bincode for efficient disk serialization, RESP for client communication.
\end{itemize}

\section{Development Environment Setup}
To run the environment, the Rust toolchain must be installed via \texttt{rustup}. 
\begin{lstlisting}[language=bash, caption=Build and Run Commands]
# Compile the workspace
cargo build --workspace --release

# Run a local 3-node cluster using the provided script
./scripts/run_cluster.sh
\end{lstlisting}

\section{Module-wise Implementation Details}
\subsection{The Pure Core (\texttt{RaftNode})}
The core is implemented as a state machine. It handles transitions seamlessly. For example, the universal term-bump rule ensures safety: if any incoming RPC has a term strictly greater than \texttt{current\_term}, the node immediately drops to the Follower role, clears its \texttt{voted\_for} field, updates its term, and emits an \texttt{Action::PersistHardState}.

\subsection{Pre-Vote Extension}
To prevent partitioned nodes from disrupting the cluster with inflated terms upon reconnection, RaftKV implements the Pre-Vote algorithm. A \texttt{PreCandidate} asks peers if they \textit{would} vote for it in the next term, without actually incrementing the term. Only upon receiving a pre-vote quorum does it increment the term and become a full \texttt{Candidate}.

\subsection{Linearizable Reads}
Reads in distributed systems are challenging due to stale leaders. RaftKV handles \texttt{GET} requests via a read-index protocol. The leader records its current \texttt{commit\_index}, verifies its leadership via a heartbeat quorum, waits for the local \texttt{applied\_index} to catch up to the recorded index, and only then serves the read from the local state machine.

\section{UI Screens / Output}
Interaction is done via \texttt{redis-cli}.
\begin{lstlisting}[caption=Redis CLI Interaction]
$ redis-cli -p 6379 PING
PONG
$ redis-cli -p 6379 SET message "Hello Raft"
OK
$ redis-cli -p 6380 GET message
"Hello Raft"
\end{lstlisting}
Note that \texttt{GET} on port 6380 (a follower) will either redirect to the leader or serve the read if read-index is handled.


% -------------------------------------------------------------------
% CHAPTER 6 — TESTING
% -------------------------------------------------------------------
\chapter{Testing}

\section{Testing Strategy}
Because of the pure/impure architectural split, testing is exceptionally rigorous. The system relies heavily on \textbf{Deterministic Simulation}. Instead of spinning up threads and relying on wall-clock time, the simulator instantiates multiple \texttt{RaftNode} instances, controls a virtual clock, and routes messages through a simulated network that can drop, delay, or reorder packets deterministically based on a PRNG seed.

\section{Types of Testing Performed}
\begin{itemize}
    \item \textbf{Unit Testing:} Tests individual functions (e.g., CRC checks, serialization formats, RESP parsing).
    \item \textbf{Deterministic Simulation:} Simulates cluster behavior over logical ticks. Tests leader election, split votes, and log convergence under partitions.
    \item \textbf{Linearizability Checking:} A custom checker implements a Wing-and-Gong style exhaustive search for short histories. It verifies that concurrent operations on the cluster resolve to a valid, linearizable history.
\end{itemize}

\section{Test Cases}
\begin{table}[H]
\centering
\begin{tabularx}{\textwidth}{|l|X|X|X|l|}
\hline
\textbf{TC ID} & \textbf{Description} & \textbf{Input} & \textbf{Expected Output} & \textbf{Status} \\ \hline
TC-01 & Clean Election & 3 nodes boot simultaneously & Single leader elected within random timeout window & Pass \\ \hline
TC-02 & Client Write & \texttt{SET K V} to Leader & Log replicated, applied, returns \texttt{OK} & Pass \\ \hline
TC-03 & Partition Isolation & Isolate leader from quorum & Leader cannot commit new writes; new leader elected in majority & Pass \\ \hline
TC-04 & Log Conflict & Reconnect partitioned old leader & Conflicting entries truncated, logs converge with new leader & Pass \\ \hline
TC-05 & Follower Write & \texttt{SET X Y} to Follower & Returns \texttt{MOVED} redirect to current Leader & Pass \\ \hline
\end{tabularx}
\caption{System Test Cases}
\end{table}

\section{Bugs \& Fixes}
During development, a classic Raft bug known as the "Figure 8" issue was carefully mitigated. A leader is restricted from advancing the commit index by counting replicas for entries not in its \textit{current} term. The \texttt{maybe\_advance\_commit} function checks that \texttt{log.term\_at(new\_commit) == current\_term} before committing, preventing stale data overwrites.

% -------------------------------------------------------------------
% CHAPTER 7 — RESULTS & DISCUSSION
% -------------------------------------------------------------------
\chapter{Results \& Discussion}

\section{System Performance}
The system demonstrated low-latency responses for key-value operations. By avoiding \texttt{fsync} on every single append and instead grouping syncs at the end of an action batch (\texttt{sync\_each\_append = false}), throughput was significantly increased without violating safety invariants, provided the runtime persists hard state prior to responding to RPCs.

\section{Architectural Triumphs}
The separation of the consensus logic (\texttt{raft-core}) from the execution environment (\texttt{raftkv-server}) proved to be the most successful decision of the project. It allowed the implementation of a deterministic simulator (\texttt{crates/sim-tests}) which made debugging split-brain scenarios and log truncation conflicts trivial. When a bug occurred, the exact seed could be replayed to recreate the failure consistently, bypassing the notorious difficulty of debugging asynchronous race conditions in distributed systems.

\section{Comparison with Existing Systems}
Compared to a standard Redis instance, RaftKV sacrifices raw throughput in exchange for strong consistency and fault tolerance. Unlike standard Redis replication, which is asynchronous and can lose writes during failover, RaftKV ensures that once a client receives an \texttt{OK} response, the data is durable on a quorum and will survive the loss of the leader.

% -------------------------------------------------------------------
% CHAPTER 8 — CONCLUSION & FUTURE SCOPE
% -------------------------------------------------------------------
\chapter{Conclusion \& Future Scope}

\section{Conclusion}
The RaftKV project successfully implemented a highly robust, deterministic Raft consensus engine in Rust. By adhering to a strict pure/impure design philosophy, the system demonstrated how complex distributed algorithms can be made highly testable. The integration of a RESP frontend bridges the gap between complex distributed systems and user-friendly interaction, allowing seamless integration with the massive Redis ecosystem. The system fulfills all primary objectives: electing leaders securely, replicating logs reliably, and applying state machine changes safely.

\section{Limitations}
While fully functional, the current implementation has a few limitations:
\begin{itemize}
    \item \textbf{Snapshot Resumption:} While the core log supports compaction, the runtime \texttt{InstallSnapshot} does not yet fully restore KV state bytes from a streaming chunked transfer.
    \item \textbf{TTL Expiry:} While TTL indices exist, the automatic replicated \texttt{Tick} proposal is not yet continuously wired in the runtime to evict keys actively without reads.
\end{itemize}

\section{Future Scope}
\begin{itemize}
    \item \textbf{Full Snapshotting:} Implement streaming snapshot transfers across the network for massive datasets.
    \item \textbf{Cloud Deployment:} Containerize and deploy via Helm charts to a Kubernetes cluster for scale testing.
    \item \textbf{Jepsen Testing:} Subject the system to the Jepsen testing framework to validate linearizability under extreme, pathological network fault injections.
    \item \textbf{Multi-Raft (Sharding):} Implement multiple Raft groups to scale write throughput horizontally, similar to TiKV.
\end{itemize}

% -------------------------------------------------------------------
% REFERENCES / BIBLIOGRAPHY
% -------------------------------------------------------------------
\renewcommand{\bibname}{References}
\begin{thebibliography}{99}
\addcontentsline{toc}{chapter}{References}

\bibitem{ongaro2014}
D. Ongaro and J. Ousterhout, \textit{"In Search of an Understandable Consensus Algorithm,"} in Proceedings of the 2014 USENIX Annual Technical Conference (USENIX ATC 14), 2014, pp. 305-319.

\bibitem{ongarothesis}
D. Ongaro, \textit{"Consensus: Bridging Theory and Practice,"} Ph.D. dissertation, Stanford University, 2014. (Introduced the Pre-Vote extension).

\bibitem{rustbook}
S. Klabnik and C. Nichols, \textit{"The Rust Programming Language,"} No Starch Press, 2018.

\bibitem{lamportpaxos}
L. Lamport, \textit{"The Part-Time Parliament,"} ACM Transactions on Computer Systems (TOCS), vol. 16, no. 2, pp. 133-169, 1998.

\bibitem{redisdoc}
Redis Ltd., \textit{"Redis Protocol Specification (RESP),"} [Online]. Available: https://redis.io/topics/protocol.

\bibitem{grpc}
Cloud Native Computing Foundation, \textit{"gRPC: A high performance, open source universal RPC framework,"} [Online]. Available: https://grpc.io/.

\bibitem{etcd}
etcd Authors, \textit{"etcd: A distributed, reliable key-value store,"} [Online]. Available: https://etcd.io/.

\bibitem{tikv}
TiKV Authors, \textit{"TiKV: A Distributed Transactional Key-Value Database,"} [Online]. Available: https://tikv.org/.

\end{thebibliography}

% -------------------------------------------------------------------
% APPENDICES
% -------------------------------------------------------------------
\appendix
\chapter{User Manual}

\section{Prerequisites}
\begin{itemize}
    \item Rust 1.70+ (\texttt{rustup})
    \item \texttt{redis-cli}
    \item \texttt{protoc} (Protocol Buffers Compiler)
\end{itemize}

\section{Running the Cluster Locally}
Navigate to the root directory and start the simulated cluster:
\begin{lstlisting}[language=bash]
$ chmod +x scripts/run_cluster.sh
$ ./scripts/run_cluster.sh
\end{lstlisting}
This will spawn three nodes listening on RESP ports 6379, 6380, and 6381.

\section{Interacting with the System}
Use the Redis CLI to connect to Node 1 (Port 6379, typical Leader):
\begin{lstlisting}[language=bash]
$ redis-cli -p 6379
127.0.0.1:6379> SET project "RaftKV"
OK
127.0.0.1:6379> GET project
"RaftKV"
127.0.0.1:6379> CLUSTER NODES
(Returns list of nodes and current leader)
\end{lstlisting}

\chapter{Source Code Overview}
This appendix contains the complete source code for the RaftKV project, structured by module.

\section{crates/kv-state-machine/src/command.rs}
\begin{lstlisting}[language=Rust]
//! Command and response types. Encoded with bincode into Raft log entries.

use serde::{Deserialize, Serialize};

/// Logical command type. The leader records `commit_ts_millis` (taken from the
/// proposing leader's monotonic clock) so that TTL evaluation is deterministic
/// across replicas (every replica sees the same timestamp in the same entry).
#[derive(Debug, Clone, Serialize, Deserialize)]
pub enum Command {
    /// SET key value [EX seconds]
    Set {
        /// Key.
        key: Vec<u8>,
        /// Value.
        value: Vec<u8>,
        /// Optional absolute expiry, in milliseconds since the proposing
        /// leader's epoch (filled in by the proposer).
        expire_at_ms: Option<u64>,
    },
    /// DEL key1 key2 ...
    Del {
        /// Keys to delete.
        keys: Vec<Vec<u8>>,
    },
    /// INCR key
    Incr {
        /// Key.
        key: Vec<u8>,
        /// Delta (can be negative for DECR).
        delta: i64,
    },
    /// EXPIRE key seconds  (absolute time computed at proposal time).
    Expire {
        /// Key.
        key: Vec<u8>,
        /// Absolute expiry ms (deterministic).
        expire_at_ms: u64,
    },
    /// PERSIST key
    Persist {
        /// Key.
        key: Vec<u8>,
    },
    /// MSET k1 v1 k2 v2 ...
    MSet {
        /// Pairs.
        pairs: Vec<(Vec<u8>, Vec<u8>)>,
    },
    /// FLUSHDB
    FlushDb,
    /// Tick: a no-op command carrying the current time so the sweeper can
    /// evict expired keys deterministically.
    Tick {
        /// Now in ms.
        now_ms: u64,
    },
}

/// Response to a command.
#[derive(Debug, Clone, Serialize, Deserialize)]
pub enum Response {
    /// "OK" string.
    Ok,
    /// Simple integer reply.
    Int(i64),
    /// Bulk reply (`Some` value or nil).
    Bulk(Option<Vec<u8>>),
    /// Multi-bulk reply.
    Array(Vec<Option<Vec<u8>>>),
    /// Error reply.
    Error(String),
}
\end{lstlisting}

\section{crates/kv-state-machine/src/lib.rs}
\begin{lstlisting}[language=Rust]
#![deny(unsafe_code)]
//! Deterministic KV state machine applied from the committed Raft log.
//!
//! NOTE: The original spec calls for RocksDB. This implementation uses sled
//! to avoid the RocksDB build toolchain dependency. Both are LSM-based and
//! provide WAL durability + atomic flush semantics; the trait-based
//! abstraction means swapping the backend later is mechanical. See
//! ARCHITECTURE.md → Design Decisions.

pub mod command;
pub mod state;

pub use command::{Command, Response};
pub use state::{KvStateMachine, StateMachineError};
\end{lstlisting}

\section{crates/kv-state-machine/src/state.rs}
\begin{lstlisting}[language=Rust]
//! sled-backed KV state machine. Applies commands deterministically.
//!
//! Trees:
//! * `kv`  : `key -> value`
//! * `ttl` : `expire_at_ms_be(8 bytes) || key -> ()`  (sorted index of TTLs)
//! * `meta`: bookkeeping (state hash, applied index)

use std::path::Path;

use parking_lot::Mutex;
use raft_core::log::LogIndex;
use sha2::{Digest, Sha256};
use sled::IVec;
use thiserror::Error;

use crate::command::{Command, Response};

const META_APPLIED: &[u8] = b"applied_index";
const META_HASH: &[u8] = b"state_hash";

/// State machine errors.
#[derive(Debug, Error)]
pub enum StateMachineError {
    /// Sled error.
    #[error("sled error: {0}")]
    Sled(#[from] sled::Error),
    /// Codec error.
    #[error("codec: {0}")]
    Codec(String),
}

/// KV state machine.
pub struct KvStateMachine {
    db: sled::Db,
    kv: sled::Tree,
    ttl: sled::Tree,
    meta: sled::Tree,
    hash: Mutex<[u8; 32]>,
}

fn ttl_key(expire_at_ms: u64, key: &[u8]) -> Vec<u8> {
    let mut k = Vec::with_capacity(8 + key.len());
    k.extend_from_slice(&expire_at_ms.to_be_bytes());
    k.extend_from_slice(key);
    k
}

impl KvStateMachine {
    /// Open or create a state machine at the given directory.
    pub fn open<P: AsRef<Path>>(path: P) -> Result<Self, StateMachineError> {
        let db = sled::Config::new().path(path).flush_every_ms(Some(200)).open()?;
        let kv = db.open_tree("kv")?;
        let ttl = db.open_tree("ttl")?;
        let meta = db.open_tree("meta")?;
        let hash = match meta.get(META_HASH)? {
            Some(b) if b.len() == 32 => {
                let mut h = [0u8; 32];
                h.copy_from_slice(&b);
                h
            }
            _ => [0u8; 32],
        };
        Ok(Self { db, kv, ttl, meta, hash: Mutex::new(hash) })
    }

    /// Last applied log index.
    pub fn applied_index(&self) -> LogIndex {
        self.meta
            .get(META_APPLIED)
            .ok()
            .flatten()
            .and_then(|b| {
                let mut buf = [0u8; 8];
                if b.len() >= 8 { buf.copy_from_slice(&b[..8]); Some(u64::from_le_bytes(buf)) } else { None }
            })
            .unwrap_or(0)
    }

    /// Snapshot of state hash (32 bytes, hex-encodable).
    pub fn state_hash(&self) -> [u8; 32] { *self.hash.lock() }

    /// Apply a deserialized command. Returns the response.
    pub fn apply(&self, applied_index: LogIndex, cmd: &Command) -> Result<Response, StateMachineError> {
        let resp = match cmd {
            Command::Set { key, value, expire_at_ms } => {
                self.kv.insert(key.as_slice(), value.as_slice())?;
                if let Some(t) = expire_at_ms {
                    self.ttl.insert(ttl_key(*t, key), &[] as &[u8])?;
                }
                self.update_hash(b"S", key, value);
                Response::Ok
            }
            Command::Del { keys } => {
                let mut n = 0i64;
                for k in keys {
                    if self.kv.remove(k.as_slice())?.is_some() { n += 1; }
                    self.update_hash(b"D", k, &[]);
                }
                Response::Int(n)
            }
            Command::Incr { key, delta } => {
                let cur = self.kv.get(key.as_slice())?.and_then(|v: IVec| {
                    std::str::from_utf8(&v).ok().and_then(|s| s.parse::<i64>().ok())
                }).unwrap_or(0);
                let new = cur.saturating_add(*delta);
                let s = new.to_string();
                self.kv.insert(key.as_slice(), s.as_bytes())?;
                self.update_hash(b"I", key, s.as_bytes());
                Response::Int(new)
            }
            Command::Expire { key, expire_at_ms } => {
                if self.kv.contains_key(key.as_slice())? {
                    self.ttl.insert(ttl_key(*expire_at_ms, key), &[] as &[u8])?;
                    self.update_hash(b"E", key, &expire_at_ms.to_le_bytes());
                    Response::Int(1)
                } else {
                    Response::Int(0)
                }
            }
            Command::Persist { key } => {
                let mut removed = 0;
                let prefix = key.as_slice();
                let to_remove: Vec<_> = self.ttl.iter().keys().filter_map(|k| k.ok()).filter(|k| {
                    if k.len() < 8 { return false; }
                    &k[8..] == prefix
                }).collect();
                for k in to_remove { self.ttl.remove(k)?; removed += 1; }
                Response::Int(if removed > 0 { 1 } else { 0 })
            }
            Command::MSet { pairs } => {
                for (k, v) in pairs {
                    self.kv.insert(k.as_slice(), v.as_slice())?;
                    self.update_hash(b"M", k, v);
                }
                Response::Ok
            }
            Command::FlushDb => {
                self.kv.clear()?;
                self.ttl.clear()?;
                *self.hash.lock() = [0u8; 32];
                Response::Ok
            }
            Command::Tick { now_ms } => {
                self.evict_expired(*now_ms)?;
                Response::Ok
            }
        };
        // Bookkeeping (atomic-ish; sled provides WAL but not transactional
        // multi-tree atomicity — for this project we accept that and rely on
        // applied_index being persisted last so a crash mid-apply re-applies).
        self.meta.insert(META_APPLIED, &applied_index.to_le_bytes())?;
        let h = *self.hash.lock();
        self.meta.insert(META_HASH, &h)?;
        Ok(resp)
    }

    /// GET (read-side). Linearizable reads must be gated by read-index in the
    /// runtime layer — this function does not enforce that.
    pub fn get(&self, key: &[u8]) -> Result<Option<Vec<u8>>, StateMachineError> {
        Ok(self.kv.get(key)?.map(|v| v.to_vec()))
    }

    /// EXISTS.
    pub fn exists(&self, key: &[u8]) -> Result<bool, StateMachineError> {
        Ok(self.kv.contains_key(key)?)
    }

    /// MGET.
    pub fn mget(&self, keys: &[Vec<u8>]) -> Result<Vec<Option<Vec<u8>>>, StateMachineError> {
        keys.iter().map(|k| self.get(k.as_slice())).collect()
    }

    /// DBSIZE.
    pub fn len(&self) -> usize { self.kv.len() }
    /// Whether store is empty.
    pub fn is_empty(&self) -> bool { self.kv.is_empty() }

    /// Force a sled flush (durability).
    pub fn flush(&self) -> Result<(), StateMachineError> {
        self.db.flush()?;
        Ok(())
    }

    fn evict_expired(&self, now_ms: u64) -> Result<(), StateMachineError> {
        // Iterate keys whose timestamp <= now and delete.
        let upper = (now_ms + 1).to_be_bytes().to_vec();
        let to_delete: Vec<_> = self
            .ttl
            .range(..upper.clone())
            .keys()
            .filter_map(|k| k.ok())
            .collect();
        for tk in to_delete {
            if tk.len() >= 8 {
                let key = &tk[8..];
                self.kv.remove(key)?;
                self.update_hash(b"X", key, &[]);
            }
            self.ttl.remove(tk)?;
        }
        Ok(())
    }

    fn update_hash(&self, op: &[u8], key: &[u8], value: &[u8]) {
        let mut h = Sha256::new();
        let cur = self.hash.lock();
        h.update(*cur);
        drop(cur);
        h.update(op);
        h.update((key.len() as u64).to_le_bytes());
        h.update(key);
        h.update((value.len() as u64).to_le_bytes());
        h.update(value);
        let out = h.finalize();
        let mut buf = [0u8; 32];
        buf.copy_from_slice(&out);
        *self.hash.lock() = buf;
    }
}

#[cfg(test)]
mod tests {
    use super::*;
    use tempfile::tempdir;

    #[test]
    fn set_get_del() {
        let dir = tempdir().unwrap();
        let sm = KvStateMachine::open(dir.path()).unwrap();
        sm.apply(1, &Command::Set { key: b"a".to_vec(), value: b"1".to_vec(), expire_at_ms: None }).unwrap();
        assert_eq!(sm.get(b"a").unwrap(), Some(b"1".to_vec()));
        sm.apply(2, &Command::Del { keys: vec![b"a".to_vec()] }).unwrap();
        assert_eq!(sm.get(b"a").unwrap(), None);
    }

    #[test]
    fn determinism() {
        // Same command sequence on two state machines yields the same hash.
        let d1 = tempdir().unwrap();
        let d2 = tempdir().unwrap();
        let s1 = KvStateMachine::open(d1.path()).unwrap();
        let s2 = KvStateMachine::open(d2.path()).unwrap();
        let cmds = vec![
            Command::Set { key: b"x".to_vec(), value: b"1".to_vec(), expire_at_ms: None },
            Command::Set { key: b"y".to_vec(), value: b"2".to_vec(), expire_at_ms: None },
            Command::Incr { key: b"x".to_vec(), delta: 5 },
        ];
        for (i, c) in cmds.iter().enumerate() {
            s1.apply(i as u64 + 1, c).unwrap();
            s2.apply(i as u64 + 1, c).unwrap();
        }
        assert_eq!(s1.state_hash(), s2.state_hash());
    }
}
\end{lstlisting}

\section{crates/linearizability-checker/src/lib.rs}
\begin{lstlisting}[language=Rust]
#![deny(unsafe_code)]
//! Linearizability checker for register histories.
//!
//! Implements a Wing-and-Gong style search: given a list of operations
//! recorded as `(invocation_ts, response_ts, op)`, find a serialization order
//! consistent with the model (here: a register supporting `Get`, `Set`, `Del`,
//! `Incr`). Operations may be concurrent — any pair of operations whose
//! `[invocation, response]` intervals overlap may be ordered either way.
//!
//! For deterministic correctness the checker is exhaustive (exponential in
//! the number of concurrent operations). It is intended for short integration
//! traces, not for production-scale Jepsen replay; for very long histories
//! use Knossos or Porcupine.

use std::collections::BTreeMap;

use serde::{Deserialize, Serialize};
use thiserror::Error;

/// One operation in a recorded history.
#[derive(Debug, Clone, Serialize, Deserialize)]
pub struct Operation {
    /// Logical client id (used only for diagnostics).
    pub client: u64,
    /// Invocation timestamp (any monotonic unit).
    pub invocation: u64,
    /// Response timestamp.
    pub response: u64,
    /// Operation kind.
    pub op: Op,
}

/// Operation variants (subset of the KV API sufficient for checker traces).
#[derive(Debug, Clone, Serialize, Deserialize)]
pub enum Op {
    /// `GET key -> value`.
    Get { key: Vec<u8>, observed: Option<Vec<u8>> },
    /// `SET key = value -> Ok`.
    Set { key: Vec<u8>, value: Vec<u8> },
    /// `DEL key -> bool` (returns true if a value was removed).
    Del { key: Vec<u8>, observed_removed: bool },
    /// `INCR key by delta -> new_value`.
    Incr { key: Vec<u8>, delta: i64, observed: i64 },
}

/// Result of a linearization attempt.
#[derive(Debug)]
pub enum Verdict {
    /// History is linearizable (a valid serialization was found).
    Linearizable,
    /// History is not linearizable.
    NotLinearizable {
        /// Best-effort offending op index.
        offending_index: usize,
    },
}

/// Errors raised by the checker.
#[derive(Debug, Error)]
pub enum CheckerError {
    /// History contained an inconsistent invocation/response order.
    #[error("inconsistent invocation/response on op {0}")]
    InconsistentTimestamps(usize),
}

/// Check whether `history` is linearizable as a register/KV store.
///
/// Algorithm: depth-first search over operation orderings constrained by
/// real-time happens-before. At each step pick a "minimum" operation (one
/// whose invocation precedes the response of every other un-linearized op)
/// and verify it is consistent with the current model state.
pub fn check(history: &[Operation]) -> Result<Verdict, CheckerError> {
    for (i, op) in history.iter().enumerate() {
        if op.invocation > op.response { return Err(CheckerError::InconsistentTimestamps(i)); }
    }
    let mut state: BTreeMap<Vec<u8>, Vec<u8>> = BTreeMap::new();
    let mut remaining: Vec<usize> = (0..history.len()).collect();
    if dfs(history, &mut remaining, &mut state) {
        Ok(Verdict::Linearizable)
    } else {
        Ok(Verdict::NotLinearizable { offending_index: 0 })
    }
}

fn dfs(
    history: &[Operation],
    remaining: &mut Vec<usize>,
    state: &mut BTreeMap<Vec<u8>, Vec<u8>>,
) -> bool {
    if remaining.is_empty() { return true; }
    // Find minimum-set: ops whose invocation <= min(response of others).
    let earliest_response = remaining.iter().map(|&i| history[i].response).min().unwrap_or(u64::MAX);
    let candidates: Vec<usize> = remaining
        .iter()
        .copied()
        .filter(|&i| history[i].invocation <= earliest_response)
        .collect();
    for c in candidates {
        if !apply_consistent(&history[c], state) { continue; }
        let saved = state.clone();
        let pos = remaining.iter().position(|&x| x == c).unwrap();
        remaining.remove(pos);
        if dfs(history, remaining, state) { return true; }
        remaining.insert(pos, c);
        *state = saved;
    }
    false
}

fn apply_consistent(op: &Operation, state: &mut BTreeMap<Vec<u8>, Vec<u8>>) -> bool {
    match &op.op {
        Op::Get { key, observed } => state.get(key).cloned() == *observed,
        Op::Set { key, value } => { state.insert(key.clone(), value.clone()); true }
        Op::Del { key, observed_removed } => {
            let had = state.remove(key).is_some();
            had == *observed_removed
        }
        Op::Incr { key, delta, observed } => {
            let cur = state.get(key)
                .and_then(|v| std::str::from_utf8(v).ok().and_then(|s| s.parse::<i64>().ok()))
                .unwrap_or(0);
            let new = cur.saturating_add(*delta);
            if new != *observed { return false; }
            state.insert(key.clone(), new.to_string().into_bytes());
            true
        }
    }
}

#[cfg(test)]
mod tests {
    use super::*;

    fn op(client: u64, t0: u64, t1: u64, op: Op) -> Operation {
        Operation { client, invocation: t0, response: t1, op }
    }

    #[test]
    fn sequential_set_get_is_linearizable() {
        let h = vec![
            op(1, 0, 1, Op::Set { key: b"x".to_vec(), value: b"1".to_vec() }),
            op(2, 2, 3, Op::Get { key: b"x".to_vec(), observed: Some(b"1".to_vec()) }),
        ];
        assert!(matches!(check(&h).unwrap(), Verdict::Linearizable));
    }

    #[test]
    fn stale_read_after_set_is_not_linearizable() {
        let h = vec![
            op(1, 0, 1, Op::Set { key: b"x".to_vec(), value: b"1".to_vec() }),
            op(2, 2, 3, Op::Get { key: b"x".to_vec(), observed: None }),
        ];
        assert!(matches!(check(&h).unwrap(), Verdict::NotLinearizable { .. }));
    }

    #[test]
    fn concurrent_set_picks_a_winner() {
        // Two concurrent SETs; subsequent GET must see one of them.
        let h = vec![
            op(1, 0, 5, Op::Set { key: b"x".to_vec(), value: b"a".to_vec() }),
            op(2, 0, 5, Op::Set { key: b"x".to_vec(), value: b"b".to_vec() }),
            op(3, 6, 7, Op::Get { key: b"x".to_vec(), observed: Some(b"a".to_vec()) }),
        ];
        assert!(matches!(check(&h).unwrap(), Verdict::Linearizable));
    }
}
\end{lstlisting}

\section{crates/raft-core/src/action.rs}
\begin{lstlisting}[language=Rust]
//! Actions the state machine asks the runtime to perform.
//!
//! The pure core never performs I/O. Every side effect goes through this
//! enum, which the runtime interprets.

use bytes::Bytes;

use crate::log::{Entry, LogIndex, Term};
use crate::message::{Message, RequestId};
use crate::node::NodeId;
use crate::state::HardState;

/// Side-effect requests emitted by the core.
#[derive(Debug, Clone)]
pub enum Action {
    /// Persist `(currentTerm, votedFor, commitIndex)` before any network reply
    /// that depends on them (Raft §5.2).
    PersistHardState(HardState),
    /// Append entries to durable log storage.
    AppendEntries(Vec<Entry>),
    /// Truncate the log so no entries with `index >= from` remain.
    TruncateLog {
        /// Truncation start (inclusive).
        from: LogIndex,
    },
    /// Send a Raft message to a peer.
    SendMessage {
        /// Destination peer.
        to: NodeId,
        /// Message body.
        msg: Message,
    },
    /// Apply a slice of committed entries to the state machine.
    ApplyCommitted {
        /// Newly applicable entries (already in the durable log).
        entries: Vec<Entry>,
    },
    /// Leader requests a snapshot be taken at or before the given index.
    TakeSnapshot {
        /// Last included index of the snapshot.
        last_included_index: LogIndex,
        /// Term at `last_included_index`.
        last_included_term: Term,
    },
    /// Install a received snapshot (replace state machine + truncate log).
    InstallSnapshot {
        /// Snapshot meta.
        last_included_index: LogIndex,
        /// Snapshot meta.
        last_included_term: Term,
        /// Snapshot bytes.
        data: Bytes,
    },
    /// Reset the election timer (election timeout starts now).
    ResetElectionTimer,
    /// Reset the heartbeat timer (only meaningful as leader).
    ResetHeartbeatTimer,
    /// Notify any read-index waiter that index `commit_index` is observed by a
    /// quorum and reads up to that point may proceed once applied.
    NotifyReadIndex {
        /// Opaque caller-supplied context.
        ctx: Bytes,
        /// Confirmed commit index at the time of issuance.
        commit_index: LogIndex,
        /// Trace correlation.
        request_id: RequestId,
    },
    /// Promotion to leader has just occurred — used by the runtime for
    /// metrics/triggers.
    BecameLeader {
        /// Term in which leadership was won.
        term: Term,
    },
    /// Demotion to follower for a given (term, leader-hint).
    BecameFollower {
        /// New current term.
        term: Term,
        /// Best-known leader, if any.
        leader: Option<NodeId>,
    },
    /// Free-form metric event for observability.
    Metric(MetricEvent),
}

/// Observable events for metrics + tracing.
#[derive(Debug, Clone)]
pub enum MetricEvent {
    /// Term was advanced.
    TermBumped {
        /// New term.
        new_term: Term,
    },
    /// An election started.
    ElectionStarted {
        /// Election term.
        term: Term,
        /// Whether this was a pre-vote.
        pre_vote: bool,
    },
    /// An election succeeded.
    ElectionWon {
        /// Won-in term.
        term: Term,
    },
    /// AppendEntries was rejected by a follower.
    AppendRejected {
        /// Peer rejecting.
        from: NodeId,
    },
    /// Commit index advanced.
    CommitAdvanced {
        /// New commit index.
        index: LogIndex,
    },
    /// Leadership transfer initiated.
    LeadershipTransferStarted {
        /// Target node id.
        target: NodeId,
    },
}
\end{lstlisting}

\section{crates/raft-core/src/config.rs}
\begin{lstlisting}[language=Rust]
//! Cluster configuration and joint-consensus state machine.
//!
//! Implements §6 of the Raft paper and §4 of the Ongaro thesis.

use std::collections::BTreeSet;

use serde::{Deserialize, Serialize};

use crate::node::NodeId;

/// Configuration phase.
#[derive(Debug, Clone, PartialEq, Eq, Serialize, Deserialize)]
pub enum ConfigState {
    /// Stable configuration `C`.
    Stable(BTreeSet<NodeId>),
    /// Joint configuration `C_old,new`.
    Joint {
        /// Old voter set.
        old: BTreeSet<NodeId>,
        /// New voter set.
        new: BTreeSet<NodeId>,
    },
}

impl Default for ConfigState {
    fn default() -> Self { Self::Stable(BTreeSet::new()) }
}

impl ConfigState {
    /// All voters relevant for any quorum check.
    pub fn voters(&self) -> Vec<NodeId> {
        match self {
            Self::Stable(s) => s.iter().copied().collect(),
            Self::Joint { old, new } => old.union(new).copied().collect(),
        }
    }

    /// Returns true if `id` is a voter in any current configuration.
    pub fn is_voter(&self, id: NodeId) -> bool {
        match self {
            Self::Stable(s) => s.contains(&id),
            Self::Joint { old, new } => old.contains(&id) || new.contains(&id),
        }
    }

    /// Quorum size for a stable configuration.
    fn quorum(set: &BTreeSet<NodeId>) -> usize { set.len() / 2 + 1 }

    /// Determine whether `acks` represents a quorum, accounting for joint
    /// configs (which require quorum in BOTH old and new).
    pub fn has_quorum(&self, acks: &BTreeSet<NodeId>) -> bool {
        match self {
            Self::Stable(s) => {
                if s.is_empty() { return false; }
                let need = Self::quorum(s);
                s.iter().filter(|id| acks.contains(id)).count() >= need
            }
            Self::Joint { old, new } => {
                let need_old = Self::quorum(old);
                let need_new = Self::quorum(new);
                let got_old = old.iter().filter(|id| acks.contains(id)).count();
                let got_new = new.iter().filter(|id| acks.contains(id)).count();
                got_old >= need_old && got_new >= need_new
            }
        }
    }
}

/// Top-level cluster configuration. Has both the committed state and any
/// pending in-flight transition.
#[derive(Debug, Clone, Default, PartialEq, Eq, Serialize, Deserialize)]
pub struct ClusterConfig {
    /// Currently active config state.
    pub state: ConfigState,
}

impl ClusterConfig {
    /// Construct a stable config from a list of voters.
    pub fn stable<I: IntoIterator<Item = NodeId>>(voters: I) -> Self {
        Self { state: ConfigState::Stable(voters.into_iter().collect()) }
    }
}

/// A configuration-change request from a client.
#[derive(Debug, Clone, PartialEq, Eq, Serialize, Deserialize)]
pub enum ConfigChange {
    /// Add a voting server.
    AddServer(NodeId),
    /// Remove a server.
    RemoveServer(NodeId),
}

#[cfg(test)]
mod tests {
    use super::*;

    fn set(ids: &[NodeId]) -> BTreeSet<NodeId> { ids.iter().copied().collect() }

    #[test]
    fn quorum_stable() {
        let cfg = ConfigState::Stable(set(&[1, 2, 3]));
        assert!(!cfg.has_quorum(&set(&[1])));
        assert!(cfg.has_quorum(&set(&[1, 2])));
        assert!(cfg.has_quorum(&set(&[1, 2, 3])));
    }

    #[test]
    fn quorum_joint_requires_both() {
        let cfg = ConfigState::Joint { old: set(&[1, 2, 3]), new: set(&[3, 4, 5]) };
        // Quorum in old (1,2) but not new (only 3)
        assert!(!cfg.has_quorum(&set(&[1, 2, 3])));
        // Quorum in both
        assert!(cfg.has_quorum(&set(&[1, 2, 3, 4])));
    }
}
\end{lstlisting}

\section{crates/raft-core/src/error.rs}
\begin{lstlisting}[language=Rust]
//! Error types for the consensus core.

use thiserror::Error;

/// Errors a client proposal may produce.
#[derive(Debug, Error, Clone, PartialEq, Eq)]
pub enum ProposeError {
    /// This node is not the leader; the runtime should redirect to the hint.
    #[error("not leader (hint: {hint:?})")]
    NotLeader {
        /// Best-known leader id, if any.
        hint: Option<u64>,
    },
    /// A configuration change is already in flight.
    #[error("configuration change already in progress")]
    ConfigChangeInProgress,
    /// The proposal was rejected because leadership transfer is in progress.
    #[error("leadership transfer in progress")]
    LeadershipTransferInProgress,
    /// Proposal payload exceeds limits.
    #[error("proposal too large: {0} bytes")]
    ProposalTooLarge(usize),
}

/// Top-level core errors (rare; most invariant violations are debug-asserted).
#[derive(Debug, Error, Clone, PartialEq, Eq)]
pub enum RaftError {
    /// Invariant violated (should never occur in correct code).
    #[error("raft invariant violated: {0}")]
    InvariantViolation(String),
    /// Bad input from the runtime.
    #[error("invalid input: {0}")]
    InvalidInput(String),
}
\end{lstlisting}

\section{crates/raft-core/src/lib.rs}
\begin{lstlisting}[language=Rust]
#![deny(unsafe_code)]
#![warn(missing_docs)]
//! Pure, deterministic Raft consensus state machine.
//!
//! This crate is intentionally free of I/O, async, clocks, and OS-level
//! randomness. It exposes a `RaftNode` that consumes [`Message`]s and produces
//! [`Action`]s for a runtime to interpret. This separation is what makes the
//! consensus logic deterministically testable (see TigerBeetle VOPR,
//! FoundationDB simulation, and the etcd/raft pure state machine).
//!
//! Reference: Ongaro & Ousterhout, "In Search of an Understandable Consensus
//! Algorithm", USENIX ATC 2014; Ongaro PhD thesis 2014.

pub mod action;
pub mod config;
pub mod error;
pub mod log;
pub mod message;
pub mod node;
pub mod progress;
pub mod role;
pub mod state;

pub use action::{Action, MetricEvent};
pub use config::{ClusterConfig, ConfigChange, ConfigState};
pub use error::{ProposeError, RaftError};
pub use log::{Entry, EntryKind, LogIndex, RaftLog, Term};
pub use message::Message;
pub use node::{NodeId, RaftConfig, RaftNode};
pub use role::Role;
pub use state::{HardState, SoftState};
\end{lstlisting}

\section{crates/raft-core/src/log.rs}
\begin{lstlisting}[language=Rust]
//! In-memory representation of the Raft log.
//!
//! The pure core keeps the log abstractly. Persistence is the runtime's
//! responsibility (via [`crate::Action::AppendEntries`] and
//! [`crate::Action::TruncateLog`]).

use serde::{Deserialize, Serialize};

/// Raft term number (monotonic).
pub type Term = u64;

/// Raft log index (1-based; 0 is the sentinel "empty log" value).
pub type LogIndex = u64;

/// Kind of log entry — discriminator for both wire and disk encoding.
#[derive(Debug, Clone, Copy, PartialEq, Eq, Serialize, Deserialize)]
pub enum EntryKind {
    /// A no-op entry leaders append on election (see Raft §8 commit safety).
    Noop,
    /// A user/state-machine command.
    Normal,
    /// Joint-consensus configuration entry (`C_old,new`).
    ConfigJoint,
    /// Final new configuration (`C_new`).
    ConfigNew,
}

/// A single Raft log entry.
#[derive(Debug, Clone, PartialEq, Eq, Serialize, Deserialize)]
pub struct Entry {
    /// Term in which the entry was created by the leader.
    pub term: Term,
    /// Position in the log (1-based).
    pub index: LogIndex,
    /// Entry kind.
    pub kind: EntryKind,
    /// Opaque payload — the state machine interprets this for `Normal`,
    /// the config module for `ConfigJoint`/`ConfigNew`.
    pub data: Vec<u8>,
}

impl Entry {
    /// Build a normal command entry.
    pub fn normal(term: Term, index: LogIndex, data: Vec<u8>) -> Self {
        Self { term, index, kind: EntryKind::Normal, data }
    }

    /// Build a no-op entry.
    pub fn noop(term: Term, index: LogIndex) -> Self {
        Self { term, index, kind: EntryKind::Noop, data: Vec::new() }
    }

    /// Approximate serialized size for batching decisions.
    pub fn approx_size(&self) -> usize {
        24 + self.data.len()
    }
}

/// Pure in-memory log. Persistence is handled by the runtime via Actions.
///
/// Invariants:
/// * `entries[i].index == base_index + i + 1`
/// * Indices are strictly monotonic.
/// * Terms are non-decreasing along the log.
#[derive(Debug, Clone, Default)]
pub struct RaftLog {
    /// Entries strictly *after* the snapshot point.
    entries: Vec<Entry>,
    /// Index of the last entry covered by the snapshot. Entries[0].index =
    /// snapshot_last_index + 1.
    snapshot_last_index: LogIndex,
    /// Term of the entry at `snapshot_last_index`.
    snapshot_last_term: Term,
}

impl RaftLog {
    /// Construct an empty log.
    pub fn new() -> Self { Self::default() }

    /// Restore from a snapshot point with no live entries.
    pub fn from_snapshot(last_index: LogIndex, last_term: Term) -> Self {
        Self { entries: Vec::new(), snapshot_last_index: last_index, snapshot_last_term: last_term }
    }

    /// Index of the last entry (or snapshot tip if log is empty).
    pub fn last_index(&self) -> LogIndex {
        self.entries.last().map_or(self.snapshot_last_index, |e| e.index)
    }

    /// Term of the last entry (or snapshot tip).
    pub fn last_term(&self) -> Term {
        self.entries.last().map_or(self.snapshot_last_term, |e| e.term)
    }

    /// First live index (i.e. snapshot_last_index + 1).
    pub fn first_index(&self) -> LogIndex {
        self.snapshot_last_index + 1
    }

    /// Snapshot tip index.
    pub fn snapshot_index(&self) -> LogIndex { self.snapshot_last_index }
    /// Snapshot tip term.
    pub fn snapshot_term(&self) -> Term { self.snapshot_last_term }

    /// Term at a given index, if known.
    pub fn term_at(&self, index: LogIndex) -> Option<Term> {
        if index == 0 {
            return Some(0);
        }
        if index == self.snapshot_last_index {
            return Some(self.snapshot_last_term);
        }
        if index < self.snapshot_last_index || index > self.last_index() {
            return None;
        }
        let off = (index - self.snapshot_last_index - 1) as usize;
        self.entries.get(off).map(|e| e.term)
    }

    /// Borrow entry at `index` (live entries only).
    pub fn get(&self, index: LogIndex) -> Option<&Entry> {
        if index <= self.snapshot_last_index { return None; }
        let off = (index - self.snapshot_last_index - 1) as usize;
        self.entries.get(off)
    }

    /// Slice of entries with indices in `[from, to)` (clamped to live range).
    pub fn slice(&self, from: LogIndex, to: LogIndex) -> &[Entry] {
        if to <= from || from > self.last_index() { return &[]; }
        let lo_idx = from.max(self.snapshot_last_index + 1);
        let hi_idx = to.min(self.last_index() + 1);
        if hi_idx <= lo_idx { return &[]; }
        let lo = (lo_idx - self.snapshot_last_index - 1) as usize;
        let hi = (hi_idx - self.snapshot_last_index - 1) as usize;
        &self.entries[lo..hi]
    }

    /// Append entries in monotonic order. Caller guarantees indices line up.
    pub fn append(&mut self, mut new_entries: Vec<Entry>) {
        if new_entries.is_empty() { return; }
        let expected = self.last_index() + 1;
        debug_assert_eq!(
            new_entries[0].index, expected,
            "append must be contiguous: expected {} got {}", expected, new_entries[0].index
        );
        // Monotonic terms invariant
        if let Some(last) = self.entries.last() {
            debug_assert!(new_entries[0].term >= last.term, "terms must be non-decreasing");
        }
        self.entries.append(&mut new_entries);
    }

    /// Truncate the log so that no entries with index >= `from` remain.
    /// Used when a follower discovers a conflict.
    pub fn truncate_from(&mut self, from: LogIndex) {
        if from <= self.snapshot_last_index {
            // Must not truncate into the snapshot.
            self.entries.clear();
            return;
        }
        let off = (from - self.snapshot_last_index - 1) as usize;
        if off < self.entries.len() {
            self.entries.truncate(off);
        }
    }

    /// Discard entries up to and including `last_included_index`, then set the
    /// snapshot pointer.
    pub fn compact_through(&mut self, last_included_index: LogIndex, last_included_term: Term) {
        if last_included_index <= self.snapshot_last_index {
            return;
        }
        if last_included_index >= self.last_index() {
            self.entries.clear();
        } else {
            let off = (last_included_index - self.snapshot_last_index) as usize;
            self.entries.drain(0..off);
        }
        self.snapshot_last_index = last_included_index;
        self.snapshot_last_term = last_included_term;
    }

    /// Number of live entries.
    pub fn len(&self) -> usize { self.entries.len() }
    /// True if no live entries.
    pub fn is_empty(&self) -> bool { self.entries.is_empty() }

    /// Determine whether a candidate's log is at least as up-to-date as ours
    /// (Raft §5.4.1).
    pub fn is_up_to_date(&self, last_log_index: LogIndex, last_log_term: Term) -> bool {
        let our_term = self.last_term();
        let our_index = self.last_index();
        if last_log_term != our_term {
            last_log_term > our_term
        } else {
            last_log_index >= our_index
        }
    }
}

#[cfg(test)]
mod tests {
    use super::*;

    #[test]
    fn append_and_query() {
        let mut log = RaftLog::new();
        assert_eq!(log.last_index(), 0);
        assert_eq!(log.last_term(), 0);

        log.append(vec![
            Entry::noop(1, 1),
            Entry::normal(1, 2, b"a".to_vec()),
            Entry::normal(2, 3, b"b".to_vec()),
        ]);
        assert_eq!(log.last_index(), 3);
        assert_eq!(log.last_term(), 2);
        assert_eq!(log.get(2).unwrap().data, b"a");
        assert_eq!(log.term_at(2), Some(1));
        assert_eq!(log.term_at(3), Some(2));
        assert_eq!(log.term_at(0), Some(0));
        assert_eq!(log.term_at(99), None);
    }

    #[test]
    fn truncate_from() {
        let mut log = RaftLog::new();
        log.append(vec![
            Entry::normal(1, 1, vec![1]),
            Entry::normal(1, 2, vec![2]),
            Entry::normal(1, 3, vec![3]),
        ]);
        log.truncate_from(2);
        assert_eq!(log.last_index(), 1);
        assert!(log.get(2).is_none());
    }

    #[test]
    fn up_to_date_check() {
        let mut log = RaftLog::new();
        log.append(vec![Entry::normal(2, 1, vec![1])]);
        // Higher term wins regardless of index.
        assert!(log.is_up_to_date(0, 3));
        // Same term, lower index loses.
        assert!(!log.is_up_to_date(0, 2));
        // Same (term, index) is a tie -> up to date.
        assert!(log.is_up_to_date(1, 2));
        // Lower term loses.
        assert!(!log.is_up_to_date(99, 1));
    }

    #[test]
    fn compaction_preserves_tail() {
        let mut log = RaftLog::new();
        log.append(vec![
            Entry::normal(1, 1, vec![1]),
            Entry::normal(1, 2, vec![2]),
            Entry::normal(1, 3, vec![3]),
            Entry::normal(1, 4, vec![4]),
        ]);
        log.compact_through(2, 1);
        assert_eq!(log.first_index(), 3);
        assert_eq!(log.last_index(), 4);
        assert_eq!(log.snapshot_index(), 2);
        assert_eq!(log.term_at(2), Some(1));
        assert_eq!(log.get(3).unwrap().data, vec![3]);
    }
}
\end{lstlisting}

\section{crates/raft-core/src/message.rs}
\begin{lstlisting}[language=Rust]
//! Internal message types — the alphabet the state machine consumes.
//!
//! These mirror but are independent from the wire `proto/raft.proto` types so
//! the core doesn't depend on prost/tonic.

use bytes::Bytes;

use crate::log::{Entry, LogIndex, Term};
use crate::node::NodeId;

/// Identifier for correlating a request with its response in tracing/metrics.
pub type RequestId = u64;

/// Internal Raft messages.
#[derive(Debug, Clone)]
pub enum Message {
    /// AppendEntries (with optional heartbeat = empty entries).
    AppendEntries {
        /// Sender (leader).
        from: NodeId,
        /// Leader's term.
        term: Term,
        /// Index immediately preceding the new entries.
        prev_log_index: LogIndex,
        /// Term of `prev_log_index`.
        prev_log_term: Term,
        /// New entries (may be empty for heartbeat).
        entries: Vec<Entry>,
        /// Leader's commit index.
        leader_commit: LogIndex,
        /// Trace correlation id.
        request_id: RequestId,
    },
    /// AppendEntries response.
    AppendEntriesResponse {
        /// Responder.
        from: NodeId,
        /// Responder's current term.
        term: Term,
        /// Whether the entries were accepted.
        success: bool,
        /// Hint for fast back-off (Ongaro thesis §3.5).
        conflict_index: LogIndex,
        /// Conflict term hint.
        conflict_term: Term,
        /// Responder's last log index — used to advance matchIndex on success.
        last_log_index: LogIndex,
        /// Trace correlation id.
        request_id: RequestId,
    },
    /// RequestVote (or PreVote when `pre_vote` is true).
    RequestVote {
        /// Sender (candidate).
        from: NodeId,
        /// Candidate's term (or proposed term for pre-vote).
        term: Term,
        /// Candidate's last log index.
        last_log_index: LogIndex,
        /// Candidate's last log term.
        last_log_term: Term,
        /// True if this is a PreVote (does not bump receiver's term).
        pre_vote: bool,
        /// Trace correlation id.
        request_id: RequestId,
    },
    /// RequestVote response.
    RequestVoteResponse {
        /// Responder.
        from: NodeId,
        /// Responder's term.
        term: Term,
        /// Whether the vote was granted.
        vote_granted: bool,
        /// Whether this corresponds to a pre-vote.
        pre_vote: bool,
        /// Trace correlation id.
        request_id: RequestId,
    },
    /// InstallSnapshot RPC (chunked transfer is the runtime's responsibility;
    /// the core sees a single logical message per snapshot).
    InstallSnapshot {
        /// Sender (leader).
        from: NodeId,
        /// Leader's term.
        term: Term,
        /// Last index covered by the snapshot.
        last_included_index: LogIndex,
        /// Term at `last_included_index`.
        last_included_term: Term,
        /// Snapshot bytes (state machine + config).
        data: Bytes,
        /// Trace correlation id.
        request_id: RequestId,
    },
    /// InstallSnapshot response.
    InstallSnapshotResponse {
        /// Responder.
        from: NodeId,
        /// Responder's term.
        term: Term,
        /// Trace correlation id.
        request_id: RequestId,
    },
    /// Leadership transfer trigger (Ongaro thesis §3.10).
    TimeoutNow {
        /// Sender (current leader).
        from: NodeId,
        /// Sender's term.
        term: Term,
    },
}
\end{lstlisting}

\section{crates/raft-core/src/node.rs}
\begin{lstlisting}[language=Rust]
//! The Raft state machine.
//!
//! Pure: no I/O, no async, no system clock, no OS RNG. The runtime drives it
//! with messages and an external monotonic time value (`u64` ticks).

use std::collections::{BTreeMap, BTreeSet};

use bytes::Bytes;
use rand::SeedableRng;
use rand::Rng;
use rand_chacha::ChaCha20Rng;
use smallvec::SmallVec;

use crate::action::{Action, MetricEvent};
use crate::config::{ClusterConfig, ConfigChange, ConfigState};
use crate::error::ProposeError;
use crate::log::{Entry, EntryKind, LogIndex, RaftLog, Term};
use crate::message::{Message, RequestId};
use crate::progress::{Progress, ProgressSet, ProgressState};
use crate::role::Role;
use crate::state::{HardState, SoftState};

/// Cluster-unique node identifier.
pub type NodeId = u64;

/// Configuration for a `RaftNode`.
#[derive(Debug, Clone)]
pub struct RaftConfig {
    /// This node's id.
    pub id: NodeId,
    /// Initial cluster membership.
    pub initial_voters: Vec<NodeId>,
    /// Election timeout in ticks (the runtime decides what a tick is; typically
    /// 1ms). The actual timeout is randomized in
    /// `[election_timeout_ticks, 2 * election_timeout_ticks)`.
    pub election_timeout_ticks: u64,
    /// Heartbeat interval in ticks.
    pub heartbeat_ticks: u64,
    /// Maximum number of entries in a single AppendEntries.
    pub max_append_entries: usize,
    /// Maximum number of in-flight pipelined AEs per peer.
    pub max_inflight_msgs: u32,
    /// PRNG seed (must be deterministic for simulation tests).
    pub rng_seed: u64,
    /// Enable pre-vote (Ongaro thesis §9.6).
    pub pre_vote: bool,
}

impl Default for RaftConfig {
    fn default() -> Self {
        Self {
            id: 1,
            initial_voters: vec![1, 2, 3],
            election_timeout_ticks: 150,
            heartbeat_ticks: 15,
            max_append_entries: 64,
            max_inflight_msgs: 256,
            rng_seed: 0xDEAD_BEEF,
            pre_vote: true,
        }
    }
}

/// The Raft state machine.
pub struct RaftNode {
    /// Local configuration.
    cfg: RaftConfig,
    /// Persisted state.
    hs: HardState,
    /// In-memory log.
    log: RaftLog,
    /// Volatile commit/apply pointers (commit_index lives in HardState).
    last_applied: LogIndex,
    /// Current role.
    role: Role,
    /// Best-known leader for the current term.
    leader_id: Option<NodeId>,
    /// Cluster configuration (latest seen, possibly uncommitted).
    config: ClusterConfig,
    /// Per-peer progress (only meaningful as Leader).
    progress: ProgressSet,
    /// Votes received in the current election (RequestVote / PreVote).
    votes: BTreeMap<NodeId, bool>,
    /// Tick counter (logical time).
    elapsed_election_ticks: u64,
    elapsed_heartbeat_ticks: u64,
    /// Randomized election timeout for the *current* term.
    randomized_election_timeout: u64,
    /// Deterministic PRNG.
    rng: ChaCha20Rng,
    /// Counter for outgoing request ids.
    next_request_id: u64,
    /// Pending leadership transfer target.
    leader_transfer_target: Option<NodeId>,
    /// Whether `C_old,new` has been observed but not yet superseded by `C_new`.
    in_joint: bool,
}

impl RaftNode {
    /// Construct a fresh node from durable state. If `hs` is default and `log`
    /// is empty, this is a brand-new cluster member.
    pub fn new(cfg: RaftConfig, hs: HardState, log: RaftLog) -> Self {
        let rng = ChaCha20Rng::seed_from_u64(cfg.rng_seed.wrapping_add(cfg.id));
        let mut me = Self {
            config: ClusterConfig::stable(cfg.initial_voters.iter().copied()),
            cfg,
            hs,
            log,
            last_applied: 0,
            role: Role::Follower,
            leader_id: None,
            progress: ProgressSet::default(),
            votes: BTreeMap::new(),
            elapsed_election_ticks: 0,
            elapsed_heartbeat_ticks: 0,
            randomized_election_timeout: 0,
            rng,
            next_request_id: 1,
            leader_transfer_target: None,
            in_joint: false,
        };
        me.reset_randomized_election_timeout();
        // last_applied >= snapshot index (we never re-apply through a snapshot)
        me.last_applied = me.log.snapshot_index();
        me
    }

    /// Soft (volatile) state snapshot.
    pub fn soft_state(&self) -> SoftState {
        SoftState { role: self.role, leader_id: self.leader_id }
    }

    /// Borrow the hard (persistent) state.
    pub fn hard_state(&self) -> &HardState { &self.hs }

    /// Borrow the cluster config.
    pub fn config(&self) -> &ClusterConfig { &self.config }

    /// This node's id.
    pub fn id(&self) -> NodeId { self.cfg.id }

    /// Current role.
    pub fn role(&self) -> Role { self.role }

    /// Last log index.
    pub fn last_log_index(&self) -> LogIndex { self.log.last_index() }

    /// Commit index.
    pub fn commit_index(&self) -> LogIndex { self.hs.commit_index }

    /// Last applied index.
    pub fn last_applied(&self) -> LogIndex { self.last_applied }

    fn next_rid(&mut self) -> RequestId {
        let r = self.next_request_id;
        self.next_request_id = self.next_request_id.wrapping_add(1);
        r
    }

    fn reset_randomized_election_timeout(&mut self) {
        let lo = self.cfg.election_timeout_ticks;
        let hi = lo * 2;
        self.randomized_election_timeout = self.rng.gen_range(lo..hi);
    }

    fn voters(&self) -> Vec<NodeId> { self.config.state.voters() }

    fn is_voter_self(&self) -> bool { self.config.state.is_voter(self.cfg.id) }

    // ===================================================================
    // Public API
    // ===================================================================

    /// Drive one tick of logical time. Runtime calls this on a fixed cadence.
    pub fn tick(&mut self) -> SmallVec<[Action; 8]> {
        let mut acts: SmallVec<[Action; 8]> = SmallVec::new();
        match self.role {
            Role::Leader => {
                self.elapsed_heartbeat_ticks += 1;
                if self.elapsed_heartbeat_ticks >= self.cfg.heartbeat_ticks {
                    self.elapsed_heartbeat_ticks = 0;
                    self.broadcast_heartbeat(&mut acts);
                }
            }
            Role::Follower | Role::Candidate | Role::PreCandidate => {
                self.elapsed_election_ticks += 1;
                if self.elapsed_election_ticks >= self.randomized_election_timeout {
                    self.elapsed_election_ticks = 0;
                    self.start_election(&mut acts);
                }
            }
        }
        acts
    }

    /// Process a Raft message from a peer.
    pub fn step(&mut self, msg: Message) -> SmallVec<[Action; 8]> {
        let mut acts: SmallVec<[Action; 8]> = SmallVec::new();
        // Universal term-bump rule (Raft §5.1): any RPC with a higher term
        // immediately reverts us to follower.
        let msg_term = match &msg {
            Message::AppendEntries { term, .. }
            | Message::AppendEntriesResponse { term, .. }
            | Message::RequestVote { term, pre_vote: false, .. }
            | Message::RequestVoteResponse { term, pre_vote: false, .. }
            | Message::InstallSnapshot { term, .. }
            | Message::InstallSnapshotResponse { term, .. }
            | Message::TimeoutNow { term, .. } => *term,
            // PreVote messages do NOT bump terms.
            Message::RequestVote { term: _, pre_vote: true, .. }
            | Message::RequestVoteResponse { term: _, pre_vote: true, .. } => 0,
        };
        if msg_term > self.hs.current_term {
            let leader_hint = match &msg {
                Message::AppendEntries { from, .. } | Message::InstallSnapshot { from, .. } => Some(*from),
                _ => None,
            };
            self.become_follower(msg_term, leader_hint, &mut acts);
        }

        match msg {
            Message::AppendEntries { from, term, prev_log_index, prev_log_term, entries, leader_commit, request_id } => {
                self.handle_append_entries(from, term, prev_log_index, prev_log_term, entries, leader_commit, request_id, &mut acts);
            }
            Message::AppendEntriesResponse { from, term, success, conflict_index, conflict_term: _, last_log_index, request_id: _ } => {
                self.handle_append_entries_response(from, term, success, conflict_index, last_log_index, &mut acts);
            }
            Message::RequestVote { from, term, last_log_index, last_log_term, pre_vote, request_id } => {
                self.handle_request_vote(from, term, last_log_index, last_log_term, pre_vote, request_id, &mut acts);
            }
            Message::RequestVoteResponse { from, term, vote_granted, pre_vote, request_id: _ } => {
                self.handle_vote_response(from, term, vote_granted, pre_vote, &mut acts);
            }
            Message::InstallSnapshot { from, term, last_included_index, last_included_term, data, request_id } => {
                self.handle_install_snapshot(from, term, last_included_index, last_included_term, data, request_id, &mut acts);
            }
            Message::InstallSnapshotResponse { from, term: _, request_id: _ } => {
                if let Some(p) = self.progress.get_mut(from) {
                    if let Some(idx) = p.pending_snapshot {
                        p.snapshot_finished(idx);
                    }
                }
            }
            Message::TimeoutNow { from: _, term: _ } => {
                if matches!(self.role, Role::Follower) && self.is_voter_self() {
                    self.start_election(&mut acts);
                }
            }
        }
        acts
    }

    /// Propose a normal entry. Only valid as leader.
    pub fn propose(&mut self, data: Vec<u8>) -> Result<SmallVec<[Action; 8]>, ProposeError> {
        if !matches!(self.role, Role::Leader) {
            return Err(ProposeError::NotLeader { hint: self.leader_id });
        }
        if self.leader_transfer_target.is_some() {
            return Err(ProposeError::LeadershipTransferInProgress);
        }
        let mut acts = SmallVec::new();
        let idx = self.log.last_index() + 1;
        let entry = Entry::normal(self.hs.current_term, idx, data);
        self.log.append(vec![entry.clone()]);
        acts.push(Action::AppendEntries(vec![entry]));
        // Leader trivially has match_index = last_log_index for self.
        if let Some(p) = self.progress.get_mut(self.cfg.id) {
            p.match_index = self.log.last_index();
            p.next_index = self.log.last_index() + 1;
        }
        self.maybe_advance_commit(&mut acts);
        self.broadcast_replication(&mut acts);
        Ok(acts)
    }

    /// Propose a configuration change. Goes through joint consensus.
    pub fn propose_config_change(&mut self, change: ConfigChange) -> Result<SmallVec<[Action; 8]>, ProposeError> {
        if !matches!(self.role, Role::Leader) {
            return Err(ProposeError::NotLeader { hint: self.leader_id });
        }
        if self.in_joint {
            return Err(ProposeError::ConfigChangeInProgress);
        }
        let mut new_voters: BTreeSet<NodeId> = match &self.config.state {
            ConfigState::Stable(s) => s.clone(),
            ConfigState::Joint { new, .. } => new.clone(),
        };
        match change {
            ConfigChange::AddServer(id) => { new_voters.insert(id); }
            ConfigChange::RemoveServer(id) => { new_voters.remove(&id); }
        }
        let old_voters: BTreeSet<NodeId> = match &self.config.state {
            ConfigState::Stable(s) => s.clone(),
            ConfigState::Joint { old, .. } => old.clone(),
        };
        let joint_state = ConfigState::Joint { old: old_voters, new: new_voters };
        let payload = bincode::serialize(&joint_state).expect("config encode");
        let idx = self.log.last_index() + 1;
        let entry = Entry { term: self.hs.current_term, index: idx, kind: EntryKind::ConfigJoint, data: payload };
        self.log.append(vec![entry.clone()]);
        self.apply_config_entry(&entry);
        let mut acts = SmallVec::new();
        acts.push(Action::AppendEntries(vec![entry]));
        if let Some(p) = self.progress.get_mut(self.cfg.id) {
            p.match_index = self.log.last_index();
            p.next_index = self.log.last_index() + 1;
        }
        self.broadcast_replication(&mut acts);
        Ok(acts)
    }

    /// Begin a linearizable read-index round. Returns actions which include
    /// heartbeats whose acks the runtime correlates with the read.
    pub fn read_index(&mut self, ctx: Bytes) -> SmallVec<[Action; 8]> {
        let mut acts = SmallVec::new();
        if !matches!(self.role, Role::Leader) {
            return acts;
        }
        let rid = self.next_rid();
        // First piggyback the request_id by issuing fresh heartbeats.
        // The runtime correlates AE responses to count quorum confirmation
        // (a simplification: any quorum ack of *any* heartbeat after this
        // point confirms commit_index).
        self.broadcast_heartbeat(&mut acts);
        acts.push(Action::NotifyReadIndex {
            ctx,
            commit_index: self.hs.commit_index,
            request_id: rid,
        });
        acts
    }

    /// Initiate leadership transfer to `target` (Ongaro thesis §3.10).
    pub fn transfer_leadership(&mut self, target: NodeId) -> SmallVec<[Action; 8]> {
        let mut acts = SmallVec::new();
        if !matches!(self.role, Role::Leader) || target == self.cfg.id {
            return acts;
        }
        if !self.config.state.is_voter(target) { return acts; }
        self.leader_transfer_target = Some(target);
        acts.push(Action::Metric(MetricEvent::LeadershipTransferStarted { target }));
        // Send TimeoutNow only if the target is caught up; otherwise wait.
        let caught_up = self.progress.get(target).is_some_and(|p| p.match_index == self.log.last_index());
        if caught_up {
            acts.push(Action::SendMessage {
                to: target,
                msg: Message::TimeoutNow { from: self.cfg.id, term: self.hs.current_term },
            });
        }
        acts
    }

    // ===================================================================
    // Election logic
    // ===================================================================

    fn become_follower(&mut self, term: Term, leader: Option<NodeId>, acts: &mut SmallVec<[Action; 8]>) {
        let term_changed = term > self.hs.current_term;
        if term_changed {
            self.hs.current_term = term;
            self.hs.voted_for = None;
            acts.push(Action::Metric(MetricEvent::TermBumped { new_term: term }));
        }
        self.role = Role::Follower;
        self.leader_id = leader;
        self.votes.clear();
        self.elapsed_election_ticks = 0;
        self.elapsed_heartbeat_ticks = 0;
        self.leader_transfer_target = None;
        self.reset_randomized_election_timeout();
        acts.push(Action::PersistHardState(self.hs.clone()));
        acts.push(Action::ResetElectionTimer);
        acts.push(Action::BecameFollower { term: self.hs.current_term, leader });
    }

    fn become_pre_candidate(&mut self, acts: &mut SmallVec<[Action; 8]>) {
        self.role = Role::PreCandidate;
        self.leader_id = None;
        self.votes.clear();
        // Pre-candidate votes count against `current_term + 1` but no state is
        // persisted yet.
        self.votes.insert(self.cfg.id, true);
        acts.push(Action::Metric(MetricEvent::ElectionStarted {
            term: self.hs.current_term + 1, pre_vote: true,
        }));
    }

    fn become_candidate(&mut self, acts: &mut SmallVec<[Action; 8]>) {
        self.hs.current_term += 1;
        self.hs.voted_for = Some(self.cfg.id);
        self.role = Role::Candidate;
        self.leader_id = None;
        self.votes.clear();
        self.votes.insert(self.cfg.id, true);
        self.elapsed_election_ticks = 0;
        self.reset_randomized_election_timeout();
        acts.push(Action::Metric(MetricEvent::TermBumped { new_term: self.hs.current_term }));
        acts.push(Action::Metric(MetricEvent::ElectionStarted { term: self.hs.current_term, pre_vote: false }));
        acts.push(Action::PersistHardState(self.hs.clone()));
    }

    fn become_leader(&mut self, acts: &mut SmallVec<[Action; 8]>) {
        self.role = Role::Leader;
        self.leader_id = Some(self.cfg.id);
        self.elapsed_heartbeat_ticks = 0;
        self.leader_transfer_target = None;

        let last_idx = self.log.last_index();
        // Reset progress for all current voters.
        let voters = self.voters();
        self.progress.peers.clear();
        for &id in &voters {
            let mut p = Progress::new(last_idx);
            if id == self.cfg.id {
                p.match_index = last_idx;
                p.next_index = last_idx + 1;
                p.state = ProgressState::Replicate;
            }
            self.progress.peers.insert(id, p);
        }

        // Append a no-op entry to commit any prior-term entries (Figure 8).
        let idx = last_idx + 1;
        let noop = Entry::noop(self.hs.current_term, idx);
        self.log.append(vec![noop.clone()]);
        if let Some(p) = self.progress.get_mut(self.cfg.id) {
            p.match_index = self.log.last_index();
            p.next_index = self.log.last_index() + 1;
        }
        acts.push(Action::AppendEntries(vec![noop]));
        acts.push(Action::BecameLeader { term: self.hs.current_term });
        acts.push(Action::Metric(MetricEvent::ElectionWon { term: self.hs.current_term }));
        acts.push(Action::ResetHeartbeatTimer);
        self.broadcast_replication(acts);
    }

    fn start_election(&mut self, acts: &mut SmallVec<[Action; 8]>) {
        if !self.is_voter_self() {
            return;
        }
        if self.cfg.pre_vote {
            self.become_pre_candidate(acts);
        } else {
            self.become_candidate(acts);
        }
        let term = if matches!(self.role, Role::PreCandidate) {
            self.hs.current_term + 1
        } else {
            self.hs.current_term
        };
        let last_log_index = self.log.last_index();
        let last_log_term = self.log.last_term();
        for peer in self.voters() {
            if peer == self.cfg.id { continue; }
            let rid = self.next_rid();
            acts.push(Action::SendMessage {
                to: peer,
                msg: Message::RequestVote {
                    from: self.cfg.id,
                    term,
                    last_log_index,
                    last_log_term,
                    pre_vote: matches!(self.role, Role::PreCandidate),
                    request_id: rid,
                },
            });
        }
        // If the cluster is a single voter (this node), trivially win.
        if self.voters().len() == 1 && self.config.state.is_voter(self.cfg.id) {
            if matches!(self.role, Role::PreCandidate) {
                self.become_candidate(acts);
            }
            self.become_leader(acts);
        }
    }

    fn handle_request_vote(
        &mut self,
        from: NodeId,
        term: Term,
        last_log_index: LogIndex,
        last_log_term: Term,
        pre_vote: bool,
        request_id: RequestId,
        acts: &mut SmallVec<[Action; 8]>,
    ) {
        let our_term = self.hs.current_term;
        // Pre-vote request term is the candidate's *proposed* term, not bumped.
        let effective_term = if pre_vote { term.saturating_sub(1) } else { term };
        let mut grant = false;

        if !pre_vote {
            // Standard RequestVote (term already handled by universal bump).
            if term < our_term {
                grant = false;
            } else {
                let can_vote = match self.hs.voted_for {
                    None => true,
                    Some(v) => v == from,
                };
                if can_vote && self.log.is_up_to_date(last_log_index, last_log_term) {
                    self.hs.voted_for = Some(from);
                    self.elapsed_election_ticks = 0;
                    grant = true;
                    acts.push(Action::PersistHardState(self.hs.clone()));
                }
            }
        } else {
            // PreVote: do not change state; only grant if we'd vote in `term`
            // and our election timer is at risk of firing.
            let leader_lease_active = self.elapsed_election_ticks < self.cfg.election_timeout_ticks
                && self.leader_id.is_some();
            if !leader_lease_active && term > our_term && self.log.is_up_to_date(last_log_index, last_log_term) {
                grant = true;
            }
            // Use the candidate's term verbatim in the response so the
            // candidate can detect competitive races.
            let _ = effective_term;
        }

        let resp_term = if pre_vote { term } else { self.hs.current_term };
        acts.push(Action::SendMessage {
            to: from,
            msg: Message::RequestVoteResponse {
                from: self.cfg.id,
                term: resp_term,
                vote_granted: grant,
                pre_vote,
                request_id,
            },
        });
    }

    fn handle_vote_response(
        &mut self,
        from: NodeId,
        term: Term,
        vote_granted: bool,
        pre_vote: bool,
        acts: &mut SmallVec<[Action; 8]>,
    ) {
        let expected_role = if pre_vote { Role::PreCandidate } else { Role::Candidate };
        if self.role != expected_role { return; }
        let expected_term = if pre_vote { self.hs.current_term + 1 } else { self.hs.current_term };
        if term != expected_term { return; }

        self.votes.insert(from, vote_granted);
        let granted_set: BTreeSet<NodeId> = self.votes.iter().filter_map(|(k, v)| v.then_some(*k)).collect();
        let denied = self.votes.values().filter(|v| !**v).count();

        if self.config.state.has_quorum(&granted_set) {
            // Won.
            if pre_vote {
                self.become_candidate(acts);
                let last_log_index = self.log.last_index();
                let last_log_term = self.log.last_term();
                for peer in self.voters() {
                    if peer == self.cfg.id { continue; }
                    let rid = self.next_rid();
                    acts.push(Action::SendMessage {
                        to: peer,
                        msg: Message::RequestVote {
                            from: self.cfg.id,
                            term: self.hs.current_term,
                            last_log_index,
                            last_log_term,
                            pre_vote: false,
                            request_id: rid,
                        },
                    });
                }
            } else {
                self.become_leader(acts);
            }
        } else if denied > self.voters().len() / 2 {
            // Cannot win — revert to follower.
            self.become_follower(self.hs.current_term, None, acts);
        }
    }

    // ===================================================================
    // Replication logic
    // ===================================================================

    fn broadcast_heartbeat(&mut self, acts: &mut SmallVec<[Action; 8]>) {
        for peer in self.voters() {
            if peer == self.cfg.id { continue; }
            self.send_append(peer, true, acts);
        }
    }

    fn broadcast_replication(&mut self, acts: &mut SmallVec<[Action; 8]>) {
        let voters = self.voters();
        for peer in voters {
            if peer == self.cfg.id { continue; }
            self.send_append(peer, false, acts);
        }
    }

    fn send_append(&mut self, peer: NodeId, force_heartbeat: bool, acts: &mut SmallVec<[Action; 8]>) {
        let max_inflight = self.cfg.max_inflight_msgs;
        let last_idx = self.log.last_index();
        let max_batch = self.cfg.max_append_entries;
        let snapshot_first = self.log.first_index();
        let snapshot_index = self.log.snapshot_index();
        let snapshot_term = self.log.snapshot_term();

        let progress = match self.progress.get_mut(peer) {
            Some(p) => p,
            None => return,
        };
        if !progress.can_send(max_inflight) && !force_heartbeat { return; }

        // Need snapshot?
        if progress.next_index < snapshot_first {
            progress.become_snapshot(snapshot_index);
            acts.push(Action::SendMessage {
                to: peer,
                msg: Message::InstallSnapshot {
                    from: self.cfg.id,
                    term: self.hs.current_term,
                    last_included_index: snapshot_index,
                    last_included_term: snapshot_term,
                    // The runtime fills the actual snapshot bytes; here we send a marker.
                    data: Bytes::new(),
                    request_id: 0,
                },
            });
            return;
        }

        let next = progress.next_index;
        let prev_index = next.saturating_sub(1);
        let prev_term = if prev_index == 0 {
            0
        } else if prev_index == snapshot_index {
            snapshot_term
        } else {
            match self.log.term_at(prev_index) {
                Some(t) => t,
                None => {
                    // Edge case: we no longer have prev_index — fall back to snapshot.
                    progress.become_snapshot(snapshot_index);
                    acts.push(Action::SendMessage {
                        to: peer,
                        msg: Message::InstallSnapshot {
                            from: self.cfg.id,
                            term: self.hs.current_term,
                            last_included_index: snapshot_index,
                            last_included_term: snapshot_term,
                            data: Bytes::new(),
                            request_id: 0,
                        },
                    });
                    return;
                }
            }
        };

        let entries = if force_heartbeat {
            Vec::new()
        } else if next > last_idx {
            Vec::new()
        } else {
            let to = (next + max_batch as u64).min(last_idx + 1);
            self.log.slice(next, to).to_vec()
        };
        progress.inflight = progress.inflight.saturating_add(1);

        let term = self.hs.current_term;
        let leader_commit = self.hs.commit_index;
        let id = self.cfg.id;
        let rid = self.next_rid();
        acts.push(Action::SendMessage {
            to: peer,
            msg: Message::AppendEntries {
                from: id,
                term,
                prev_log_index: prev_index,
                prev_log_term: prev_term,
                entries,
                leader_commit,
                request_id: rid,
            },
        });
    }

    #[allow(clippy::too_many_arguments)]
    fn handle_append_entries(
        &mut self,
        from: NodeId,
        term: Term,
        prev_log_index: LogIndex,
        prev_log_term: Term,
        entries: Vec<Entry>,
        leader_commit: LogIndex,
        request_id: RequestId,
        acts: &mut SmallVec<[Action; 8]>,
    ) {
        // Term older than ours: reject.
        if term < self.hs.current_term {
            acts.push(Action::SendMessage {
                to: from,
                msg: Message::AppendEntriesResponse {
                    from: self.cfg.id,
                    term: self.hs.current_term,
                    success: false,
                    conflict_index: 0,
                    conflict_term: 0,
                    last_log_index: self.log.last_index(),
                    request_id,
                },
            });
            return;
        }

        // Same/higher term -> recognize sender as leader for this term.
        self.leader_id = Some(from);
        if !matches!(self.role, Role::Follower) {
            self.role = Role::Follower;
        }
        self.elapsed_election_ticks = 0;

        // Log matching: prev_log_index must exist with matching term.
        let our_prev_term = if prev_log_index == 0 { Some(0) } else { self.log.term_at(prev_log_index) };
        let matches = our_prev_term == Some(prev_log_term);
        if !matches {
            // Construct fast back-off hint.
            let (conflict_index, conflict_term) = if prev_log_index > self.log.last_index() {
                (self.log.last_index() + 1, 0)
            } else {
                let t = self.log.term_at(prev_log_index).unwrap_or(0);
                // Find earliest index in `t`
                let mut ci = prev_log_index;
                while ci > self.log.first_index() {
                    if self.log.term_at(ci - 1) != Some(t) { break; }
                    ci -= 1;
                }
                (ci, t)
            };
            acts.push(Action::Metric(MetricEvent::AppendRejected { from: self.cfg.id }));
            acts.push(Action::SendMessage {
                to: from,
                msg: Message::AppendEntriesResponse {
                    from: self.cfg.id,
                    term: self.hs.current_term,
                    success: false,
                    conflict_index,
                    conflict_term,
                    last_log_index: self.log.last_index(),
                    request_id,
                },
            });
            return;
        }

        // Append / overwrite: walk new entries, find the first conflict, then
        // truncate and append.
        if !entries.is_empty() {
            let mut conflict_at: Option<LogIndex> = None;
            for e in &entries {
                match self.log.term_at(e.index) {
                    Some(existing) if existing == e.term => continue,
                    Some(_) => { conflict_at = Some(e.index); break; }
                    None => { conflict_at = Some(e.index); break; }
                }
            }
            if let Some(ci) = conflict_at {
                if ci <= self.log.last_index() {
                    self.log.truncate_from(ci);
                    acts.push(Action::TruncateLog { from: ci });
                }
                let to_append: Vec<Entry> = entries.into_iter().filter(|e| e.index >= ci).collect();
                if !to_append.is_empty() {
                    // Apply config entries observed by follower.
                    for e in &to_append {
                        if matches!(e.kind, EntryKind::ConfigJoint | EntryKind::ConfigNew) {
                            self.apply_config_entry(e);
                        }
                    }
                    self.log.append(to_append.clone());
                    acts.push(Action::AppendEntries(to_append));
                }
            }
        }

        // Advance commit index.
        let new_commit = leader_commit.min(self.log.last_index());
        if new_commit > self.hs.commit_index {
            self.hs.commit_index = new_commit;
            acts.push(Action::Metric(MetricEvent::CommitAdvanced { index: new_commit }));
            acts.push(Action::PersistHardState(self.hs.clone()));
            self.flush_apply(acts);
        }

        acts.push(Action::SendMessage {
            to: from,
            msg: Message::AppendEntriesResponse {
                from: self.cfg.id,
                term: self.hs.current_term,
                success: true,
                conflict_index: 0,
                conflict_term: 0,
                last_log_index: self.log.last_index(),
                request_id,
            },
        });
    }

    fn handle_append_entries_response(
        &mut self,
        from: NodeId,
        term: Term,
        success: bool,
        conflict_index: LogIndex,
        last_log_index: LogIndex,
        acts: &mut SmallVec<[Action; 8]>,
    ) {
        if !matches!(self.role, Role::Leader) || term != self.hs.current_term {
            return;
        }
        let Some(progress) = self.progress.get_mut(from) else { return; };
        progress.recent_active = true;
        progress.inflight = progress.inflight.saturating_sub(1);

        if success {
            progress.maybe_update(last_log_index);
            // Possibly advance commit.
            self.maybe_advance_commit(acts);
            // Continue replicating in case there is more.
            self.send_append(from, false, acts);
            // Leadership transfer hand-off.
            if let Some(target) = self.leader_transfer_target {
                if target == from {
                    if let Some(p) = self.progress.get(target) {
                        if p.match_index == self.log.last_index() {
                            acts.push(Action::SendMessage {
                                to: target,
                                msg: Message::TimeoutNow { from: self.cfg.id, term: self.hs.current_term },
                            });
                        }
                    }
                }
            }
        } else {
            progress.maybe_decr_to(conflict_index);
            self.send_append(from, false, acts);
        }
    }

    fn maybe_advance_commit(&mut self, acts: &mut SmallVec<[Action; 8]>) {
        // Only commit entries from current term (Figure 8 caveat).
        let mut indices: Vec<LogIndex> = self
            .progress
            .iter()
            .filter(|(id, _)| self.config.state.is_voter(**id))
            .map(|(_, p)| p.match_index)
            .collect();
        if indices.is_empty() { return; }
        indices.sort_unstable();
        // For joint consensus, require quorum across BOTH sets independently.
        let new_commit = match &self.config.state {
            ConfigState::Stable(s) => {
                if s.is_empty() { return; }
                let mut sorted: Vec<LogIndex> = s
                    .iter()
                    .filter_map(|id| self.progress.get(*id).map(|p| p.match_index))
                    .collect();
                sorted.sort_unstable();
                let mid = (sorted.len() - 1) / 2;
                sorted[sorted.len() - 1 - mid]
            }
            ConfigState::Joint { old, new } => {
                let mut so: Vec<LogIndex> = old.iter().filter_map(|id| self.progress.get(*id).map(|p| p.match_index)).collect();
                let mut sn: Vec<LogIndex> = new.iter().filter_map(|id| self.progress.get(*id).map(|p| p.match_index)).collect();
                if so.is_empty() || sn.is_empty() { return; }
                so.sort_unstable();
                sn.sort_unstable();
                let mo = so[so.len() - 1 - (so.len() - 1) / 2];
                let mn = sn[sn.len() - 1 - (sn.len() - 1) / 2];
                mo.min(mn)
            }
        };

        if new_commit > self.hs.commit_index {
            // Only commit if entry@new_commit is from current term.
            if self.log.term_at(new_commit) == Some(self.hs.current_term) {
                self.hs.commit_index = new_commit;
                acts.push(Action::Metric(MetricEvent::CommitAdvanced { index: new_commit }));
                acts.push(Action::PersistHardState(self.hs.clone()));
                self.flush_apply(acts);
                self.maybe_finish_joint(acts);
            }
        }
    }

    fn flush_apply(&mut self, acts: &mut SmallVec<[Action; 8]>) {
        if self.hs.commit_index > self.last_applied {
            let from = self.last_applied + 1;
            let to = self.hs.commit_index + 1;
            let entries = self.log.slice(from, to).to_vec();
            if !entries.is_empty() {
                self.last_applied = self.hs.commit_index;
                acts.push(Action::ApplyCommitted { entries });
            }
        }
    }

    // ===================================================================
    // Snapshot handling
    // ===================================================================

    #[allow(clippy::too_many_arguments)]
    fn handle_install_snapshot(
        &mut self,
        from: NodeId,
        term: Term,
        last_included_index: LogIndex,
        last_included_term: Term,
        data: Bytes,
        request_id: RequestId,
        acts: &mut SmallVec<[Action; 8]>,
    ) {
        if term < self.hs.current_term {
            acts.push(Action::SendMessage {
                to: from,
                msg: Message::InstallSnapshotResponse { from: self.cfg.id, term: self.hs.current_term, request_id },
            });
            return;
        }
        self.leader_id = Some(from);
        self.role = Role::Follower;
        self.elapsed_election_ticks = 0;

        // If our log already covers the snapshot, ignore.
        if last_included_index <= self.log.snapshot_index() {
            acts.push(Action::SendMessage {
                to: from,
                msg: Message::InstallSnapshotResponse { from: self.cfg.id, term: self.hs.current_term, request_id },
            });
            return;
        }

        self.log.compact_through(last_included_index, last_included_term);
        self.last_applied = last_included_index;
        if self.hs.commit_index < last_included_index {
            self.hs.commit_index = last_included_index;
            acts.push(Action::PersistHardState(self.hs.clone()));
        }
        acts.push(Action::InstallSnapshot { last_included_index, last_included_term, data });
        acts.push(Action::SendMessage {
            to: from,
            msg: Message::InstallSnapshotResponse { from: self.cfg.id, term: self.hs.current_term, request_id },
        });
    }

    /// Inform the core that a local snapshot has been taken at `index/term`.
    /// The core will compact its in-memory log accordingly.
    pub fn on_snapshot_taken(&mut self, last_included_index: LogIndex, last_included_term: Term) {
        self.log.compact_through(last_included_index, last_included_term);
    }

    // ===================================================================
    // Configuration handling
    // ===================================================================

    fn apply_config_entry(&mut self, entry: &Entry) {
        match entry.kind {
            EntryKind::ConfigJoint => {
                if let Ok(state) = bincode::deserialize::<ConfigState>(&entry.data) {
                    self.config.state = state;
                    self.in_joint = true;
                    let voters = self.voters();
                    if matches!(self.role, Role::Leader) {
                        self.progress.ensure(&voters, self.log.last_index());
                        self.progress.retain_in(&voters);
                    }
                }
            }
            EntryKind::ConfigNew => {
                if let Ok(state) = bincode::deserialize::<ConfigState>(&entry.data) {
                    self.config.state = state;
                    self.in_joint = false;
                    let voters = self.voters();
                    if matches!(self.role, Role::Leader) {
                        self.progress.ensure(&voters, self.log.last_index());
                        self.progress.retain_in(&voters);
                    }
                }
            }
            _ => {}
        }
    }

    fn maybe_finish_joint(&mut self, acts: &mut SmallVec<[Action; 8]>) {
        if !self.in_joint || !matches!(self.role, Role::Leader) { return; }
        // Once `C_old,new` is committed, append `C_new`.
        let new_state = match &self.config.state {
            ConfigState::Joint { new, .. } => ConfigState::Stable(new.clone()),
            _ => return,
        };
        let payload = bincode::serialize(&new_state).expect("config encode");
        let idx = self.log.last_index() + 1;
        let entry = Entry { term: self.hs.current_term, index: idx, kind: EntryKind::ConfigNew, data: payload };
        self.log.append(vec![entry.clone()]);
        self.apply_config_entry(&entry);
        if let Some(p) = self.progress.get_mut(self.cfg.id) {
            p.match_index = self.log.last_index();
            p.next_index = self.log.last_index() + 1;
        }
        acts.push(Action::AppendEntries(vec![entry]));
        self.broadcast_replication(acts);
    }
}

#[cfg(test)]
mod tests {
    use super::*;

    fn cfg(id: NodeId, voters: Vec<NodeId>) -> RaftConfig {
        RaftConfig {
            id, initial_voters: voters,
            election_timeout_ticks: 10,
            heartbeat_ticks: 1,
            max_append_entries: 8,
            max_inflight_msgs: 16,
            rng_seed: 42, pre_vote: true,
        }
    }

    #[test]
    fn single_node_elects_self() {
        let mut n = RaftNode::new(cfg(1, vec![1]), HardState::default(), RaftLog::new());
        // tick until election fires
        for _ in 0..50 { let _ = n.tick(); }
        assert!(matches!(n.role(), Role::Leader));
    }

    #[test]
    fn three_node_election_with_majority() {
        let mut a = RaftNode::new(cfg(1, vec![1, 2, 3]), HardState::default(), RaftLog::new());
        // Force election.
        for _ in 0..50 { let _ = a.tick(); }
        // a should be Candidate (or PreCandidate) waiting for votes.
        assert!(matches!(a.role(), Role::Candidate | Role::PreCandidate));
        // Simulate granted vote from peer 2.
        let term = a.hard_state().current_term;
        let _ = a.step(Message::RequestVoteResponse {
            from: 2, term: if matches!(a.role(), Role::PreCandidate) { term + 1 } else { term },
            vote_granted: true, pre_vote: matches!(a.role(), Role::PreCandidate), request_id: 0,
        });
        // If pre-vote, that promotes to candidate; need another grant.
        if matches!(a.role(), Role::Candidate) {
            let term = a.hard_state().current_term;
            let _ = a.step(Message::RequestVoteResponse {
                from: 2, term, vote_granted: true, pre_vote: false, request_id: 0,
            });
        }
        assert!(matches!(a.role(), Role::Leader), "expected leader after majority votes");
    }

    #[test]
    fn append_entries_truncates_conflicting_suffix() {
        let mut n = RaftNode::new(cfg(2, vec![1, 2, 3]), HardState::default(), RaftLog::new());
        // Pretend leader 1 sends entries with prev=0.
        let entries = vec![
            Entry::normal(1, 1, b"a".to_vec()),
            Entry::normal(1, 2, b"b".to_vec()),
        ];
        let _ = n.step(Message::AppendEntries {
            from: 1, term: 1, prev_log_index: 0, prev_log_term: 0,
            entries, leader_commit: 0, request_id: 1,
        });
        assert_eq!(n.last_log_index(), 2);
        // New leader (term 2) overwrites index 2.
        let _ = n.step(Message::AppendEntries {
            from: 1, term: 2, prev_log_index: 1, prev_log_term: 1,
            entries: vec![Entry::normal(2, 2, b"B".to_vec())],
            leader_commit: 0, request_id: 2,
        });
        assert_eq!(n.last_log_index(), 2);
    }
}
\end{lstlisting}

\section{crates/raft-core/src/progress.rs}
\begin{lstlisting}[language=Rust]
//! Per-peer replication progress tracker (leader-side).
//!
//! Mirrors the etcd/raft `Progress` design. Each peer is in exactly one of
//! three states:
//!
//! * `Probe`     — slow, one-at-a-time AppendEntries until a match is found.
//! * `Replicate` — fast pipelined AppendEntries.
//! * `Snapshot`  — awaiting completion of an InstallSnapshot.

use crate::log::LogIndex;
use crate::node::NodeId;
use std::collections::BTreeMap;

/// Replication state for a single follower.
#[derive(Debug, Clone, PartialEq, Eq)]
pub enum ProgressState {
    /// Probing for the matching log point with one outstanding AE.
    Probe,
    /// Pipelining AEs.
    Replicate,
    /// Waiting for an InstallSnapshot to complete.
    Snapshot,
}

/// Per-peer progress. The leader maintains one of these per voter / learner.
#[derive(Debug, Clone)]
pub struct Progress {
    /// Index of the next log entry to send to this follower.
    pub next_index: LogIndex,
    /// Highest log entry known to be replicated on this follower.
    pub match_index: LogIndex,
    /// Current state.
    pub state: ProgressState,
    /// Whether the follower has responded since we became leader (used for
    /// read-index quorum confirmation).
    pub recent_active: bool,
    /// Number of in-flight AEs (very small bounded value used for backpressure).
    pub inflight: u32,
    /// Pending snapshot index, if any.
    pub pending_snapshot: Option<LogIndex>,
    /// Whether this peer is a non-voting learner.
    pub is_learner: bool,
}

impl Progress {
    /// Create a fresh progress with `next_index = last_log_index + 1`.
    pub fn new(last_log_index: LogIndex) -> Self {
        Self {
            next_index: last_log_index + 1,
            match_index: 0,
            state: ProgressState::Probe,
            recent_active: false,
            inflight: 0,
            pending_snapshot: None,
            is_learner: false,
        }
    }

    /// Note: use this when (re)becoming leader to start cautiously.
    pub fn reset_to_probe(&mut self, last_log_index: LogIndex) {
        self.next_index = last_log_index + 1;
        self.match_index = 0;
        self.state = ProgressState::Probe;
        self.inflight = 0;
        self.pending_snapshot = None;
    }

    /// Apply a successful AppendEntries response.
    pub fn maybe_update(&mut self, last_log_index_on_peer: LogIndex) -> bool {
        let updated = if last_log_index_on_peer > self.match_index {
            self.match_index = last_log_index_on_peer;
            true
        } else { false };
        if last_log_index_on_peer + 1 > self.next_index {
            self.next_index = last_log_index_on_peer + 1;
        }
        if matches!(self.state, ProgressState::Probe) {
            self.state = ProgressState::Replicate;
        }
        updated
    }

    /// Apply a rejected AppendEntries response with conflict hints.
    /// Returns true if `next_index` was decreased.
    pub fn maybe_decr_to(&mut self, conflict_index: LogIndex) -> bool {
        let target = conflict_index.max(1);
        if target < self.next_index {
            self.next_index = target;
            self.state = ProgressState::Probe;
            self.inflight = 0;
            true
        } else {
            false
        }
    }

    /// Mark a snapshot as in flight.
    pub fn become_snapshot(&mut self, snap_index: LogIndex) {
        self.state = ProgressState::Snapshot;
        self.pending_snapshot = Some(snap_index);
        self.inflight = 0;
    }

    /// Snapshot finished installing — resume probing from snap_index + 1.
    pub fn snapshot_finished(&mut self, snap_index: LogIndex) {
        self.match_index = self.match_index.max(snap_index);
        self.next_index = snap_index + 1;
        self.state = ProgressState::Probe;
        self.pending_snapshot = None;
    }

    /// Whether we may send an AppendEntries to this peer right now.
    pub fn can_send(&self, max_inflight: u32) -> bool {
        match self.state {
            ProgressState::Snapshot => false,
            ProgressState::Probe => self.inflight == 0,
            ProgressState::Replicate => self.inflight < max_inflight,
        }
    }
}

/// Map of peer-id -> Progress.
#[derive(Debug, Default)]
pub struct ProgressSet {
    /// Underlying storage.
    pub peers: BTreeMap<NodeId, Progress>,
}

impl ProgressSet {
    /// Insert or refresh entries for a list of voters.
    pub fn ensure(&mut self, ids: &[NodeId], last_log_index: LogIndex) {
        for &id in ids {
            self.peers.entry(id).or_insert_with(|| Progress::new(last_log_index));
        }
    }

    /// Borrow a peer's progress mutably.
    pub fn get_mut(&mut self, id: NodeId) -> Option<&mut Progress> { self.peers.get_mut(&id) }

    /// Borrow a peer's progress.
    pub fn get(&self, id: NodeId) -> Option<&Progress> { self.peers.get(&id) }

    /// Iterate all peers.
    pub fn iter(&self) -> impl Iterator<Item = (&NodeId, &Progress)> { self.peers.iter() }

    /// Drop peers no longer in the cluster.
    pub fn retain_in(&mut self, ids: &[NodeId]) {
        let keep: std::collections::BTreeSet<_> = ids.iter().copied().collect();
        self.peers.retain(|k, _| keep.contains(k));
    }
}
\end{lstlisting}

\section{crates/raft-core/src/role.rs}
\begin{lstlisting}[language=Rust]
//! Roles a Raft server may occupy.

use serde::{Deserialize, Serialize};

/// Raft roles. `PreCandidate` is the pre-vote extension (Ongaro thesis §9.6).
#[derive(Debug, Clone, Copy, PartialEq, Eq, Serialize, Deserialize)]
pub enum Role {
    /// Default role; redirects clients to the leader.
    Follower,
    /// Probing peers to determine election viability without bumping the term.
    PreCandidate,
    /// Has bumped term and is collecting votes.
    Candidate,
    /// Replicating entries.
    Leader,
}

impl Role {
    /// Short string for logs/metrics.
    pub fn as_str(self) -> &'static str {
        match self {
            Self::Follower => "follower",
            Self::PreCandidate => "precandidate",
            Self::Candidate => "candidate",
            Self::Leader => "leader",
        }
    }
}
\end{lstlisting}

\section{crates/raft-core/src/state.rs}
\begin{lstlisting}[language=Rust]
//! Persistent and volatile state of a Raft node.

use serde::{Deserialize, Serialize};

use crate::log::{LogIndex, Term};
use crate::node::NodeId;
use crate::role::Role;

/// State that **must** be persisted before responding to election/replication
/// RPCs (Raft §5.2, §5.3).
#[derive(Debug, Clone, Default, PartialEq, Eq, Serialize, Deserialize)]
pub struct HardState {
    /// Latest term server has seen (initialized to 0, increases monotonically).
    pub current_term: Term,
    /// Candidate id that received vote in the current term, or `None`.
    pub voted_for: Option<NodeId>,
    /// Highest log index known to be committed. Persisting this is optional for
    /// safety (Raft) but useful for fast restart.
    pub commit_index: LogIndex,
}

/// Volatile state observable by the runtime / metrics layer.
#[derive(Debug, Clone, PartialEq, Eq)]
pub struct SoftState {
    /// Current role.
    pub role: Role,
    /// Best-known leader for the current term.
    pub leader_id: Option<NodeId>,
}

impl Default for SoftState {
    fn default() -> Self {
        Self { role: Role::Follower, leader_id: None }
    }
}
\end{lstlisting}

\section{crates/raft-net/build.rs}
\begin{lstlisting}[language=Rust]
fn main() -> Result<(), Box<dyn std::error::Error>> {
    if std::env::var_os("PROTOC").is_none() {
        if let Ok(p) = protoc_bin_vendored::protoc_bin_path() {
            std::env::set_var("PROTOC", p);
        }
    }
    let proto_dir = "../../proto";
    tonic_build::configure()
        .build_server(true)
        .build_client(true)
        .compile_protos(
            &[
                format!("{}/raft.proto", proto_dir),
                format!("{}/membership.proto", proto_dir),
                format!("{}/admin.proto", proto_dir),
            ],
            &[proto_dir.to_string()],
        )?;
    Ok(())
}
\end{lstlisting}

\section{crates/raft-net/src/client.rs}
\begin{lstlisting}[language=Rust]
//! Outbound peer client.
//!
//! Each Raft RPC has exactly one synchronous reply. The peer client sends the
//! outbound request and feeds the response back into the local inbox so the
//! core can step on it. `*Response` variants in the message stream are no-ops
//! at this layer (they're already on the wire by the time we see them).

use std::sync::Arc;
use std::time::Duration;

use tokio::sync::Mutex;
use tonic::transport::{Channel, Endpoint};

use crate::convert::{
    append_request_to_pb, append_response_from_pb, snapshot_to_pb,
    vote_request_to_pb, vote_response_from_pb,
};
use crate::pb::raft as pb;
use crate::pb::raft::raft_client::RaftClient;
use crate::InboxTx;
use raft_core::{Message, NodeId};

/// One-per-peer client that handles connection lifecycle and message dispatch.
#[derive(Clone)]
pub struct PeerClient {
    /// Peer id (informational).
    pub peer_id: NodeId,
    /// gRPC endpoint URL, e.g. `http://node2:7000`.
    pub endpoint: String,
    /// Inbox to deliver synchronous Raft responses into.
    pub inbox: InboxTx,
    inner: Arc<Mutex<Option<RaftClient<Channel>>>>,
}

impl PeerClient {
    /// Create a client without connecting yet.
    pub fn new(peer_id: NodeId, endpoint: String, inbox: InboxTx) -> Self {
        Self { peer_id, endpoint, inbox, inner: Arc::new(Mutex::new(None)) }
    }

    async fn connect(&self) -> Result<RaftClient<Channel>, tonic::transport::Error> {
        let ep: Endpoint = Endpoint::from_shared(self.endpoint.clone())?
            .connect_timeout(Duration::from_millis(500))
            .timeout(Duration::from_secs(2));
        let chan = ep.connect().await?;
        Ok(RaftClient::new(chan))
    }

    async fn ensure(&self) -> Option<RaftClient<Channel>> {
        let mut g = self.inner.lock().await;
        if g.is_none() {
            match self.connect().await {
                Ok(c) => *g = Some(c),
                Err(e) => {
                    tracing::debug!(peer = self.peer_id, error = %e, "peer connect failed");
                    return None;
                }
            }
        }
        g.clone()
    }

    /// Send a single Raft message to the peer; deliver any synchronous
    /// response into the local inbox.
    pub async fn send(&self, msg: Message) {
        let Some(mut client) = self.ensure().await else { return; };
        let dropped = match &msg {
            Message::AppendEntries { .. } => {
                let Some(req) = append_request_to_pb(&msg) else { return; };
                match client.append_entries(req).await {
                    Ok(resp) => {
                        let resp_msg = append_response_from_pb(resp.into_inner(), self.peer_id);
                        let _ = self.inbox.send(resp_msg);
                        false
                    }
                    Err(e) => { tracing::debug!(peer = self.peer_id, error = %e, "AE rpc failed"); true }
                }
            }
            Message::RequestVote { pre_vote, .. } => {
                let Some(req) = vote_request_to_pb(&msg) else { return; };
                let pre = *pre_vote;
                let r = if pre { client.pre_vote(req).await } else { client.request_vote(req).await };
                match r {
                    Ok(resp) => {
                        let resp_msg = vote_response_from_pb(resp.into_inner(), self.peer_id, pre);
                        let _ = self.inbox.send(resp_msg);
                        false
                    }
                    Err(e) => { tracing::debug!(peer = self.peer_id, error = %e, "vote rpc failed"); true }
                }
            }
            Message::AppendEntriesResponse { .. } | Message::RequestVoteResponse { .. } | Message::InstallSnapshotResponse { .. } => {
                // These responses already flowed back over the inbound RPC's
                // synchronous reply path. Nothing to do here.
                false
            }
            Message::InstallSnapshot { .. } => {
                let Some(chunk) = snapshot_to_pb(&msg) else { return; };
                let stream = futures::stream::iter(vec![chunk]);
                match client.install_snapshot(stream).await {
                    Ok(_) => false,
                    Err(e) => { tracing::debug!(peer = self.peer_id, error = %e, "snapshot rpc failed"); true }
                }
            }
            Message::TimeoutNow { from, term } => {
                let req = pb::TimeoutNowRequest { term: *term, leader_id: *from };
                match client.timeout_now(req).await {
                    Ok(_) => false,
                    Err(e) => { tracing::debug!(peer = self.peer_id, error = %e, "timeout_now rpc failed"); true }
                }
            }
        };
        if dropped {
            *self.inner.lock().await = None;
        }
    }
}
\end{lstlisting}

\section{crates/raft-net/src/convert.rs}
\begin{lstlisting}[language=Rust]
//! Conversions between protobuf wire types and `raft_core::Message`.

use bytes::Bytes;
use raft_core::log::{Entry, EntryKind};
use raft_core::Message;

use crate::pb::raft as pb;

fn entry_kind_to_pb(k: EntryKind) -> i32 {
    match k {
        EntryKind::Noop => pb::EntryType::Noop as i32,
        EntryKind::Normal => pb::EntryType::Normal as i32,
        EntryKind::ConfigJoint => pb::EntryType::ConfigJoint as i32,
        EntryKind::ConfigNew => pb::EntryType::ConfigNew as i32,
    }
}

fn entry_kind_from_pb(t: i32) -> EntryKind {
    match pb::EntryType::try_from(t).unwrap_or(pb::EntryType::Normal) {
        pb::EntryType::Noop => EntryKind::Noop,
        pb::EntryType::Normal => EntryKind::Normal,
        pb::EntryType::ConfigJoint => EntryKind::ConfigJoint,
        pb::EntryType::ConfigNew => EntryKind::ConfigNew,
    }
}

/// Convert a core entry to wire form.
pub fn entry_to_pb(e: &Entry) -> pb::LogEntry {
    pb::LogEntry { term: e.term, index: e.index, entry_type: entry_kind_to_pb(e.kind), data: e.data.clone() }
}

/// Convert a wire entry to core form.
pub fn entry_from_pb(e: pb::LogEntry) -> Entry {
    Entry { term: e.term, index: e.index, kind: entry_kind_from_pb(e.entry_type), data: e.data }
}

fn rid_to_bytes(rid: u64) -> Vec<u8> { rid.to_le_bytes().to_vec() }
fn rid_from_bytes(b: &[u8]) -> u64 {
    let mut buf = [0u8; 8];
    let n = b.len().min(8);
    buf[..n].copy_from_slice(&b[..n]);
    u64::from_le_bytes(buf)
}

/// Convert AppendEntries request from wire.
pub fn append_request_from_pb(r: pb::AppendEntriesRequest) -> Message {
    Message::AppendEntries {
        from: r.leader_id,
        term: r.term,
        prev_log_index: r.prev_log_index,
        prev_log_term: r.prev_log_term,
        entries: r.entries.into_iter().map(entry_from_pb).collect(),
        leader_commit: r.leader_commit,
        request_id: rid_from_bytes(&r.request_id),
    }
}

/// Convert AppendEntries response from wire.
pub fn append_response_from_pb(r: pb::AppendEntriesResponse, from: u64) -> Message {
    Message::AppendEntriesResponse {
        from,
        term: r.term,
        success: r.success,
        conflict_index: r.conflict_index,
        conflict_term: r.conflict_term,
        last_log_index: r.last_log_index,
        request_id: rid_from_bytes(&r.request_id),
    }
}

/// Convert RequestVote/PreVote request from wire.
pub fn vote_request_from_pb(r: pb::RequestVoteRequest) -> Message {
    Message::RequestVote {
        from: r.candidate_id,
        term: r.term,
        last_log_index: r.last_log_index,
        last_log_term: r.last_log_term,
        pre_vote: r.pre_vote,
        request_id: rid_from_bytes(&r.request_id),
    }
}

/// Convert vote response from wire.
pub fn vote_response_from_pb(r: pb::RequestVoteResponse, from: u64, pre_vote: bool) -> Message {
    Message::RequestVoteResponse {
        from,
        term: r.term,
        vote_granted: r.vote_granted,
        pre_vote,
        request_id: rid_from_bytes(&r.request_id),
    }
}

/// Build a wire AppendEntries request.
pub fn append_request_to_pb(msg: &Message) -> Option<pb::AppendEntriesRequest> {
    if let Message::AppendEntries { from, term, prev_log_index, prev_log_term, entries, leader_commit, request_id } = msg {
        Some(pb::AppendEntriesRequest {
            term: *term,
            leader_id: *from,
            prev_log_index: *prev_log_index,
            prev_log_term: *prev_log_term,
            entries: entries.iter().map(entry_to_pb).collect(),
            leader_commit: *leader_commit,
            request_id: rid_to_bytes(*request_id),
        })
    } else { None }
}

/// Build a wire AppendEntries response.
pub fn append_response_to_pb(msg: &Message) -> Option<pb::AppendEntriesResponse> {
    if let Message::AppendEntriesResponse { term, success, conflict_index, conflict_term, last_log_index, request_id, .. } = msg {
        Some(pb::AppendEntriesResponse {
            term: *term,
            success: *success,
            conflict_index: *conflict_index,
            conflict_term: *conflict_term,
            last_log_index: *last_log_index,
            request_id: rid_to_bytes(*request_id),
        })
    } else { None }
}

/// Build a wire RequestVote request.
pub fn vote_request_to_pb(msg: &Message) -> Option<pb::RequestVoteRequest> {
    if let Message::RequestVote { from, term, last_log_index, last_log_term, pre_vote, request_id } = msg {
        Some(pb::RequestVoteRequest {
            term: *term,
            candidate_id: *from,
            last_log_index: *last_log_index,
            last_log_term: *last_log_term,
            pre_vote: *pre_vote,
            request_id: rid_to_bytes(*request_id),
        })
    } else { None }
}

/// Build a wire RequestVote response.
pub fn vote_response_to_pb(msg: &Message) -> Option<pb::RequestVoteResponse> {
    if let Message::RequestVoteResponse { term, vote_granted, request_id, .. } = msg {
        Some(pb::RequestVoteResponse {
            term: *term,
            vote_granted: *vote_granted,
            request_id: rid_to_bytes(*request_id),
        })
    } else { None }
}

/// Build a wire InstallSnapshot single-chunk message.
pub fn snapshot_to_pb(msg: &Message) -> Option<pb::InstallSnapshotChunk> {
    if let Message::InstallSnapshot { from, term, last_included_index, last_included_term, data, request_id } = msg {
        Some(pb::InstallSnapshotChunk {
            term: *term,
            leader_id: *from,
            last_included_index: *last_included_index,
            last_included_term: *last_included_term,
            offset: 0,
            data: data.to_vec(),
            done: true,
            config: vec![],
            request_id: rid_to_bytes(*request_id),
        })
    } else { None }
}

/// Reconstruct a single InstallSnapshot message from a chunk stream.
pub fn snapshot_from_chunks(chunks: Vec<pb::InstallSnapshotChunk>) -> Option<Message> {
    let head = chunks.first()?;
    let from = head.leader_id;
    let term = head.term;
    let last_included_index = head.last_included_index;
    let last_included_term = head.last_included_term;
    let request_id = rid_from_bytes(&head.request_id);
    let mut buf = Vec::new();
    for c in chunks { buf.extend_from_slice(&c.data); }
    Some(Message::InstallSnapshot {
        from,
        term,
        last_included_index,
        last_included_term,
        data: Bytes::from(buf),
        request_id,
    })
}
\end{lstlisting}

\section{crates/raft-net/src/lib.rs}
\begin{lstlisting}[language=Rust]
#![deny(unsafe_code)]
//! gRPC transport for Raft RPCs.
//!
//! Wire types live in the generated module `pb` below. Conversion helpers
//! map between the generated types and `raft_core::Message`.

#[allow(clippy::all)]
#[allow(missing_docs)]
pub mod pb {
    pub mod raft {
        tonic::include_proto!("raftkv.raft.v1");
    }
    pub mod membership {
        tonic::include_proto!("raftkv.membership.v1");
    }
    pub mod admin {
        tonic::include_proto!("raftkv.admin.v1");
    }
}

pub mod convert;
pub mod client;
pub mod server;

pub use client::PeerClient;
pub use server::{MessageProcessor, RaftServer};

/// Outbox for Raft messages produced by the core (consumed by the network layer).
pub type Outbox = tokio::sync::mpsc::UnboundedSender<(raft_core::NodeId, raft_core::Message)>;
/// Inbox of incoming Raft messages from peers.
pub type Inbox = tokio::sync::mpsc::UnboundedReceiver<raft_core::Message>;
/// Inbox sender side.
pub type InboxTx = tokio::sync::mpsc::UnboundedSender<raft_core::Message>;
\end{lstlisting}

\section{crates/raft-net/src/server.rs}
\begin{lstlisting}[language=Rust]
//! gRPC service implementation.
//!
//! The server is structured around a `MessageProcessor` trait: gRPC handlers
//! convert wire messages to `raft_core::Message`, hand them to the processor,
//! and return the response inline. The processor implementation lives in the
//! runtime — it locks the core, runs `step()`, and pulls the
//! `Action::SendMessage { to: <caller> }` whose body is the response.
//!
//! This is the *correct* mapping from Raft RPC semantics to gRPC: each Raft
//! RPC has exactly one synchronous reply, so it goes back over the same gRPC
//! call rather than through the outbound peer client.

use std::sync::Arc;

use async_trait::async_trait;
use tonic::{Request, Response, Status};

use crate::convert::{
    append_request_from_pb, snapshot_from_chunks, vote_request_from_pb,
};
use crate::pb::raft as pb;
use crate::pb::raft::raft_server::Raft as RaftService;
use raft_core::Message;

/// Synchronous message-processing hook the runtime implements.
#[async_trait]
pub trait MessageProcessor: Send + Sync + 'static {
    /// Step the core with `msg` and return the synchronous response message
    /// the core asked us to send back to `from`, if any.
    async fn process(&self, from: u64, msg: Message) -> Option<Message>;
}

/// gRPC service.
pub struct RaftServer {
    processor: Arc<dyn MessageProcessor>,
    /// Local node id.
    pub local_id: u64,
}

impl RaftServer {
    /// Construct a new server attached to a processor.
    pub fn new(local_id: u64, processor: Arc<dyn MessageProcessor>) -> Self {
        Self { processor, local_id }
    }
}

fn unwrap_append_response(msg: Option<Message>, term_fallback: u64, last_log_index: u64, request_id: Vec<u8>) -> pb::AppendEntriesResponse {
    if let Some(Message::AppendEntriesResponse { term, success, conflict_index, conflict_term, last_log_index, request_id: rid, .. }) = msg {
        pb::AppendEntriesResponse {
            term, success, conflict_index, conflict_term, last_log_index,
            request_id: rid.to_le_bytes().to_vec(),
        }
    } else {
        pb::AppendEntriesResponse {
            term: term_fallback,
            success: false,
            conflict_index: 0,
            conflict_term: 0,
            last_log_index,
            request_id,
        }
    }
}

fn unwrap_vote_response(msg: Option<Message>, term_fallback: u64, request_id: Vec<u8>) -> pb::RequestVoteResponse {
    if let Some(Message::RequestVoteResponse { term, vote_granted, request_id: rid, .. }) = msg {
        pb::RequestVoteResponse {
            term, vote_granted,
            request_id: rid.to_le_bytes().to_vec(),
        }
    } else {
        pb::RequestVoteResponse { term: term_fallback, vote_granted: false, request_id }
    }
}

#[tonic::async_trait]
impl RaftService for RaftServer {
    async fn append_entries(
        &self,
        request: Request<pb::AppendEntriesRequest>,
    ) -> Result<Response<pb::AppendEntriesResponse>, Status> {
        let req = request.into_inner();
        let request_id = req.request_id.clone();
        let term = req.term;
        let last_log_index_hint = req.prev_log_index + req.entries.len() as u64;
        let from = req.leader_id;
        let msg = append_request_from_pb(req);
        let resp_msg = self.processor.process(from, msg).await;
        let pb_resp = unwrap_append_response(resp_msg, term, last_log_index_hint, request_id);
        Ok(Response::new(pb_resp))
    }

    async fn request_vote(
        &self,
        request: Request<pb::RequestVoteRequest>,
    ) -> Result<Response<pb::RequestVoteResponse>, Status> {
        let req = request.into_inner();
        let request_id = req.request_id.clone();
        let term = req.term;
        let from = req.candidate_id;
        let msg = vote_request_from_pb(req);
        let resp_msg = self.processor.process(from, msg).await;
        Ok(Response::new(unwrap_vote_response(resp_msg, term, request_id)))
    }

    async fn pre_vote(
        &self,
        request: Request<pb::RequestVoteRequest>,
    ) -> Result<Response<pb::RequestVoteResponse>, Status> {
        self.request_vote(request).await
    }

    async fn install_snapshot(
        &self,
        request: Request<tonic::Streaming<pb::InstallSnapshotChunk>>,
    ) -> Result<Response<pb::InstallSnapshotResponse>, Status> {
        let mut stream = request.into_inner();
        let mut chunks = Vec::new();
        while let Some(c) = stream.message().await? {
            chunks.push(c);
        }
        let mut term = 0;
        let mut request_id = vec![];
        if let Some(msg) = snapshot_from_chunks(chunks) {
            if let Message::InstallSnapshot { from, term: t, request_id: rid, .. } = &msg {
                term = *t;
                request_id = rid.to_le_bytes().to_vec();
                let _ = self.processor.process(*from, msg).await;
            }
        }
        Ok(Response::new(pb::InstallSnapshotResponse { term, request_id }))
    }

    async fn timeout_now(
        &self,
        request: Request<pb::TimeoutNowRequest>,
    ) -> Result<Response<pb::TimeoutNowResponse>, Status> {
        let req = request.into_inner();
        let from = req.leader_id;
        let msg = Message::TimeoutNow { from, term: req.term };
        let _ = self.processor.process(from, msg).await;
        Ok(Response::new(pb::TimeoutNowResponse { accepted: true }))
    }
}
\end{lstlisting}

\section{crates/raft-storage/src/error.rs}
\begin{lstlisting}[language=Rust]
//! Storage error types.

use thiserror::Error;

/// All storage errors.
#[derive(Debug, Error)]
pub enum StorageError {
    /// I/O failed.
    #[error("io error: {0}")]
    Io(#[from] std::io::Error),
    /// Sled metadata error.
    #[error("sled error: {0}")]
    Sled(#[from] sled::Error),
    /// Encoding/decoding failed.
    #[error("codec error: {0}")]
    Codec(String),
    /// CRC mismatch detected on read.
    #[error("crc mismatch at offset {offset} in {file}")]
    Crc {
        /// File path.
        file: String,
        /// Byte offset.
        offset: u64,
    },
    /// Tried to read an index outside the log.
    #[error("log index {0} out of range")]
    OutOfRange(u64),
    /// Invariant violated.
    #[error("invariant: {0}")]
    Invariant(&'static str),
}
\end{lstlisting}

\section{crates/raft-storage/src/lib.rs}
\begin{lstlisting}[language=Rust]
#![deny(unsafe_code)]
//! Durable storage for the Raft consensus log + metadata.
//!
//! Contains:
//! * [`segmented_log::SegmentedLog`] — append-only log with CRC32C-checked
//!   records, segment rotation, configurable fsync, and torn-write recovery.
//! * [`meta::MetaStore`] — sled-backed `(currentTerm, votedFor, commitIndex)`.
//! * [`snapshot_store::SnapshotStore`] — atomic snapshot file management.

pub mod error;
pub mod meta;
pub mod segment;
pub mod segmented_log;
pub mod snapshot_store;

pub use error::StorageError;
pub use meta::MetaStore;
pub use segmented_log::SegmentedLog;
pub use snapshot_store::{SnapshotMeta, SnapshotStore};
\end{lstlisting}

\section{crates/raft-storage/src/meta.rs}
\begin{lstlisting}[language=Rust]
//! Sled-backed Raft persistent metadata: `(currentTerm, votedFor, commitIndex)`
//! plus snapshot bookmarks.

use std::path::Path;

use raft_core::state::HardState;
use raft_core::log::{LogIndex, Term};

use crate::error::StorageError;

const KEY_HARD_STATE: &[u8] = b"hs";
const KEY_SNAP_INDEX: &[u8] = b"snap_index";
const KEY_SNAP_TERM: &[u8] = b"snap_term";

/// Sled-backed metadata store.
pub struct MetaStore {
    db: sled::Db,
}

impl MetaStore {
    /// Open or create the metadata db at `path`.
    pub fn open<P: AsRef<Path>>(path: P) -> Result<Self, StorageError> {
        let db = sled::Config::new().path(path).flush_every_ms(Some(50)).open()?;
        Ok(Self { db })
    }

    /// Persist `HardState` and fsync. Must complete before responding to RPCs
    /// that depend on the state.
    pub fn save_hard_state(&self, hs: &HardState) -> Result<(), StorageError> {
        let bytes = bincode::serialize(hs).map_err(|e| StorageError::Codec(e.to_string()))?;
        self.db.insert(KEY_HARD_STATE, bytes)?;
        self.db.flush()?;
        Ok(())
    }

    /// Read `HardState`, or return default.
    pub fn load_hard_state(&self) -> Result<HardState, StorageError> {
        match self.db.get(KEY_HARD_STATE)? {
            Some(b) => bincode::deserialize(&b).map_err(|e| StorageError::Codec(e.to_string())),
            None => Ok(HardState::default()),
        }
    }

    /// Persist snapshot meta `(last_included_index, last_included_term)`.
    pub fn save_snapshot_pointer(&self, idx: LogIndex, term: Term) -> Result<(), StorageError> {
        self.db.insert(KEY_SNAP_INDEX, &idx.to_le_bytes())?;
        self.db.insert(KEY_SNAP_TERM, &term.to_le_bytes())?;
        self.db.flush()?;
        Ok(())
    }

    /// Read snapshot meta if present.
    pub fn load_snapshot_pointer(&self) -> Result<Option<(LogIndex, Term)>, StorageError> {
        let i = self.db.get(KEY_SNAP_INDEX)?;
        let t = self.db.get(KEY_SNAP_TERM)?;
        match (i, t) {
            (Some(i), Some(t)) => {
                let idx = u64::from_le_bytes(i[..8].try_into().unwrap_or([0; 8]));
                let term = u64::from_le_bytes(t[..8].try_into().unwrap_or([0; 8]));
                Ok(Some((idx, term)))
            }
            _ => Ok(None),
        }
    }
}
\end{lstlisting}

\section{crates/raft-storage/src/segment.rs}
\begin{lstlisting}[language=Rust]
//! Single segment file for the segmented Raft log.
//!
//! On-disk record format:
//!
//! ```text
//! [u32 little-endian: payload length]
//! [u32 little-endian: crc32c of payload]
//! [u8;  payload (bincode-encoded raft_core::Entry)]
//! ```
//!
//! A torn write at the tail (length read but not all payload bytes flushed)
//! is detected by:
//!   * length-prefix that runs past EOF, OR
//!   * crc32c mismatch
//! Either case truncates the segment at the start of the bad record.

use std::fs::{File, OpenOptions};
use std::io::{Read, Seek, SeekFrom, Write};
use std::path::{Path, PathBuf};

use raft_core::log::{Entry, LogIndex};

use crate::error::StorageError;

/// Append-only log segment.
pub struct Segment {
    pub(crate) path: PathBuf,
    pub(crate) base_index: LogIndex,
    pub(crate) file: File,
    pub(crate) write_offset: u64,
    pub(crate) index: Vec<(LogIndex, u64)>, // (index, file_offset)
    pub(crate) max_size_bytes: u64,
    pub(crate) sync_each: bool,
}

impl Segment {
    /// Create a new segment file (truncating any existing file).
    pub fn create(path: &Path, base_index: LogIndex, max_size_bytes: u64, sync_each: bool) -> Result<Self, StorageError> {
        let file = OpenOptions::new().read(true).write(true).create(true).truncate(true).open(path)?;
        Ok(Self {
            path: path.to_path_buf(),
            base_index,
            file,
            write_offset: 0,
            index: Vec::new(),
            max_size_bytes,
            sync_each,
        })
    }

    /// Open an existing segment, scanning it to recover the index and detect
    /// torn writes. Returns the loaded segment plus the highest valid index.
    pub fn open(path: &Path, base_index: LogIndex, max_size_bytes: u64, sync_each: bool) -> Result<(Self, Option<LogIndex>), StorageError> {
        let mut file = OpenOptions::new().read(true).write(true).open(path)?;
        let len = file.metadata()?.len();
        file.seek(SeekFrom::Start(0))?;

        let mut buf = Vec::new();
        file.read_to_end(&mut buf)?;
        let mut offset: usize = 0;
        let mut index: Vec<(LogIndex, u64)> = Vec::new();
        let mut last_good_off: u64 = 0;
        let mut highest: Option<LogIndex> = None;

        while offset + 8 <= buf.len() {
            let length = u32::from_le_bytes(buf[offset..offset + 4].try_into().unwrap()) as usize;
            let crc = u32::from_le_bytes(buf[offset + 4..offset + 8].try_into().unwrap());
            let payload_start = offset + 8;
            let payload_end = payload_start + length;
            if payload_end > buf.len() { break; }
            let payload = &buf[payload_start..payload_end];
            let actual = crc32c::crc32c(payload);
            if actual != crc { break; }
            let entry: Entry = match bincode::deserialize(payload) {
                Ok(e) => e,
                Err(_) => break,
            };
            index.push((entry.index, offset as u64));
            highest = Some(entry.index);
            offset = payload_end;
            last_good_off = offset as u64;
        }
        // Truncate any garbage tail.
        if last_good_off < len {
            file.set_len(last_good_off)?;
        }
        file.seek(SeekFrom::End(0))?;
        Ok((
            Self {
                path: path.to_path_buf(),
                base_index,
                file,
                write_offset: last_good_off,
                index,
                max_size_bytes,
                sync_each,
            },
            highest,
        ))
    }

    /// Append one entry.
    pub fn append(&mut self, entry: &Entry) -> Result<(), StorageError> {
        let payload = bincode::serialize(entry).map_err(|e| StorageError::Codec(e.to_string()))?;
        let crc = crc32c::crc32c(&payload);
        let length = u32::try_from(payload.len()).map_err(|_| StorageError::Invariant("entry too large"))?;
        let mut hdr = [0u8; 8];
        hdr[..4].copy_from_slice(&length.to_le_bytes());
        hdr[4..].copy_from_slice(&crc.to_le_bytes());

        let off = self.write_offset;
        self.file.write_all(&hdr)?;
        self.file.write_all(&payload)?;
        self.write_offset += 8 + payload.len() as u64;
        self.index.push((entry.index, off));
        if self.sync_each {
            self.file.sync_data()?;
        }
        Ok(())
    }

    /// fsync (caller-controlled, used by group commit).
    pub fn sync(&mut self) -> Result<(), StorageError> {
        self.file.sync_data()?;
        Ok(())
    }

    /// Whether this segment is at or above its max size.
    pub fn is_full(&self) -> bool { self.write_offset >= self.max_size_bytes }

    /// Last log index in this segment, if any.
    pub fn last_index(&self) -> Option<LogIndex> { self.index.last().map(|(i, _)| *i) }

    /// Read entry at `index`, if present in this segment.
    pub fn read(&mut self, index: LogIndex) -> Result<Option<Entry>, StorageError> {
        let pos = match self.index.iter().find(|(i, _)| *i == index) {
            Some((_, p)) => *p,
            None => return Ok(None),
        };
        self.file.seek(SeekFrom::Start(pos))?;
        let mut hdr = [0u8; 8];
        self.file.read_exact(&mut hdr)?;
        let length = u32::from_le_bytes(hdr[..4].try_into().unwrap()) as usize;
        let crc = u32::from_le_bytes(hdr[4..].try_into().unwrap());
        let mut buf = vec![0u8; length];
        self.file.read_exact(&mut buf)?;
        if crc32c::crc32c(&buf) != crc {
            return Err(StorageError::Crc { file: self.path.display().to_string(), offset: pos });
        }
        let entry: Entry = bincode::deserialize(&buf).map_err(|e| StorageError::Codec(e.to_string()))?;
        self.file.seek(SeekFrom::End(0))?;
        Ok(Some(entry))
    }

    /// Truncate so that no entries with index >= `from` remain.
    pub fn truncate_from(&mut self, from: LogIndex) -> Result<(), StorageError> {
        // Find first entry index >= from
        let split = self.index.iter().position(|(i, _)| *i >= from);
        if let Some(pos) = split {
            let off = self.index[pos].1;
            self.index.truncate(pos);
            self.file.set_len(off)?;
            self.file.seek(SeekFrom::End(0))?;
            self.write_offset = off;
        }
        Ok(())
    }
}
\end{lstlisting}

\section{crates/raft-storage/src/segmented\_log.rs}
\begin{lstlisting}[language=Rust]
//! Append-only segmented log with crash-safe recovery.
//!
//! Disk layout: `<dir>/00000000000000000001.log`, where the number is the
//! 1-based index of the first entry stored in that segment.

use std::fs;
use std::path::{Path, PathBuf};

use parking_lot::Mutex;
use raft_core::log::{Entry, LogIndex};

use crate::error::StorageError;
use crate::segment::Segment;

/// Configuration for a segmented log.
#[derive(Debug, Clone)]
pub struct SegmentedLogConfig {
    /// Directory containing segment files.
    pub dir: PathBuf,
    /// Roll segments at this size (default 64 MiB).
    pub max_segment_bytes: u64,
    /// fsync after every append (slow, safest) or only on `flush()` (fast).
    pub sync_each_append: bool,
}

impl Default for SegmentedLogConfig {
    fn default() -> Self {
        Self {
            dir: PathBuf::from("./data/raft-log"),
            max_segment_bytes: 64 * 1024 * 1024,
            sync_each_append: false,
        }
    }
}

/// Append-only segmented log.
pub struct SegmentedLog {
    cfg: SegmentedLogConfig,
    inner: Mutex<Inner>,
}

struct Inner {
    segments: Vec<Segment>,
    last_index: LogIndex,
    /// Snapshot tip below which entries are not retained.
    snapshot_index: LogIndex,
}

fn segment_filename(base_index: LogIndex) -> String { format!("{:020}.log", base_index) }

fn parse_segment_index(p: &Path) -> Option<LogIndex> {
    let stem = p.file_stem()?.to_str()?;
    stem.parse::<u64>().ok()
}

impl SegmentedLog {
    /// Open or create the log directory.
    pub fn open(cfg: SegmentedLogConfig) -> Result<Self, StorageError> {
        fs::create_dir_all(&cfg.dir)?;
        let mut entries: Vec<PathBuf> = fs::read_dir(&cfg.dir)?
            .filter_map(|e| e.ok().map(|e| e.path()))
            .filter(|p| p.extension().and_then(|s| s.to_str()) == Some("log"))
            .collect();
        entries.sort();

        let mut segments: Vec<Segment> = Vec::new();
        let mut last_index: LogIndex = 0;
        for path in entries {
            let base = parse_segment_index(&path).ok_or(StorageError::Invariant("bad segment filename"))?;
            let (seg, hi) = Segment::open(&path, base, cfg.max_segment_bytes, cfg.sync_each_append)?;
            if let Some(h) = hi { last_index = last_index.max(h); }
            segments.push(seg);
        }
        if segments.is_empty() {
            let path = cfg.dir.join(segment_filename(1));
            let seg = Segment::create(&path, 1, cfg.max_segment_bytes, cfg.sync_each_append)?;
            segments.push(seg);
        }
        Ok(Self {
            cfg,
            inner: Mutex::new(Inner { segments, last_index, snapshot_index: 0 }),
        })
    }

    /// Last index stored in the log.
    pub fn last_index(&self) -> LogIndex { self.inner.lock().last_index }

    /// First index that may still be served from the log (snapshot+1).
    pub fn first_index(&self) -> LogIndex { self.inner.lock().snapshot_index + 1 }

    /// Append a batch of entries (single fsync at end).
    pub fn append(&self, entries: &[Entry]) -> Result<(), StorageError> {
        if entries.is_empty() { return Ok(()); }
        let mut inner = self.inner.lock();
        let expected = inner.last_index + 1;
        if entries[0].index != expected {
            return Err(StorageError::Invariant("non-contiguous append"));
        }
        for e in entries {
            // Roll segment if full.
            let needs_roll = inner
                .segments
                .last()
                .map(Segment::is_full)
                .unwrap_or(true);
            if needs_roll {
                let path = self.cfg.dir.join(segment_filename(e.index));
                let seg = Segment::create(&path, e.index, self.cfg.max_segment_bytes, self.cfg.sync_each_append)?;
                inner.segments.push(seg);
            }
            let seg = inner.segments.last_mut().expect("segment");
            seg.append(e)?;
            inner.last_index = e.index;
        }
        if let Some(seg) = inner.segments.last_mut() { seg.sync()?; }
        Ok(())
    }

    /// Force an fsync.
    pub fn flush(&self) -> Result<(), StorageError> {
        let mut inner = self.inner.lock();
        if let Some(seg) = inner.segments.last_mut() { seg.sync()?; }
        Ok(())
    }

    /// Read a single entry by index.
    pub fn read(&self, index: LogIndex) -> Result<Option<Entry>, StorageError> {
        let mut inner = self.inner.lock();
        if index <= inner.snapshot_index || index > inner.last_index { return Ok(None); }
        // Find the right segment (largest base_index <= index).
        let pos = match inner.segments.iter().rposition(|s| s.base_index <= index) {
            Some(p) => p,
            None => return Ok(None),
        };
        inner.segments[pos].read(index)
    }

    /// Read entries with indices in `[from, to)`. Caller-supplied buffer to
    /// avoid allocation churn in hot paths.
    pub fn read_range(&self, from: LogIndex, to: LogIndex, out: &mut Vec<Entry>) -> Result<(), StorageError> {
        out.clear();
        if to <= from { return Ok(()); }
        for i in from..to {
            if let Some(e) = self.read(i)? { out.push(e); } else { break; }
        }
        Ok(())
    }

    /// Truncate so no entry with `index >= from` remains.
    pub fn truncate_from(&self, from: LogIndex) -> Result<(), StorageError> {
        let mut inner = self.inner.lock();
        if from > inner.last_index { return Ok(()); }
        // Drop whole segments whose base_index >= from
        while let Some(seg) = inner.segments.last() {
            if seg.base_index >= from {
                let p = seg.path.clone();
                inner.segments.pop();
                drop(fs::remove_file(p));
            } else { break; }
        }
        if let Some(seg) = inner.segments.last_mut() {
            seg.truncate_from(from)?;
            inner.last_index = seg.last_index().unwrap_or(seg.base_index.saturating_sub(1));
        } else {
            // recreate empty initial segment
            let path = self.cfg.dir.join(segment_filename(from.max(1)));
            let seg = Segment::create(&path, from.max(1), self.cfg.max_segment_bytes, self.cfg.sync_each_append)?;
            inner.segments.push(seg);
            inner.last_index = from.saturating_sub(1);
        }
        Ok(())
    }

    /// Drop entries up to `last_included_index` (inclusive). Used after a
    /// snapshot.
    pub fn compact(&self, last_included_index: LogIndex) -> Result<(), StorageError> {
        let mut inner = self.inner.lock();
        // Drop full segments whose final index <= last_included_index, but keep
        // at least one segment around.
        let mut to_drop: Vec<PathBuf> = Vec::new();
        while inner.segments.len() > 1 {
            let drop_first = match inner.segments.first().and_then(|s| s.last_index()) {
                Some(li) => li <= last_included_index,
                None => false,
            };
            if drop_first {
                let s = inner.segments.remove(0);
                to_drop.push(s.path.clone());
            } else { break; }
        }
        for p in to_drop { let _ = fs::remove_file(p); }
        inner.snapshot_index = last_included_index;
        Ok(())
    }

    /// Inform the log of the current snapshot tip without compacting (used on
    /// restart after a snapshot install).
    pub fn set_snapshot_tip(&self, idx: LogIndex) {
        let mut inner = self.inner.lock();
        inner.snapshot_index = idx;
        inner.last_index = inner.last_index.max(idx);
    }
}

#[cfg(test)]
mod tests {
    use super::*;
    use raft_core::log::Entry;
    use tempfile::tempdir;

    #[test]
    fn append_and_read_roundtrip() {
        let dir = tempdir().unwrap();
        let log = SegmentedLog::open(SegmentedLogConfig {
            dir: dir.path().to_path_buf(),
            max_segment_bytes: 1024,
            sync_each_append: false,
        }).unwrap();
        let entries = (1u64..=20).map(|i| Entry::normal(1, i, vec![i as u8; 50])).collect::<Vec<_>>();
        log.append(&entries).unwrap();
        assert_eq!(log.last_index(), 20);
        for i in 1u64..=20 {
            assert_eq!(log.read(i).unwrap().unwrap().index, i);
        }
    }

    #[test]
    fn truncate_drops_tail() {
        let dir = tempdir().unwrap();
        let log = SegmentedLog::open(SegmentedLogConfig { dir: dir.path().into(), max_segment_bytes: 64*1024, sync_each_append: false }).unwrap();
        let entries: Vec<_> = (1..=10).map(|i| Entry::normal(1, i, vec![1])).collect();
        log.append(&entries).unwrap();
        log.truncate_from(6).unwrap();
        assert_eq!(log.last_index(), 5);
        assert!(log.read(6).unwrap().is_none());
    }

    #[test]
    fn reopen_recovers_entries() {
        let dir = tempdir().unwrap();
        {
            let log = SegmentedLog::open(SegmentedLogConfig { dir: dir.path().into(), max_segment_bytes: 64*1024, sync_each_append: true }).unwrap();
            let entries: Vec<_> = (1..=5).map(|i| Entry::normal(1, i, vec![1])).collect();
            log.append(&entries).unwrap();
        }
        let log = SegmentedLog::open(SegmentedLogConfig { dir: dir.path().into(), max_segment_bytes: 64*1024, sync_each_append: false }).unwrap();
        assert_eq!(log.last_index(), 5);
    }
}
\end{lstlisting}

\section{crates/raft-storage/src/snapshot\_store.rs}
\begin{lstlisting}[language=Rust]
//! Snapshot storage: write-to-temp + atomic rename. Crash-safe.

use std::fs::{self, File, OpenOptions};
use std::io::{Read, Write};
use std::path::{Path, PathBuf};

use raft_core::log::{LogIndex, Term};
use serde::{Deserialize, Serialize};

use crate::error::StorageError;

/// Metadata for a snapshot file.
#[derive(Debug, Clone, Serialize, Deserialize)]
pub struct SnapshotMeta {
    /// Last log index covered.
    pub last_included_index: LogIndex,
    /// Term at `last_included_index`.
    pub last_included_term: Term,
    /// Optional cluster config bytes.
    pub config: Vec<u8>,
}

/// Snapshot file management with atomic rename.
pub struct SnapshotStore {
    dir: PathBuf,
}

impl SnapshotStore {
    /// Open or create the snapshot directory.
    pub fn open<P: AsRef<Path>>(path: P) -> Result<Self, StorageError> {
        fs::create_dir_all(&path)?;
        Ok(Self { dir: path.as_ref().to_path_buf() })
    }

    /// Path of the latest snapshot, if any.
    pub fn latest(&self) -> Option<(SnapshotMeta, PathBuf)> {
        let mut entries: Vec<_> = fs::read_dir(&self.dir).ok()?.filter_map(|e| e.ok()).collect();
        entries.sort_by_key(|e| e.file_name());
        let last = entries.into_iter().rev().find(|e| {
            e.path().extension().and_then(|s| s.to_str()) == Some("snap")
        })?;
        let meta_path = last.path().with_extension("meta");
        let mut buf = Vec::new();
        File::open(&meta_path).ok()?.read_to_end(&mut buf).ok()?;
        let meta: SnapshotMeta = bincode::deserialize(&buf).ok()?;
        Some((meta, last.path()))
    }

    /// Begin writing a snapshot to a temp file. Returns a handle.
    pub fn begin_write(&self, meta: &SnapshotMeta) -> Result<SnapshotWriter, StorageError> {
        let stem = format!("{:020}-{}", meta.last_included_index, meta.last_included_term);
        let snap_path = self.dir.join(format!("{stem}.snap"));
        let tmp_path = self.dir.join(format!("{stem}.snap.tmp"));
        let meta_path = self.dir.join(format!("{stem}.meta"));
        let file = OpenOptions::new().create(true).truncate(true).write(true).open(&tmp_path)?;
        Ok(SnapshotWriter { tmp_path, snap_path, meta_path, meta: meta.clone(), file })
    }

    /// List snapshots (oldest first).
    pub fn list(&self) -> Vec<PathBuf> {
        let mut out: Vec<_> = fs::read_dir(&self.dir)
            .into_iter()
            .flatten()
            .filter_map(|e| e.ok().map(|e| e.path()))
            .filter(|p| p.extension().and_then(|s| s.to_str()) == Some("snap"))
            .collect();
        out.sort();
        out
    }

    /// Remove all but the most recent N snapshots.
    pub fn keep_last(&self, n: usize) -> Result<(), StorageError> {
        let snaps = self.list();
        if snaps.len() <= n { return Ok(()); }
        let drop_n = snaps.len() - n;
        for p in &snaps[..drop_n] {
            let _ = fs::remove_file(p);
            let _ = fs::remove_file(p.with_extension("meta"));
        }
        Ok(())
    }
}

/// A snapshot in the process of being written.
pub struct SnapshotWriter {
    tmp_path: PathBuf,
    snap_path: PathBuf,
    meta_path: PathBuf,
    meta: SnapshotMeta,
    file: File,
}

impl SnapshotWriter {
    /// Write a chunk.
    pub fn write_chunk(&mut self, data: &[u8]) -> Result<(), StorageError> {
        self.file.write_all(data)?;
        Ok(())
    }

    /// fsync and atomically install the snapshot.
    pub fn finish(self) -> Result<(), StorageError> {
        self.file.sync_all()?;
        // Write meta first (also via temp+rename for atomicity).
        let meta_tmp = self.meta_path.with_extension("meta.tmp");
        let bytes = bincode::serialize(&self.meta).map_err(|e| StorageError::Codec(e.to_string()))?;
        std::fs::write(&meta_tmp, bytes)?;
        std::fs::rename(&meta_tmp, &self.meta_path)?;
        std::fs::rename(&self.tmp_path, &self.snap_path)?;
        Ok(())
    }
}
\end{lstlisting}

\section{crates/raftkv-cli/src/main.rs}
\begin{lstlisting}[language=Rust]
//! `raftkv-cli` — operator CLI for the cluster.
//!
//! Communicates with a node over its RESP port. The CLI focuses on
//! observability commands (INFO, CLUSTER NODES) plus the standard KV ops so
//! it can serve as a smoke-test client.

#![deny(unsafe_code)]

use std::time::Duration;

use anyhow::{Context, Result};
use clap::{Parser, Subcommand};
use tokio::io::{AsyncReadExt, AsyncWriteExt};
use tokio::net::TcpStream;

#[derive(Debug, Parser)]
#[command(name = "raftkv-cli", version, about = "raftkv operator CLI")]
struct Cli {
    /// Node RESP endpoint, e.g. 127.0.0.1:6379.
    #[arg(long, default_value = "127.0.0.1:6379", env = "RAFTKV_ADDR")]
    addr: String,
    #[command(subcommand)]
    cmd: Cmd,
}

#[derive(Debug, Subcommand)]
enum Cmd {
    /// Show cluster status.
    Status,
    /// List cluster members.
    Members,
    /// SET a key.
    Set {
        /// Key.
        key: String,
        /// Value.
        value: String,
    },
    /// GET a key.
    Get {
        /// Key.
        key: String,
    },
    /// DEL keys.
    Del {
        /// Keys to delete.
        keys: Vec<String>,
    },
    /// PING the server.
    Ping,
}

#[tokio::main]
async fn main() -> Result<()> {
    tracing_subscriber::fmt()
        .with_env_filter(tracing_subscriber::EnvFilter::try_from_default_env()
            .unwrap_or_else(|_| tracing_subscriber::EnvFilter::new("info")))
        .init();

    let cli = Cli::parse();
    let mut sock = TcpStream::connect(&cli.addr).await
        .with_context(|| format!("connecting to {}", cli.addr))?;
    sock.set_nodelay(true).ok();

    let bytes = match cli.cmd {
        Cmd::Status => resp_array(&["INFO", "raft"]),
        Cmd::Members => resp_array(&["CLUSTER", "NODES"]),
        Cmd::Set { key, value } => resp_array(&["SET", &key, &value]),
        Cmd::Get { key } => resp_array(&["GET", &key]),
        Cmd::Del { keys } => {
            let mut v = vec!["DEL".to_string()];
            v.extend(keys);
            let s: Vec<&str> = v.iter().map(|s| s.as_str()).collect();
            resp_array(&s)
        }
        Cmd::Ping => resp_array(&["PING"]),
    };
    sock.write_all(&bytes).await?;
    let mut buf = vec![0u8; 64 * 1024];
    let n = tokio::time::timeout(Duration::from_secs(5), sock.read(&mut buf)).await??;
    print!("{}", String::from_utf8_lossy(&buf[..n]));
    Ok(())
}

fn resp_array(parts: &[&str]) -> Vec<u8> {
    let mut out = Vec::new();
    out.extend_from_slice(format!("*{}\r\n", parts.len()).as_bytes());
    for p in parts {
        out.extend_from_slice(format!("${}\r\n", p.len()).as_bytes());
        out.extend_from_slice(p.as_bytes());
        out.extend_from_slice(b"\r\n");
    }
    out
}
\end{lstlisting}

\section{crates/raftkv-server/src/config.rs}
\begin{lstlisting}[language=Rust]
//! Server configuration.

use std::path::PathBuf;

use serde::{Deserialize, Serialize};

/// One peer in the cluster.
#[derive(Debug, Clone, Serialize, Deserialize)]
pub struct Peer {
    /// Cluster-unique node id.
    pub id: u64,
    /// gRPC address (e.g. http://node2:7000).
    pub raft_addr: String,
}

/// Top-level server config.
#[derive(Debug, Clone, Serialize, Deserialize)]
pub struct ServerConfig {
    /// This node's id.
    pub id: u64,
    /// gRPC bind address (Raft peer-to-peer).
    pub raft_listen: String,
    /// RESP bind address (clients).
    pub client_listen: String,
    /// Prometheus metrics scrape endpoint.
    pub metrics_listen: String,
    /// Initial voters (cluster bootstrap). Must be the same on all nodes at
    /// first boot.
    pub peers: Vec<Peer>,
    /// Data directory.
    pub data_dir: PathBuf,
    /// Election timeout in milliseconds (low end of randomized range).
    pub election_timeout_ms: u64,
    /// Heartbeat interval in milliseconds.
    pub heartbeat_ms: u64,
    /// Tick granularity in milliseconds.
    pub tick_ms: u64,
    /// Snapshot threshold: trigger when log exceeds N entries past last snap.
    pub snapshot_entries_threshold: u64,
    /// Enable pre-vote.
    pub pre_vote: bool,
}

impl Default for ServerConfig {
    fn default() -> Self {
        Self {
            id: 1,
            raft_listen: "0.0.0.0:7001".into(),
            client_listen: "0.0.0.0:6379".into(),
            metrics_listen: "0.0.0.0:9100".into(),
            peers: vec![Peer { id: 1, raft_addr: "http://127.0.0.1:7001".into() }],
            data_dir: PathBuf::from("./data"),
            election_timeout_ms: 300,
            heartbeat_ms: 50,
            tick_ms: 10,
            snapshot_entries_threshold: 10_000,
            pre_vote: true,
        }
    }
}

impl ServerConfig {
    /// Load config from a TOML file.
    pub fn from_toml_file(path: &std::path::Path) -> anyhow::Result<Self> {
        let s = std::fs::read_to_string(path)?;
        Ok(toml::from_str(&s)?)
    }
}
\end{lstlisting}

\section{crates/raftkv-server/src/main.rs}
\begin{lstlisting}[language=Rust]
//! `raftkv-server` binary entry point. Wires together:
//!  * Configuration loading (TOML + CLI override)
//!  * Tracing (env-filtered)
//!  * The runtime (raft core driver)
//!  * gRPC peer server (raft-net)
//!  * RESP client server (resp-server)
//!  * Prometheus metrics endpoint
//!
//! On SIGINT/SIGTERM we cancel the listeners and let the runtime drain.

#![deny(unsafe_code)]

use std::path::PathBuf;
use std::sync::Arc;

use clap::Parser;
use tracing_subscriber::EnvFilter;

mod config;
mod runtime;

use config::ServerConfig;
use runtime::{ClientHandler, Runtime};

#[derive(Debug, Parser)]
#[command(name = "raftkv-server", version)]
struct Cli {
    /// Path to TOML config.
    #[arg(long, env = "RAFTKV_CONFIG")]
    config: Option<PathBuf>,
    /// Override node id.
    #[arg(long, env = "RAFTKV_ID")]
    id: Option<u64>,
    /// Override raft listen address.
    #[arg(long, env = "RAFTKV_RAFT_LISTEN")]
    raft_listen: Option<String>,
    /// Override client (RESP) listen address.
    #[arg(long, env = "RAFTKV_CLIENT_LISTEN")]
    client_listen: Option<String>,
    /// Override metrics listen address.
    #[arg(long, env = "RAFTKV_METRICS_LISTEN")]
    metrics_listen: Option<String>,
    /// Override data dir.
    #[arg(long, env = "RAFTKV_DATA_DIR")]
    data_dir: Option<PathBuf>,
}

#[tokio::main]
async fn main() -> anyhow::Result<()> {
    init_tracing();
    let cli = Cli::parse();

    let mut cfg = match cli.config {
        Some(p) => ServerConfig::from_toml_file(&p)?,
        None => ServerConfig::default(),
    };
    if let Some(v) = cli.id { cfg.id = v; }
    if let Some(v) = cli.raft_listen { cfg.raft_listen = v; }
    if let Some(v) = cli.client_listen { cfg.client_listen = v; }
    if let Some(v) = cli.metrics_listen { cfg.metrics_listen = v; }
    if let Some(v) = cli.data_dir { cfg.data_dir = v; }

    tracing::info!(node = cfg.id, raft = %cfg.raft_listen, client = %cfg.client_listen, "starting raftkv-server");

    let raft_addr: std::net::SocketAddr = strip_scheme(&cfg.raft_listen).parse()?;
    let client_addr = cfg.client_listen.clone();
    let metrics_addr: std::net::SocketAddr = strip_scheme(&cfg.metrics_listen).parse()?;

    let (rt, handles) = Runtime::new(cfg).await?;
    let local_id = rt.local_id();

    rt.clone().spawn(handles);

    let processor: Arc<dyn raft_net::MessageProcessor> = rt.clone();
    let raft_server = raft_net::server::RaftServer::new(local_id, processor);
    let raft_grpc = tokio::spawn(async move {
        let svc = raft_net::pb::raft::raft_server::RaftServer::new(raft_server);
        if let Err(e) = tonic::transport::Server::builder().add_service(svc).serve(raft_addr).await {
            tracing::error!(error = %e, "gRPC server exited");
        }
    });

    let handler: Arc<ClientHandler> = Arc::new(ClientHandler::new(rt.clone()));
    let client_handler = handler.clone();
    let resp_task = tokio::spawn(async move {
        if let Err(e) = resp_server::server::serve(&client_addr, client_handler).await {
            tracing::error!(error = %e, "RESP server exited");
        }
    });

    let metrics_task = tokio::spawn(async move {
        if let Err(e) = serve_metrics(metrics_addr).await {
            tracing::warn!(error = %e, "metrics server exited");
        }
    });

    tokio::signal::ctrl_c().await.ok();
    tracing::info!("shutdown signal received");
    raft_grpc.abort();
    resp_task.abort();
    metrics_task.abort();
    Ok(())
}

fn init_tracing() {
    let filter = EnvFilter::try_from_default_env().unwrap_or_else(|_| EnvFilter::new("info,raftkv=info"));
    tracing_subscriber::fmt()
        .with_env_filter(filter)
        .with_target(true)
        .with_level(true)
        .init();
}

fn strip_scheme(s: &str) -> &str {
    s.strip_prefix("http://").or_else(|| s.strip_prefix("https://")).unwrap_or(s)
}

/// Minimal Prometheus text-format endpoint over plain TCP.
async fn serve_metrics(addr: std::net::SocketAddr) -> std::io::Result<()> {
    use tokio::io::AsyncWriteExt;
    use tokio::net::TcpListener;
    let listener = TcpListener::bind(addr).await?;
    tracing::info!(%addr, "metrics endpoint listening");
    loop {
        let (mut sock, _) = listener.accept().await?;
        tokio::spawn(async move {
            let body = match prometheus::TextEncoder::new()
                .encode_to_string(&prometheus::gather())
            {
                Ok(s) => s,
                Err(_) => String::new(),
            };
            let resp = format!(
                "HTTP/1.1 200 OK\r\nContent-Type: text/plain; version=0.0.4\r\nContent-Length: {}\r\n\r\n{}",
                body.len(), body
            );
            let _ = sock.write_all(resp.as_bytes()).await;
            let _ = sock.shutdown().await;
        });
    }
}
\end{lstlisting}

\section{crates/raftkv-server/src/runtime.rs}
\begin{lstlisting}[language=Rust]
//! The runtime drives `raft_core::RaftNode`. It owns:
//!
//! * The pure state machine (locked behind a Mutex).
//! * The durable log + meta + snapshot stores.
//! * The KV state machine.
//! * Outbound peer clients and an inbox for incoming peer responses.
//! * Channels for client proposals and read-index requests.
//!
//! Side-effect mapping: every `Action` produced by the core is interpreted
//! here.

use std::collections::HashMap;
use std::sync::Arc;
use std::time::{Duration, SystemTime, UNIX_EPOCH};

use async_trait::async_trait;
use kv_state_machine::{Command, KvStateMachine, Response};
use parking_lot::Mutex;
use raft_core::{
    Action, ConfigChange, Entry, EntryKind, Message, MetricEvent, NodeId, ProposeError,
    RaftConfig, RaftLog, RaftNode, Role,
};
use raft_net::{client::PeerClient, server::MessageProcessor, Inbox, InboxTx};
use raft_storage::{MetaStore, SegmentedLog, SnapshotMeta, SnapshotStore};
use resp_server::handler::{CommandHandler, ProposeFailure};
use resp_server::codec::RespFrame;
use tokio::sync::{mpsc, oneshot, Notify};

use crate::config::ServerConfig;

/// Per-proposal waiter info.
struct PendingProposal {
    expected_index: Option<u64>,
    tx: oneshot::Sender<Result<Response, ProposeFailure>>,
}

/// Read-index waiter (linearizable read).
struct PendingRead {
    commit_target: u64,
    tx: oneshot::Sender<RespFrame>,
    f: Box<dyn FnOnce(&KvStateMachine) -> RespFrame + Send>,
}

/// Proposal request enqueued from the client task.
enum ProposalKind {
    Command(Command),
    ConfigChange(ConfigChange),
}

struct ProposalRequest {
    kind: ProposalKind,
    tx: oneshot::Sender<Result<Response, ProposeFailure>>,
}

struct ReadRequest {
    f: Box<dyn FnOnce(&KvStateMachine) -> RespFrame + Send>,
    tx: oneshot::Sender<RespFrame>,
}

/// The runtime, shared with client tasks via Arc.
pub struct Runtime {
    cfg: ServerConfig,
    node: Mutex<RaftNode>,
    log: Arc<SegmentedLog>,
    meta: Arc<MetaStore>,
    snaps: Arc<SnapshotStore>,
    sm: Arc<KvStateMachine>,
    peers: HashMap<NodeId, PeerClient>,
    proposal_tx: mpsc::UnboundedSender<ProposalRequest>,
    read_tx: mpsc::UnboundedSender<ReadRequest>,
    notify: Arc<Notify>,
    soft: Mutex<SoftCache>,
    addr_book: HashMap<NodeId, String>,
    pending: Mutex<HashMap<u64, PendingProposal>>,
    pending_reads: Mutex<Vec<PendingRead>>,
    last_term: Mutex<u64>,
}

#[derive(Default)]
struct SoftCache {
    role: String,
    term: u64,
    leader: Option<NodeId>,
    commit_index: u64,
    applied_index: u64,
    last_log_index: u64,
}

impl Runtime {
    /// Create the runtime, opening durable stores and constructing the core.
    pub async fn new(cfg: ServerConfig) -> anyhow::Result<(Arc<Self>, RuntimeHandles)> {
        std::fs::create_dir_all(&cfg.data_dir)?;
        let log = Arc::new(SegmentedLog::open(raft_storage::segmented_log::SegmentedLogConfig {
            dir: cfg.data_dir.join("log"),
            max_segment_bytes: 64 * 1024 * 1024,
            sync_each_append: false,
        })?);
        let meta = Arc::new(MetaStore::open(cfg.data_dir.join("meta"))?);
        let snaps = Arc::new(SnapshotStore::open(cfg.data_dir.join("snapshots"))?);
        let sm = Arc::new(KvStateMachine::open(cfg.data_dir.join("kv"))?);

        let hs = meta.load_hard_state()?;
        let snap_ptr = meta.load_snapshot_pointer()?.unwrap_or((0, 0));
        let mut raft_log = RaftLog::from_snapshot(snap_ptr.0, snap_ptr.1);
        let mut buf = Vec::new();
        let last = log.last_index();
        let first = log.first_index().max(snap_ptr.0 + 1);
        if last >= first {
            log.read_range(first, last + 1, &mut buf)?;
            raft_log.append(buf);
        }

        let node = RaftNode::new(
            RaftConfig {
                id: cfg.id,
                initial_voters: cfg.peers.iter().map(|p| p.id).collect(),
                election_timeout_ticks: (cfg.election_timeout_ms / cfg.tick_ms.max(1)).max(1),
                heartbeat_ticks: (cfg.heartbeat_ms / cfg.tick_ms.max(1)).max(1),
                max_append_entries: 256,
                max_inflight_msgs: 1024,
                rng_seed: 0x600D_F00D ^ cfg.id,
                pre_vote: cfg.pre_vote,
            },
            hs,
            raft_log,
        );

        let (proposal_tx, proposal_rx) = mpsc::unbounded_channel();
        let (read_tx, read_rx) = mpsc::unbounded_channel();
        let (inbox_tx, inbox_rx): (InboxTx, Inbox) = mpsc::unbounded_channel();
        let notify = Arc::new(Notify::new());

        let mut peers: HashMap<NodeId, PeerClient> = HashMap::new();
        let mut addr_book: HashMap<NodeId, String> = HashMap::new();
        for p in &cfg.peers {
            addr_book.insert(p.id, p.raft_addr.clone());
            if p.id != cfg.id {
                peers.insert(p.id, PeerClient::new(p.id, p.raft_addr.clone(), inbox_tx.clone()));
            }
        }

        let rt = Arc::new(Self {
            cfg,
            node: Mutex::new(node),
            log,
            meta,
            snaps,
            sm,
            peers,
            proposal_tx,
            read_tx,
            notify,
            soft: Mutex::new(SoftCache::default()),
            addr_book,
            pending: Mutex::new(HashMap::new()),
            pending_reads: Mutex::new(Vec::new()),
            last_term: Mutex::new(0),
        });

        Ok((rt.clone(), RuntimeHandles { inbox_tx, inbox_rx, proposal_rx, read_rx }))
    }

    /// Spawn the runtime loop.
    pub fn spawn(self: Arc<Self>, handles: RuntimeHandles) {
        let RuntimeHandles { inbox_tx: _, mut inbox_rx, mut proposal_rx, mut read_rx } = handles;
        let me = self.clone();
        tokio::spawn(async move {
            let tick = Duration::from_millis(me.cfg.tick_ms);
            let mut tick_iv = tokio::time::interval(tick);
            tick_iv.set_missed_tick_behavior(tokio::time::MissedTickBehavior::Delay);

            loop {
                tokio::select! {
                    _ = tick_iv.tick() => {
                        let actions = me.node.lock().tick();
                        me.execute_actions(actions).await;
                    }
                    Some(msg) = inbox_rx.recv() => {
                        let actions = me.node.lock().step(msg);
                        me.execute_actions(actions).await;
                    }
                    Some(req) = proposal_rx.recv() => {
                        me.handle_proposal(req).await;
                    }
                    Some(req) = read_rx.recv() => {
                        me.handle_read(req).await;
                    }
                }
                me.fail_stale_proposals();
            }
        });
    }

    fn fail_stale_proposals(&self) {
        let (term, role) = {
            let n = self.node.lock();
            (n.hard_state().current_term, n.role())
        };
        let mut last = self.last_term.lock();
        if term != *last && !matches!(role, Role::Leader) {
            let mut pending = self.pending.lock();
            for (_, p) in pending.drain() {
                let _ = p.tx.send(Err(ProposeFailure::NotLeader { hint: None }));
            }
            *last = term;
        }
    }

    fn now_ms() -> u64 {
        SystemTime::now().duration_since(UNIX_EPOCH).map(|d| d.as_millis() as u64).unwrap_or(0)
    }

    async fn handle_proposal(self: &Arc<Self>, req: ProposalRequest) {
        let acts_res = match req.kind {
            ProposalKind::Command(cmd) => {
                let bytes = bincode::serialize(&cmd).unwrap_or_default();
                let mut node = self.node.lock();
                let prev_last = node.last_log_index();
                match node.propose(bytes) {
                    Ok(acts) => Ok((prev_last + 1, acts)),
                    Err(e) => Err(e),
                }
            }
            ProposalKind::ConfigChange(cc) => {
                let mut node = self.node.lock();
                let prev_last = node.last_log_index();
                match node.propose_config_change(cc) {
                    Ok(acts) => Ok((prev_last + 1, acts)),
                    Err(e) => Err(e),
                }
            }
        };
        match acts_res {
            Ok((idx, acts)) => {
                self.pending.lock().insert(idx, PendingProposal { expected_index: Some(idx), tx: req.tx });
                self.execute_actions(acts).await;
            }
            Err(e) => {
                let f = match e {
                    ProposeError::NotLeader { hint } => ProposeFailure::NotLeader {
                        hint: hint.and_then(|h| self.addr_book.get(&h).cloned()),
                    },
                    other => ProposeFailure::Other(other.to_string()),
                };
                let _ = req.tx.send(Err(f));
            }
        }
    }

    async fn handle_read(self: &Arc<Self>, req: ReadRequest) {
        let acts;
        let commit_target;
        let is_leader;
        {
            let mut node = self.node.lock();
            is_leader = matches!(node.role(), Role::Leader);
            commit_target = node.commit_index();
            acts = if is_leader { node.read_index(bytes::Bytes::new()) } else { Default::default() };
        }
        if !is_leader {
            let voters_only_self = self.cfg.peers.len() == 1;
            if voters_only_self {
                let resp = (req.f)(&self.sm);
                let _ = req.tx.send(resp);
                return;
            }
            let _ = req.tx.send(RespFrame::err("MOVED 0 ".to_string()));
            return;
        }
        self.pending_reads.lock().push(PendingRead { commit_target, tx: req.tx, f: req.f });
        self.execute_actions(acts).await;
        self.drain_ready_reads();
    }

    fn drain_ready_reads(&self) {
        let applied = self.sm.applied_index();
        let mut reads = self.pending_reads.lock();
        let mut still: Vec<PendingRead> = Vec::with_capacity(reads.len());
        for r in reads.drain(..) {
            if applied >= r.commit_target {
                let resp = (r.f)(&self.sm);
                let _ = r.tx.send(resp);
            } else {
                still.push(r);
            }
        }
        *reads = still;
    }

    async fn execute_actions(self: &Arc<Self>, actions: smallvec::SmallVec<[Action; 8]>) {
        for act in actions {
            match act {
                Action::PersistHardState(hs) => {
                    if let Err(e) = self.meta.save_hard_state(&hs) {
                        tracing::error!(error = %e, "failed to persist hard state");
                    }
                }
                Action::AppendEntries(entries) => {
                    if let Err(e) = self.log.append(&entries) {
                        tracing::error!(error = %e, "log append failed");
                    }
                }
                Action::TruncateLog { from } => {
                    if let Err(e) = self.log.truncate_from(from) {
                        tracing::error!(error = %e, "log truncate failed");
                    }
                }
                Action::SendMessage { to, msg } => {
                    if let Some(client) = self.peers.get(&to).cloned() {
                        tokio::spawn(async move { client.send(msg).await });
                    }
                }
                Action::ApplyCommitted { entries } => {
                    self.apply_entries(entries).await;
                }
                Action::TakeSnapshot { last_included_index, last_included_term } => {
                    if let Err(e) = self.take_snapshot(last_included_index, last_included_term) {
                        tracing::error!(error = %e, "snapshot failed");
                    }
                }
                Action::InstallSnapshot { last_included_index, last_included_term, .. } => {
                    let _ = self.meta.save_snapshot_pointer(last_included_index, last_included_term);
                    self.log.set_snapshot_tip(last_included_index);
                    self.node.lock().on_snapshot_taken(last_included_index, last_included_term);
                }
                Action::ResetElectionTimer | Action::ResetHeartbeatTimer => {}
                Action::NotifyReadIndex { .. } => {
                    self.drain_ready_reads();
                }
                Action::BecameLeader { term } => {
                    tracing::info!(term, node = self.cfg.id, "became leader");
                    self.refresh_soft();
                }
                Action::BecameFollower { term, leader } => {
                    tracing::info!(term, leader = ?leader, node = self.cfg.id, "became follower");
                    self.refresh_soft();
                }
                Action::Metric(ev) => self.observe(ev),
            }
        }
        self.drain_ready_reads();
        self.refresh_soft();
        self.notify.notify_waiters();
    }

    async fn apply_entries(self: &Arc<Self>, entries: Vec<Entry>) {
        for e in entries {
            let idx = e.index;
            let resp = match e.kind {
                EntryKind::Noop => Response::Ok,
                EntryKind::Normal => match bincode::deserialize::<Command>(&e.data) {
                    Ok(cmd) => match self.sm.apply(idx, &cmd) {
                        Ok(r) => r,
                        Err(err) => Response::Error(err.to_string()),
                    },
                    Err(_) => Response::Error("bad command encoding".into()),
                },
                EntryKind::ConfigJoint | EntryKind::ConfigNew => Response::Ok,
            };
            if let Some(p) = self.pending.lock().remove(&idx) {
                let _ = p.tx.send(Ok(resp));
            }
        }
    }

    fn take_snapshot(&self, last_included_index: u64, last_included_term: u64) -> anyhow::Result<()> {
        let meta = SnapshotMeta {
            last_included_index, last_included_term, config: vec![],
        };
        let writer = self.snaps.begin_write(&meta)?;
        writer.finish()?;
        self.meta.save_snapshot_pointer(last_included_index, last_included_term)?;
        self.log.compact(last_included_index)?;
        self.snaps.keep_last(3).ok();
        Ok(())
    }

    fn refresh_soft(&self) {
        let n = self.node.lock();
        let mut s = self.soft.lock();
        s.role = n.role().as_str().into();
        s.term = n.hard_state().current_term;
        s.leader = n.soft_state().leader_id;
        s.commit_index = n.commit_index();
        s.applied_index = self.sm.applied_index();
        s.last_log_index = n.last_log_index();
    }

    fn observe(&self, ev: MetricEvent) {
        match ev {
            MetricEvent::TermBumped { new_term } => tracing::info!(term = new_term, "term bumped"),
            MetricEvent::ElectionStarted { term, pre_vote } => tracing::info!(term, pre_vote, "election started"),
            MetricEvent::ElectionWon { term } => tracing::info!(term, "election won"),
            MetricEvent::CommitAdvanced { index } => tracing::debug!(index, "commit advanced"),
            MetricEvent::AppendRejected { from } => tracing::debug!(from, "append rejected"),
            MetricEvent::LeadershipTransferStarted { target } => tracing::info!(target, "leadership transfer"),
        }
    }

    /// Local node id.
    pub fn local_id(&self) -> NodeId { self.cfg.id }
}

/// gRPC handlers route inbound Raft RPCs through here. Each Raft RPC has
/// exactly one synchronous reply, so we step the core, harvest the response
/// `Action::SendMessage { to: <caller> }`, and return the body to the gRPC
/// handler — which sends it back over the same gRPC call.
#[async_trait]
impl MessageProcessor for Runtime {
    async fn process(&self, from: u64, msg: Message) -> Option<Message> {
        let actions = self.node.lock().step(msg);

        let mut response: Option<Message> = None;
        let mut other_actions: smallvec::SmallVec<[Action; 8]> = smallvec::SmallVec::new();
        for act in actions {
            match act {
                Action::SendMessage { to, msg } if to == from && response.is_none() && is_response_kind(&msg) => {
                    response = Some(msg);
                }
                a => other_actions.push(a),
            }
        }

        // Build an Arc-self for execute_actions reuse via a ref-counted shim.
        // We can't trivially get an `Arc<Self>` from `&self`, so we inline the
        // effects we care about here.
        for act in other_actions {
            match act {
                Action::PersistHardState(hs) => {
                    if let Err(e) = self.meta.save_hard_state(&hs) {
                        tracing::error!(error = %e, "failed to persist hard state");
                    }
                }
                Action::AppendEntries(entries) => {
                    if let Err(e) = self.log.append(&entries) {
                        tracing::error!(error = %e, "log append failed");
                    }
                }
                Action::TruncateLog { from } => {
                    if let Err(e) = self.log.truncate_from(from) {
                        tracing::error!(error = %e, "log truncate failed");
                    }
                }
                Action::SendMessage { to, msg } => {
                    if let Some(client) = self.peers.get(&to).cloned() {
                        tokio::spawn(async move { client.send(msg).await });
                    }
                }
                Action::ApplyCommitted { entries } => {
                    for e in entries {
                        let idx = e.index;
                        let resp = match e.kind {
                            EntryKind::Noop => Response::Ok,
                            EntryKind::Normal => match bincode::deserialize::<Command>(&e.data) {
                                Ok(cmd) => match self.sm.apply(idx, &cmd) {
                                    Ok(r) => r,
                                    Err(err) => Response::Error(err.to_string()),
                                },
                                Err(_) => Response::Error("bad command encoding".into()),
                            },
                            EntryKind::ConfigJoint | EntryKind::ConfigNew => Response::Ok,
                        };
                        if let Some(p) = self.pending.lock().remove(&idx) {
                            let _ = p.tx.send(Ok(resp));
                        }
                    }
                }
                Action::TakeSnapshot { last_included_index, last_included_term } => {
                    if let Err(e) = self.take_snapshot(last_included_index, last_included_term) {
                        tracing::error!(error = %e, "snapshot failed");
                    }
                }
                Action::InstallSnapshot { last_included_index, last_included_term, .. } => {
                    let _ = self.meta.save_snapshot_pointer(last_included_index, last_included_term);
                    self.log.set_snapshot_tip(last_included_index);
                    self.node.lock().on_snapshot_taken(last_included_index, last_included_term);
                }
                Action::ResetElectionTimer | Action::ResetHeartbeatTimer => {}
                Action::NotifyReadIndex { .. } => {
                    self.drain_ready_reads();
                }
                Action::BecameLeader { term } => {
                    tracing::info!(term, node = self.cfg.id, "became leader");
                    self.refresh_soft();
                }
                Action::BecameFollower { term, leader } => {
                    tracing::info!(term, leader = ?leader, node = self.cfg.id, "became follower");
                    self.refresh_soft();
                }
                Action::Metric(ev) => self.observe(ev),
            }
        }

        self.drain_ready_reads();
        self.refresh_soft();
        self.notify.notify_waiters();

        response
    }
}

fn is_response_kind(m: &Message) -> bool {
    matches!(
        m,
        Message::AppendEntriesResponse { .. }
            | Message::RequestVoteResponse { .. }
            | Message::InstallSnapshotResponse { .. }
    )
}

/// Channels owned by the run loop.
pub struct RuntimeHandles {
    /// Network -> runtime inbox sender.
    pub inbox_tx: InboxTx,
    /// Network -> runtime inbox receiver.
    pub inbox_rx: Inbox,
    /// Proposals receiver.
    pub(crate) proposal_rx: mpsc::UnboundedReceiver<ProposalRequest>,
    /// Reads receiver.
    pub(crate) read_rx: mpsc::UnboundedReceiver<ReadRequest>,
}

/// Implements the resp-server CommandHandler trait via the runtime channels.
pub struct ClientHandler {
    rt: Arc<Runtime>,
}

impl ClientHandler {
    /// Construct.
    pub fn new(rt: Arc<Runtime>) -> Self { Self { rt } }
}

#[async_trait]
impl CommandHandler for ClientHandler {
    async fn propose(&self, cmd: Command) -> Result<Response, ProposeFailure> {
        let (tx, rx) = oneshot::channel();
        let req = ProposalRequest { kind: ProposalKind::Command(cmd), tx };
        if self.rt.proposal_tx.send(req).is_err() {
            return Err(ProposeFailure::Other("runtime gone".into()));
        }
        match tokio::time::timeout(Duration::from_secs(5), rx).await {
            Ok(Ok(res)) => res,
            Ok(Err(_)) => Err(ProposeFailure::Other("dropped".into())),
            Err(_) => Err(ProposeFailure::Timeout),
        }
    }

    async fn linearizable_read<F>(&self, f: F) -> Result<RespFrame, ProposeFailure>
    where F: FnOnce(&KvStateMachine) -> RespFrame + Send + 'static {
        let (tx, rx) = oneshot::channel();
        let req = ReadRequest { f: Box::new(f), tx };
        if self.rt.read_tx.send(req).is_err() {
            return Err(ProposeFailure::Other("runtime gone".into()));
        }
        match tokio::time::timeout(Duration::from_secs(5), rx).await {
            Ok(Ok(frame)) => Ok(frame),
            Ok(Err(_)) => Err(ProposeFailure::Other("dropped".into())),
            Err(_) => Err(ProposeFailure::Timeout),
        }
    }

    fn is_leader(&self) -> bool {
        matches!(self.rt.node.lock().role(), Role::Leader)
    }

    fn info(&self, _section: Option<&str>) -> String {
        let s = self.rt.soft.lock();
        format!(
            "# Server\nraftkv_version:0.1.0\n# Raft\nrole:{}\nterm:{}\nleader_id:{}\ncommit_index:{}\napplied_index:{}\nlast_log_index:{}\n",
            s.role, s.term, s.leader.unwrap_or(0), s.commit_index, s.applied_index, s.last_log_index,
        )
    }

    fn cluster_nodes(&self) -> String {
        self.rt
            .addr_book
            .iter()
            .map(|(id, addr)| format!("{} {}", id, addr))
            .collect::<Vec<_>>()
            .join("\n")
    }
}
\end{lstlisting}

\section{crates/resp-server/src/codec.rs}
\begin{lstlisting}[language=Rust]
//! RESP2 codec. Decodes inline + multibulk; encodes simple/error/integer/bulk/array.

use bytes::{Buf, BufMut, BytesMut};
use std::io;
use tokio_util::codec::{Decoder, Encoder};

/// A single RESP frame.
#[derive(Debug, Clone, PartialEq, Eq)]
pub enum RespFrame {
    /// `+...\r\n`
    Simple(String),
    /// `-...\r\n`
    Error(String),
    /// `:N\r\n`
    Integer(i64),
    /// `$N\r\n<N bytes>\r\n` or nil.
    Bulk(Option<Vec<u8>>),
    /// `*N\r\n...`
    Array(Vec<RespFrame>),
}

impl RespFrame {
    /// Convenience: build a simple OK reply.
    pub fn ok() -> Self { Self::Simple("OK".into()) }
    /// Convenience: build a nil bulk reply.
    pub fn nil() -> Self { Self::Bulk(None) }
    /// Convenience: build an error.
    pub fn err(s: impl Into<String>) -> Self { Self::Error(s.into()) }
    /// Convenience: build an integer.
    pub fn int(n: i64) -> Self { Self::Integer(n) }
    /// Convenience: build a bulk string.
    pub fn bulk(b: impl Into<Vec<u8>>) -> Self { Self::Bulk(Some(b.into())) }
}

/// Tokio codec implementing RESP2.
#[derive(Default, Clone)]
pub struct RespCodec;

fn find_crlf(buf: &[u8]) -> Option<usize> {
    buf.windows(2).position(|w| w == b"\r\n")
}

fn parse_int(buf: &[u8]) -> io::Result<i64> {
    std::str::from_utf8(buf)
        .map_err(|_| io::Error::new(io::ErrorKind::InvalidData, "non-utf8 integer"))?
        .parse::<i64>()
        .map_err(|_| io::Error::new(io::ErrorKind::InvalidData, "bad integer"))
}

fn try_decode(buf: &[u8]) -> io::Result<Option<(RespFrame, usize)>> {
    if buf.is_empty() { return Ok(None); }
    match buf[0] {
        b'+' | b'-' | b':' | b'$' | b'*' => decode_typed(buf),
        _ => decode_inline(buf),
    }
}

fn decode_typed(buf: &[u8]) -> io::Result<Option<(RespFrame, usize)>> {
    let kind = buf[0];
    let line_end = match find_crlf(&buf[1..]) {
        Some(p) => p + 1,
        None => return Ok(None),
    };
    let body = &buf[1..line_end];
    match kind {
        b'+' => {
            let s = std::str::from_utf8(body)
                .map_err(|_| io::Error::new(io::ErrorKind::InvalidData, "non-utf8"))?
                .to_string();
            Ok(Some((RespFrame::Simple(s), line_end + 2)))
        }
        b'-' => {
            let s = std::str::from_utf8(body)
                .map_err(|_| io::Error::new(io::ErrorKind::InvalidData, "non-utf8"))?
                .to_string();
            Ok(Some((RespFrame::Error(s), line_end + 2)))
        }
        b':' => {
            let n = parse_int(body)?;
            Ok(Some((RespFrame::Integer(n), line_end + 2)))
        }
        b'$' => {
            let n = parse_int(body)?;
            let after_hdr = line_end + 2;
            if n < 0 { return Ok(Some((RespFrame::Bulk(None), after_hdr))); }
            let n_us = n as usize;
            if buf.len() < after_hdr + n_us + 2 { return Ok(None); }
            let payload = buf[after_hdr..after_hdr + n_us].to_vec();
            Ok(Some((RespFrame::Bulk(Some(payload)), after_hdr + n_us + 2)))
        }
        b'*' => {
            let n = parse_int(body)?;
            let mut consumed = line_end + 2;
            if n < 0 { return Ok(Some((RespFrame::Array(vec![]), consumed))); }
            let mut out = Vec::with_capacity(n as usize);
            for _ in 0..n {
                match try_decode(&buf[consumed..])? {
                    Some((f, c)) => { out.push(f); consumed += c; }
                    None => return Ok(None),
                }
            }
            Ok(Some((RespFrame::Array(out), consumed)))
        }
        _ => unreachable!(),
    }
}

fn decode_inline(buf: &[u8]) -> io::Result<Option<(RespFrame, usize)>> {
    let line_end = match find_crlf(buf) {
        Some(p) => p,
        None => return Ok(None),
    };
    let line = &buf[..line_end];
    let mut parts: Vec<RespFrame> = Vec::new();
    let s = std::str::from_utf8(line)
        .map_err(|_| io::Error::new(io::ErrorKind::InvalidData, "non-utf8 inline"))?;
    for tok in s.split_whitespace() {
        parts.push(RespFrame::Bulk(Some(tok.as_bytes().to_vec())));
    }
    Ok(Some((RespFrame::Array(parts), line_end + 2)))
}

impl Decoder for RespCodec {
    type Item = RespFrame;
    type Error = io::Error;

    fn decode(&mut self, src: &mut BytesMut) -> Result<Option<Self::Item>, Self::Error> {
        match try_decode(src.as_ref())? {
            Some((frame, consumed)) => {
                src.advance(consumed);
                Ok(Some(frame))
            }
            None => Ok(None),
        }
    }
}

fn encode_into(frame: &RespFrame, out: &mut BytesMut) {
    match frame {
        RespFrame::Simple(s) => {
            out.put_u8(b'+');
            out.extend_from_slice(s.as_bytes());
            out.put_slice(b"\r\n");
        }
        RespFrame::Error(s) => {
            out.put_u8(b'-');
            out.extend_from_slice(s.as_bytes());
            out.put_slice(b"\r\n");
        }
        RespFrame::Integer(n) => {
            out.put_u8(b':');
            out.extend_from_slice(n.to_string().as_bytes());
            out.put_slice(b"\r\n");
        }
        RespFrame::Bulk(None) => {
            out.put_slice(b"$-1\r\n");
        }
        RespFrame::Bulk(Some(b)) => {
            out.put_u8(b'$');
            out.extend_from_slice(b.len().to_string().as_bytes());
            out.put_slice(b"\r\n");
            out.extend_from_slice(b);
            out.put_slice(b"\r\n");
        }
        RespFrame::Array(arr) => {
            out.put_u8(b'*');
            out.extend_from_slice(arr.len().to_string().as_bytes());
            out.put_slice(b"\r\n");
            for f in arr { encode_into(f, out); }
        }
    }
}

impl Encoder<RespFrame> for RespCodec {
    type Error = io::Error;
    fn encode(&mut self, item: RespFrame, dst: &mut BytesMut) -> Result<(), Self::Error> {
        encode_into(&item, dst);
        Ok(())
    }
}

#[cfg(test)]
mod tests {
    use super::*;

    fn roundtrip(f: RespFrame) {
        let mut codec = RespCodec;
        let mut buf = BytesMut::new();
        codec.encode(f.clone(), &mut buf).unwrap();
        let decoded = codec.decode(&mut buf).unwrap().unwrap();
        assert_eq!(decoded, f);
    }

    #[test]
    fn roundtrip_all() {
        roundtrip(RespFrame::Simple("OK".into()));
        roundtrip(RespFrame::Error("ERR foo".into()));
        roundtrip(RespFrame::Integer(42));
        roundtrip(RespFrame::Bulk(None));
        roundtrip(RespFrame::Bulk(Some(b"hello".to_vec())));
        roundtrip(RespFrame::Array(vec![
            RespFrame::Bulk(Some(b"SET".to_vec())),
            RespFrame::Bulk(Some(b"k".to_vec())),
            RespFrame::Bulk(Some(b"v".to_vec())),
        ]));
    }

    #[test]
    fn inline_command_decodes() {
        let mut codec = RespCodec;
        let mut buf = BytesMut::from(&b"PING\r\n"[..]);
        let f = codec.decode(&mut buf).unwrap().unwrap();
        match f {
            RespFrame::Array(parts) => {
                assert_eq!(parts.len(), 1);
                if let RespFrame::Bulk(Some(b)) = &parts[0] { assert_eq!(b, b"PING"); }
                else { panic!("not bulk"); }
            }
            _ => panic!("not array"),
        }
    }
}
\end{lstlisting}

\section{crates/resp-server/src/commands.rs}
\begin{lstlisting}[language=Rust]
//! Parse RESP frames into commands.

use crate::codec::RespFrame;

/// Parsed RESP command.
#[derive(Debug, Clone)]
pub enum ClientCommand {
    /// PING
    Ping(Option<Vec<u8>>),
    /// ECHO msg
    Echo(Vec<u8>),
    /// GET key
    Get(Vec<u8>),
    /// SET key value [EX seconds]
    Set {
        /// Key.
        key: Vec<u8>,
        /// Value.
        value: Vec<u8>,
        /// TTL seconds, if any.
        ex_seconds: Option<u64>,
    },
    /// DEL key1 key2 ...
    Del(Vec<Vec<u8>>),
    /// EXISTS key1 ...
    Exists(Vec<Vec<u8>>),
    /// INCR key
    Incr(Vec<u8>),
    /// DECR key
    Decr(Vec<u8>),
    /// MSET k1 v1 ...
    MSet(Vec<(Vec<u8>, Vec<u8>)>),
    /// MGET k1 ...
    MGet(Vec<Vec<u8>>),
    /// EXPIRE key seconds
    Expire { key: Vec<u8>, seconds: u64 },
    /// TTL key
    Ttl(Vec<u8>),
    /// PERSIST key
    Persist(Vec<u8>),
    /// FLUSHDB
    FlushDb,
    /// DBSIZE
    DbSize,
    /// INFO [section]
    Info(Option<String>),
    /// CLUSTER NODES / CLUSTER INFO
    ClusterNodes,
    /// CLUSTER INFO
    ClusterInfo,
    /// COMMAND
    Command,
    /// CLIENT subcommands (best-effort)
    ClientNoop,
    /// SELECT n  (best-effort: only DB 0 supported)
    Select(u64),
    /// QUIT
    Quit,
    /// Unsupported / unknown
    Unknown(String),
}

fn extract_bulk(f: &RespFrame) -> Option<&[u8]> {
    if let RespFrame::Bulk(Some(b)) = f { Some(b.as_slice()) } else { None }
}

fn upper(b: &[u8]) -> String {
    String::from_utf8_lossy(b).to_ascii_uppercase()
}

/// Parse a top-level RESP frame as a client command.
pub fn parse(frame: &RespFrame) -> ClientCommand {
    let parts: Vec<&[u8]> = match frame {
        RespFrame::Array(items) => items.iter().filter_map(extract_bulk).collect(),
        _ => return ClientCommand::Unknown("expected array".into()),
    };
    if parts.is_empty() { return ClientCommand::Unknown("empty".into()); }
    let cmd = upper(parts[0]);
    let args = &parts[1..];
    match cmd.as_str() {
        "PING" => ClientCommand::Ping(args.first().map(|b| b.to_vec())),
        "ECHO" if args.len() == 1 => ClientCommand::Echo(args[0].to_vec()),
        "GET" if args.len() == 1 => ClientCommand::Get(args[0].to_vec()),
        "SET" if args.len() >= 2 => {
            let mut ex: Option<u64> = None;
            let mut i = 2;
            while i < args.len() {
                match upper(args[i]).as_str() {
                    "EX" if i + 1 < args.len() => {
                        if let Ok(n) = std::str::from_utf8(args[i+1]).unwrap_or("").parse::<u64>() {
                            ex = Some(n);
                        }
                        i += 2;
                    }
                    "PX" if i + 1 < args.len() => {
                        if let Ok(n) = std::str::from_utf8(args[i+1]).unwrap_or("").parse::<u64>() {
                            ex = Some(n / 1000);
                        }
                        i += 2;
                    }
                    _ => i += 1,
                }
            }
            ClientCommand::Set { key: args[0].to_vec(), value: args[1].to_vec(), ex_seconds: ex }
        }
        "DEL" if !args.is_empty() => ClientCommand::Del(args.iter().map(|b| b.to_vec()).collect()),
        "EXISTS" if !args.is_empty() => ClientCommand::Exists(args.iter().map(|b| b.to_vec()).collect()),
        "INCR" if args.len() == 1 => ClientCommand::Incr(args[0].to_vec()),
        "DECR" if args.len() == 1 => ClientCommand::Decr(args[0].to_vec()),
        "MSET" if !args.is_empty() && args.len() % 2 == 0 => {
            let mut pairs = Vec::with_capacity(args.len() / 2);
            for chunk in args.chunks(2) {
                pairs.push((chunk[0].to_vec(), chunk[1].to_vec()));
            }
            ClientCommand::MSet(pairs)
        }
        "MGET" if !args.is_empty() => ClientCommand::MGet(args.iter().map(|b| b.to_vec()).collect()),
        "EXPIRE" if args.len() == 2 => {
            let secs = std::str::from_utf8(args[1]).unwrap_or("").parse::<u64>().unwrap_or(0);
            ClientCommand::Expire { key: args[0].to_vec(), seconds: secs }
        }
        "TTL" if args.len() == 1 => ClientCommand::Ttl(args[0].to_vec()),
        "PERSIST" if args.len() == 1 => ClientCommand::Persist(args[0].to_vec()),
        "FLUSHDB" => ClientCommand::FlushDb,
        "DBSIZE" => ClientCommand::DbSize,
        "INFO" => ClientCommand::Info(args.first().map(|b| String::from_utf8_lossy(b).to_string())),
        "CLUSTER" if !args.is_empty() => match upper(args[0]).as_str() {
            "NODES" => ClientCommand::ClusterNodes,
            "INFO" => ClientCommand::ClusterInfo,
            _ => ClientCommand::Unknown(format!("CLUSTER {}", upper(args[0]))),
        },
        "COMMAND" => ClientCommand::Command,
        "CLIENT" => ClientCommand::ClientNoop,
        "SELECT" if args.len() == 1 => {
            let n = std::str::from_utf8(args[0]).unwrap_or("0").parse::<u64>().unwrap_or(0);
            ClientCommand::Select(n)
        }
        "QUIT" => ClientCommand::Quit,
        _ => ClientCommand::Unknown(cmd),
    }
}
\end{lstlisting}

\section{crates/resp-server/src/handler.rs}
\begin{lstlisting}[language=Rust]
//! Translate parsed RESP commands into either:
//!  * Local read against the state machine (gated by linearizable read-index
//!    in the runtime), or
//!  * A `Command` proposed to Raft for replication.

use std::sync::Arc;

use async_trait::async_trait;
use kv_state_machine::{Command, KvStateMachine, Response};

use crate::codec::RespFrame;
use crate::commands::ClientCommand;

/// Trait implemented by the runtime: handles proposing to Raft and waiting for
/// the apply notification.
#[async_trait]
pub trait CommandHandler: Send + Sync + 'static {
    /// Propose a command to Raft and await its application; return the
    /// state-machine response.
    async fn propose(&self, cmd: Command) -> Result<Response, ProposeFailure>;
    /// Read-index → wait until applied >= read_index → execute closure on the
    /// state machine.
    async fn linearizable_read<F>(&self, f: F) -> Result<RespFrame, ProposeFailure>
    where F: FnOnce(&KvStateMachine) -> RespFrame + Send + 'static;
    /// Whether this node is currently the leader.
    fn is_leader(&self) -> bool;
    /// Build the INFO string.
    fn info(&self, section: Option<&str>) -> String;
    /// Build CLUSTER NODES output.
    fn cluster_nodes(&self) -> String;
}

/// Failures the handler may return.
#[derive(Debug, Clone, thiserror::Error)]
pub enum ProposeFailure {
    /// Not leader; client is given a hint.
    #[error("not leader")]
    NotLeader {
        /// Optional leader hint string.
        hint: Option<String>,
    },
    /// Timed out waiting for commit/apply.
    #[error("timeout")]
    Timeout,
    /// Generic error.
    #[error("{0}")]
    Other(String),
}

/// Glue: convert a command into a frame.
pub fn response_to_frame(r: Response) -> RespFrame {
    match r {
        Response::Ok => RespFrame::ok(),
        Response::Int(n) => RespFrame::int(n),
        Response::Bulk(b) => RespFrame::Bulk(b),
        Response::Array(a) => RespFrame::Array(a.into_iter().map(RespFrame::Bulk).collect()),
        Response::Error(e) => RespFrame::err(e),
    }
}

/// Top-level dispatch.
pub async fn dispatch<H: CommandHandler + ?Sized>(
    handler: Arc<H>,
    cmd: ClientCommand,
    now_ms: u64,
) -> RespFrame {
    match cmd {
        ClientCommand::Ping(arg) => match arg {
            Some(b) => RespFrame::Bulk(Some(b)),
            None => RespFrame::Simple("PONG".into()),
        },
        ClientCommand::Echo(b) => RespFrame::Bulk(Some(b)),
        ClientCommand::Quit => RespFrame::ok(),
        ClientCommand::Select(_) => RespFrame::ok(),
        ClientCommand::Command => RespFrame::Array(vec![]),
        ClientCommand::ClientNoop => RespFrame::ok(),
        ClientCommand::Info(section) => RespFrame::Bulk(Some(handler.info(section.as_deref()).into_bytes())),
        ClientCommand::ClusterNodes => RespFrame::Bulk(Some(handler.cluster_nodes().into_bytes())),
        ClientCommand::ClusterInfo => RespFrame::Bulk(Some(b"cluster_enabled:1\r\ncluster_state:ok\r\n".to_vec())),
        ClientCommand::DbSize => match handler.linearizable_read(|sm| RespFrame::int(sm.len() as i64)).await {
            Ok(f) => f,
            Err(_) => RespFrame::err("ERR not leader or timeout"),
        },
        ClientCommand::Get(key) => {
            match handler.linearizable_read(move |sm| match sm.get(&key) {
                Ok(Some(v)) => RespFrame::Bulk(Some(v)),
                Ok(None) => RespFrame::nil(),
                Err(e) => RespFrame::err(format!("ERR {e}")),
            }).await {
                Ok(f) => f,
                Err(ProposeFailure::NotLeader { hint }) => RespFrame::err(format!("MOVED 0 {}", hint.unwrap_or_default())),
                Err(_) => RespFrame::err("ERR timeout"),
            }
        }
        ClientCommand::MGet(keys) => {
            match handler.linearizable_read(move |sm| match sm.mget(&keys) {
                Ok(vs) => RespFrame::Array(vs.into_iter().map(RespFrame::Bulk).collect()),
                Err(e) => RespFrame::err(format!("ERR {e}")),
            }).await {
                Ok(f) => f,
                Err(_) => RespFrame::err("ERR timeout"),
            }
        }
        ClientCommand::Exists(keys) => {
            match handler.linearizable_read(move |sm| {
                let mut n = 0i64;
                for k in &keys {
                    if sm.exists(k).unwrap_or(false) { n += 1; }
                }
                RespFrame::int(n)
            }).await {
                Ok(f) => f,
                Err(_) => RespFrame::err("ERR timeout"),
            }
        }
        ClientCommand::Set { key, value, ex_seconds } => {
            let expire_at_ms = ex_seconds.map(|s| now_ms + s * 1000);
            match handler.propose(Command::Set { key, value, expire_at_ms }).await {
                Ok(r) => response_to_frame(r),
                Err(ProposeFailure::NotLeader { hint }) => RespFrame::err(format!("MOVED 0 {}", hint.unwrap_or_default())),
                Err(e) => RespFrame::err(format!("ERR {e}")),
            }
        }
        ClientCommand::Del(keys) => match handler.propose(Command::Del { keys }).await {
            Ok(r) => response_to_frame(r),
            Err(e) => RespFrame::err(format!("ERR {e}")),
        },
        ClientCommand::Incr(key) => match handler.propose(Command::Incr { key, delta: 1 }).await {
            Ok(r) => response_to_frame(r),
            Err(e) => RespFrame::err(format!("ERR {e}")),
        },
        ClientCommand::Decr(key) => match handler.propose(Command::Incr { key, delta: -1 }).await {
            Ok(r) => response_to_frame(r),
            Err(e) => RespFrame::err(format!("ERR {e}")),
        },
        ClientCommand::MSet(pairs) => match handler.propose(Command::MSet { pairs }).await {
            Ok(r) => response_to_frame(r),
            Err(e) => RespFrame::err(format!("ERR {e}")),
        },
        ClientCommand::Expire { key, seconds } => {
            let expire_at_ms = now_ms + seconds * 1000;
            match handler.propose(Command::Expire { key, expire_at_ms }).await {
                Ok(r) => response_to_frame(r),
                Err(e) => RespFrame::err(format!("ERR {e}")),
            }
        }
        ClientCommand::Ttl(_key) => RespFrame::int(-1),
        ClientCommand::Persist(key) => match handler.propose(Command::Persist { key }).await {
            Ok(r) => response_to_frame(r),
            Err(e) => RespFrame::err(format!("ERR {e}")),
        },
        ClientCommand::FlushDb => match handler.propose(Command::FlushDb).await {
            Ok(r) => response_to_frame(r),
            Err(e) => RespFrame::err(format!("ERR {e}")),
        },
        ClientCommand::Unknown(c) => RespFrame::err(format!("ERR unknown command '{}'", c)),
    }
}
\end{lstlisting}

\section{crates/resp-server/src/lib.rs}
\begin{lstlisting}[language=Rust]
#![deny(unsafe_code)]
//! RESP2/RESP3 protocol server for a Raft-replicated KV.

pub mod codec;
pub mod commands;
pub mod handler;
pub mod server;

pub use codec::{RespCodec, RespFrame};
pub use handler::CommandHandler;
pub use server::serve;
\end{lstlisting}

\section{crates/resp-server/src/server.rs}
\begin{lstlisting}[language=Rust]
//! RESP TCP server.

use std::sync::Arc;
use std::time::{SystemTime, UNIX_EPOCH};

use futures::{SinkExt, StreamExt};
use tokio::net::TcpListener;
use tokio_util::codec::Framed;

use crate::codec::{RespCodec, RespFrame};
use crate::commands::parse;
use crate::handler::{dispatch, CommandHandler};

/// Run a RESP server on `addr`. Each connection is handled by a spawned task.
pub async fn serve<H: CommandHandler + 'static>(
    addr: &str,
    handler: Arc<H>,
) -> std::io::Result<()> {
    let listener = TcpListener::bind(addr).await?;
    tracing::info!(addr = addr, "RESP server listening");
    loop {
        let (sock, peer) = match listener.accept().await {
            Ok(x) => x,
            Err(e) => { tracing::warn!(error = %e, "accept failed"); continue; }
        };
        let handler = handler.clone();
        tokio::spawn(async move {
            sock.set_nodelay(true).ok();
            let mut framed = Framed::new(sock, RespCodec);
            tracing::debug!(client = %peer, "client connected");
            while let Some(frame_res) = framed.next().await {
                match frame_res {
                    Ok(frame) => {
                        let cmd = parse(&frame);
                        let now_ms = SystemTime::now()
                            .duration_since(UNIX_EPOCH)
                            .map(|d| d.as_millis() as u64)
                            .unwrap_or(0);
                        let resp: RespFrame = dispatch(handler.clone(), cmd, now_ms).await;
                        if let Err(e) = framed.send(resp).await {
                            tracing::debug!(error = %e, "send failed");
                            break;
                        }
                    }
                    Err(e) => {
                        tracing::debug!(error = %e, "client read error");
                        break;
                    }
                }
            }
        });
    }
}
\end{lstlisting}

\section{crates/sim-tests/src/lib.rs}
\begin{lstlisting}[language=Rust]
#![deny(unsafe_code)]
//! Deterministic simulation harness around `raft_core::RaftNode`.
//!
//! The harness is fully synchronous — it advances a logical clock by ticks and
//! delivers messages from a network that the test can perturb (drop, delay,
//! reorder, partition). Tests built on top of this are 100% reproducible from
//! a single PRNG seed.
//!
//! See `tests/` for concrete invariants (election liveness, log matching,
//! commit safety under partitions).

use std::collections::{BTreeMap, VecDeque};

use rand::Rng;
use rand::SeedableRng;
use rand_chacha::ChaCha20Rng;

use raft_core::{
    Action, HardState, Message, NodeId, RaftConfig, RaftLog, RaftNode, Role,
};

/// One simulated cluster node: a `RaftNode` plus an inbox.
pub struct SimNode {
    /// The node id.
    pub id: NodeId,
    /// The pure raft state machine.
    pub node: RaftNode,
    /// Pending in-bound messages (delivery queue).
    pub inbox: VecDeque<Message>,
    /// Whether this node is currently partitioned (drops all in/out).
    pub partitioned: bool,
}

impl SimNode {
    /// Construct a simulated node from a config.
    pub fn new(cfg: RaftConfig) -> Self {
        let id = cfg.id;
        Self {
            id,
            node: RaftNode::new(cfg, HardState::default(), RaftLog::new()),
            inbox: VecDeque::new(),
            partitioned: false,
        }
    }
}

/// In-flight network message (with a delivery time).
pub struct InFlight {
    /// Source node.
    pub from: NodeId,
    /// Destination node.
    pub to: NodeId,
    /// Delivery time in absolute ticks.
    pub deliver_at: u64,
    /// Message body.
    pub msg: Message,
}

/// The simulation cluster.
pub struct Sim {
    /// Per-node state.
    pub nodes: BTreeMap<NodeId, SimNode>,
    /// In-flight messages (will be delivered when `now >= deliver_at`).
    pub network: Vec<InFlight>,
    /// Current tick.
    pub now: u64,
    /// Network delay range (in ticks).
    pub delay_min: u64,
    /// Upper bound (exclusive).
    pub delay_max: u64,
    /// Per-step drop probability (0..=1).
    pub drop_p: f64,
    /// Deterministic PRNG.
    pub rng: ChaCha20Rng,
}

impl Sim {
    /// Build a simulation cluster.
    pub fn new(ids: &[NodeId], seed: u64) -> Self {
        let mut nodes = BTreeMap::new();
        for &id in ids {
            let cfg = RaftConfig {
                id,
                initial_voters: ids.to_vec(),
                election_timeout_ticks: 10,
                heartbeat_ticks: 1,
                max_append_entries: 16,
                max_inflight_msgs: 64,
                rng_seed: seed.wrapping_add(id),
                pre_vote: true,
            };
            nodes.insert(id, SimNode::new(cfg));
        }
        Self {
            nodes,
            network: Vec::new(),
            now: 0,
            delay_min: 1,
            delay_max: 3,
            drop_p: 0.0,
            rng: ChaCha20Rng::seed_from_u64(seed),
        }
    }

    /// Advance time by one tick on every node and deliver due messages.
    pub fn step(&mut self) {
        self.now += 1;
        // Tick every node.
        let ids: Vec<NodeId> = self.nodes.keys().copied().collect();
        for id in &ids {
            let acts = self.nodes.get_mut(id).unwrap().node.tick();
            self.dispatch(*id, acts);
        }
        // Deliver due messages.
        let mut keep = Vec::with_capacity(self.network.len());
        let due: Vec<InFlight> = std::mem::take(&mut self.network)
            .into_iter()
            .filter_map(|m| if m.deliver_at <= self.now { Some(m) } else { keep.push(m); None })
            .collect();
        self.network = keep;
        for m in due {
            if let Some(node) = self.nodes.get_mut(&m.to) {
                if node.partitioned { continue; }
                let acts = node.node.step(m.msg);
                self.dispatch(m.to, acts);
            }
        }
    }

    /// Run until predicate is satisfied or `max_ticks` elapses.
    pub fn run_until<F: Fn(&Self) -> bool>(&mut self, pred: F, max_ticks: u64) -> bool {
        for _ in 0..max_ticks {
            if pred(self) { return true; }
            self.step();
        }
        pred(self)
    }

    /// Find the current leader, if any.
    pub fn leader(&self) -> Option<NodeId> {
        self.nodes.values().find(|n| matches!(n.node.role(), Role::Leader)).map(|n| n.id)
    }

    fn dispatch<I: IntoIterator<Item = Action>>(&mut self, from: NodeId, acts: I) {
        let from_partitioned = self.nodes.get(&from).map(|n| n.partitioned).unwrap_or(false);
        for a in acts {
            if let Action::SendMessage { to, msg } = a {
                if from_partitioned { continue; }
                if self.rng.gen_bool(self.drop_p) { continue; }
                let lo = self.delay_min;
                let hi = self.delay_max.max(lo + 1);
                let d = self.rng.gen_range(lo..hi);
                self.network.push(InFlight { from, to, deliver_at: self.now + d, msg });
            }
        }
    }
}

#[cfg(test)]
mod tests {
    use super::*;

    #[test]
    fn three_node_cluster_elects_a_leader() {
        let mut sim = Sim::new(&[1, 2, 3], 42);
        let elected = sim.run_until(|s| s.leader().is_some(), 500);
        assert!(elected, "expected a leader within 500 ticks");
    }

    #[test]
    fn deterministic_replay() {
        let mut a = Sim::new(&[1, 2, 3, 4, 5], 0xCAFE);
        let mut b = Sim::new(&[1, 2, 3, 4, 5], 0xCAFE);
        for _ in 0..1000 { a.step(); b.step(); }
        assert_eq!(a.leader(), b.leader());
    }
}
\end{lstlisting}

\section{proto/admin.proto}
\begin{lstlisting}[language=protobuf]
syntax = "proto3";

package raftkv.admin.v1;

service Admin {
  rpc Status (StatusRequest) returns (StatusResponse);
  rpc TriggerSnapshot (TriggerSnapshotRequest) returns (TriggerSnapshotResponse);
}

message StatusRequest {}

message StatusResponse {
  uint64 node_id = 1;
  string role = 2;
  uint64 term = 3;
  uint64 commit_index = 4;
  uint64 applied_index = 5;
  uint64 last_log_index = 6;
  uint64 leader_id = 7;
  repeated NodeInfo members = 8;
}

message NodeInfo {
  uint64 id = 1;
  string addr = 2;
  bool   voter = 3;
}

message TriggerSnapshotRequest {}
message TriggerSnapshotResponse { bool accepted = 1; }
\end{lstlisting}

\section{proto/membership.proto}
\begin{lstlisting}[language=protobuf]
syntax = "proto3";

package raftkv.membership.v1;

service Membership {
  rpc AddServer (AddServerRequest) returns (AddServerResponse);
  rpc RemoveServer (RemoveServerRequest) returns (RemoveServerResponse);
  rpc TransferLeadership (TransferLeadershipRequest) returns (TransferLeadershipResponse);
}

message Server {
  uint64 id = 1;
  string addr = 2;
  bool   voter = 3;
}

message AddServerRequest {
  Server server = 1;
}

message AddServerResponse {
  Status status = 1;
  string leader_hint = 2;
}

message RemoveServerRequest {
  uint64 id = 1;
}

message RemoveServerResponse {
  Status status = 1;
  string leader_hint = 2;
}

message TransferLeadershipRequest {
  uint64 target_id = 1;
}

message TransferLeadershipResponse {
  Status status = 1;
}

enum Status {
  STATUS_UNSPECIFIED = 0;
  STATUS_OK = 1;
  STATUS_NOT_LEADER = 2;
  STATUS_TIMEOUT = 3;
  STATUS_REJECTED = 4;
}
\end{lstlisting}

\section{proto/raft.proto}
\begin{lstlisting}[language=protobuf]
syntax = "proto3";

package raftkv.raft.v1;

// Wire-level Raft RPCs. Mirrors the message types in raft-core::message but
// preserves stable field numbering for cross-version compatibility.

service Raft {
  rpc AppendEntries (AppendEntriesRequest) returns (AppendEntriesResponse);
  rpc RequestVote   (RequestVoteRequest)   returns (RequestVoteResponse);
  rpc PreVote       (RequestVoteRequest)   returns (RequestVoteResponse);
  rpc InstallSnapshot (stream InstallSnapshotChunk) returns (InstallSnapshotResponse);
  rpc TimeoutNow    (TimeoutNowRequest)    returns (TimeoutNowResponse);
}

message LogEntry {
  uint64 term  = 1;
  uint64 index = 2;
  EntryType entry_type = 3;
  bytes data = 4;
}

enum EntryType {
  ENTRY_TYPE_NORMAL        = 0;
  ENTRY_TYPE_NOOP          = 1;
  ENTRY_TYPE_CONFIG_JOINT  = 2;
  ENTRY_TYPE_CONFIG_NEW    = 3;
}

message AppendEntriesRequest {
  uint64 term = 1;
  uint64 leader_id = 2;
  uint64 prev_log_index = 3;
  uint64 prev_log_term = 4;
  repeated LogEntry entries = 5;
  uint64 leader_commit = 6;
  bytes  request_id = 7;
}

message AppendEntriesResponse {
  uint64 term = 1;
  bool   success = 2;
  // Conflict information for fast back-off (Ongaro thesis §3.5)
  uint64 conflict_index = 3;
  uint64 conflict_term  = 4;
  uint64 last_log_index = 5;
  bytes  request_id = 6;
}

message RequestVoteRequest {
  uint64 term = 1;
  uint64 candidate_id = 2;
  uint64 last_log_index = 3;
  uint64 last_log_term = 4;
  bool   pre_vote = 5;
  bytes  request_id = 6;
}

message RequestVoteResponse {
  uint64 term = 1;
  bool   vote_granted = 2;
  bytes  request_id = 3;
}

message InstallSnapshotChunk {
  uint64 term = 1;
  uint64 leader_id = 2;
  uint64 last_included_index = 3;
  uint64 last_included_term = 4;
  uint64 offset = 5;
  bytes  data = 6;
  bool   done = 7;
  bytes  config = 8;
  bytes  request_id = 9;
}

message InstallSnapshotResponse {
  uint64 term = 1;
  bytes  request_id = 2;
}

message TimeoutNowRequest {
  uint64 term = 1;
  uint64 leader_id = 2;
}

message TimeoutNowResponse {
  bool accepted = 1;
}
\end{lstlisting}

\end{document}